% Chapter 1, Section 1 _Linear Algebra_ Jim Hefferon
%  http://joshua.smcvt.edu/linearalgebra
%  2001-Jun-09
\chapter{线性方程组}\leavevmode  % because no preamble?
\section{求解线性方程组}
线性方程组在科学和数学中很常见。
来自高中科学的这两个例子\cite{Onan}
可以让我们体会它们是如何产生的。

第一个例子来自
\hypertarget{ex:Statics}{静力学}。\index{Statics problem}
假设我们有三个物体，
已知其中一个的质量为$2$~千克，
而想求出两个未知的质量。
用一根米尺做实验，得到下面两个平衡状态。
\begin{center}
  \includegraphics{gr/mp/ch1.1}
  \qquad
  \includegraphics{gr/mp/ch1.2}
\end{center}
要使两边平衡，必须有左边的力矩之和等于右边的力矩之和，其中
一个物体的力矩是它的质量乘以它到
支点的距离。
这就给出了一个含两个线性方程的方程组。
\begin{align*}  % odd formatting to match chem problem below
      40h + 15c &= 100  \\
            25c &= 50+50h
\end{align*}

第二个例子
来自化学。\index{Chemistry problem}
在可控条件下，我们可以将甲苯$\hbox{C}_7\hbox{H}_8$与
硝酸$\hbox{H}\hbox{N}\hbox{O}_3$混合，生成
三硝基甲苯$\hbox{C}_7\hbox{H}_5\hbox{O}_6\hbox{N}_3$，
同时得到副产物水
（条件必须受到极好的控制\Dash 三硝基甲苯
更为人熟知的名称是TNT）。
我们应该按怎样的比例混合它们？
反应前出现的每种元素的原子个数
\begin{equation*}
    x\, {\rm C}_7{\rm H}_8\ +\ y\,{\rm H}{\rm N}{\rm O}_3
    \quad\longrightarrow\quad
    z\,{\rm C}_7{\rm H}_5{\rm O}_6{\rm N}_3\ +\ w\,{\rm H}_2{\rm O}
\tag*{}\end{equation*}
必须等于反应后出现的原子个数。
把它依次应用于元素 C、H、N 和 O，就得到
下面的方程组。
\begin{align*}  % odd formatting to state more naturally
      7x      &= 7z  \\
      8x +1y  &= 5z+2w  \\
      1y      &= 3z  \\
      3y      &= 6z+1w
\end{align*}

这两个例子都归结为求解一个方程组。
在每个方程组中，方程只含有每个变量的一次幂。
本章将说明如何求解任何这样的方程组。















\subsection{高斯消元法}
\begin{definition}\label{def:linearcombination}
%<*df:linearcombination>
\( x_1 \), \ldots, \( x_n \) 的一个\definend{线性组合}\index{linear combination}形如
\begin{equation*}
   a_1x_1+a_2x_2+a_3x_3+\cdots+a_nx_n
\end{equation*}
其中数\( a_1, \ldots ,a_n\in\Re \)称为该组合的
\definend{系数}\index{linear equation!coefficients}。
%</df:linearcombination>
%<*df:linearequations>
变量 $x_1$, \ldots, $x_n$ 的一个\definend{线性方程}\index{linear equation}
形如
$a_1x_1+a_2x_2+a_3x_3+\cdots+a_nx_n=d$
其中
\( d\in\Re \) 称为\definend{常数}\index{linear equation!constant}。

\( n \) 元组 \( (s_1,s_2,\ldots ,s_n)\in\Re^n \) 是该方程的一个
\definend{解}\index{linear equation!solution of} %
（或者说它\definend{满足}该方程），如果用数
$s_1$, \ldots, $s_n$ 代替各变量
得到一个真的陈述：
$a_1s_1+a_2s_2+\cdots+a_ns_n=d$。
一个\definend{线性方程组}\index{linear equation!system of}%
\index{system of linear equations}
\begin{equation*}
  \begin{linsys}{4}
    a_{1,1}x_1 &+ &a_{1,2}x_2  &+  &\cdots &+ &a_{1,n}x_n &=  &d_1  \\
    a_{2,1}x_1 &+ &a_{2,2}x_2  &+  &\cdots &+ &a_{2,n}x_n &=  &d_2  \\
             &  &           &   &       &  &          &\vdotswithin{=}  \\
    a_{m,1}x_1 &+ &a_{m,2}x_2  &+  &\cdots &+ &a_{m,n}x_n &=  &d_m
  \end{linsys}
\end{equation*}
有解
\( (s_1,s_2,\ldots ,s_n) \)，如果该 $n$ 元组是所有
方程的解。
%</df:linearequations>
\end{definition}

\begin{example}
组合 \( 3x_1 + 2x_2 \) 是 $x_1$ 与 $x_2$ 的线性组合。
组合 \( 3x_1^2 + 2x_2 \) 不是 $x_1$ 与 $x_2$ 的线性函数，
\( 3x_1 + 2\sin(x_2) \) 也不是。
\end{example}

我们通常取 $x_1$, \ldots, $x_n$ 彼此不同，因为在含有重复项的和式中，
我们可以重新整理使各项唯一，如 $2x+3y+4x=6x+3y$。
我们有时写出系数为零的项，如
$x-2y+0z$，有时则将其省略，视哪种方便而定。

\begin{example}
有序对 \( (-1,5) \) 是下面这个方程组的解。
\begin{equation*}
  \begin{linsys}{2}
    3x_1 &+ &2x_2 &= &7  \\
    -x_1 &+ &x_2  &= &6
  \end{linsys}
\end{equation*}
相比之下，\( (5,-1) \) 不是解。
\end{example}

求出由全部解组成的集合称为
\definend{求解}\index{system of linear equations!solving}
方程组。
我们不需要猜测或碰运气；
有一个总是有效的算法。
这个算法就是
\definend{高斯消元法}\index{Gauss's Method}%
\index{system of linear equations!Gauss's Method}
（或称\definend{高斯消去法}\index{Gaussian elimination}%
\index{system of linear equations!Gaussian elimination}
或\definend{线性消去法}\index{linear elimination}%
\index{system of linear equations!linear elimination}%
\index{system of linear equations!elimination}%
\index{elimination, Gaussian}）。
% It transforms the system, step by step, into one
% with a form that we can easily solve.
% We will first illustrate how it goes and then we will see the
% formal statement.

\begin{example}
为求解方程组
\begin{equation*}
  \begin{linsys}{3}
                     &   &      &   &3x_3  &=  &9  \\
                 x_1 &+  &5x_2  &-  &2x_3  &=  &2  \\
      \frac{1}{3}x_1 &+  &2x_2  &   &      &=  &3
  \end{linsys}
\end{equation*}
我们一步一步地变换它，直到它化成
我们容易求解的形式。

第一个变换
通过对换第一行与第三行来重写方程组。
\begin{align*}
  \quad
  &\grstep{ \text{对换第1行与第3行} }
  \begin{linsys}{3}
        \frac{1}{3}x_1 &+  &2x_2  &   &      &=  &3  \\
                   x_1 &+  &5x_2  &-  &2x_3  &=  &2  \\
                       &   &      &   &3x_3  &=  &9
    \end{linsys}
\end{align*}
第二个变换将第一行倍乘~$3$。
\begin{align*}
\quad
  &\grstep{ \text{将第1行倍乘3} }
  \begin{linsys}{3}
        x_1 &+  &6x_2  &   &      &=  &9  \\
        x_1 &+  &5x_2  &-  &2x_3  &=  &2  \\
            &   &      &   &3x_3  &=  &9
   \end{linsys}
\end{align*}
第三个变换是本例中唯一非平凡的变换。
我们在心中将第一行的两边乘以 \( -1 \)，
再在心中把它加到第二行上，
并把所得结果写成新的第二行。
\begin{align*}
  &\grstep{ \text{将第1行的\(-1\)倍加到第2行} }
  \begin{linsys}{3}
        x_1 &+  &6x_2  &   &      &=  &9  \\
            &   &-x_2  &-  &2x_3  &=  &-7 \\
            &   &      &   &3x_3  &=  &9
   \end{linsys}
\end{align*}
这些步骤已把方程组化成
我们容易求出每个变量之值的形式。
最下面的方程表明 \( x_3=3 \)。
用 $3$ 代替中间方程中的 \( x_3 \)，可知 \( x_2=1 \)。
再将这两个值代入最上面的方程，
得到 \( x_1=3 \)。
于是方程组有唯一解；
解集为\set{(3,1,3)}。
\end{example}

我们将在全书使用高斯消元法。
它快捷而容易。
我们现在来证明
它还是安全的：
高斯消元法从不丢失解，也从不
引入多余的解，因此
一个元组在我们应用该方法之前是方程组的解，当且仅当
在应用之后它仍是方程组的解。

\begin{theorem}[高斯消元法]
\index{linear equation!solution of!Gauss's Method} \label{th:GaussMethod}
%<*th:GaussMethod>
若一个线性方程组经由下述操作之一
变成另一个方程组
\begin{enumerate}
  \setlength{\itemsep}{0ex}
  \item
    一个方程与另一个方程对换
  \item
    一个方程的两边乘以一个非零常数
  \item
    一个方程被替换成它本身与另一个方程的倍数之和
\end{enumerate}
则这两个方程组有相同的解集。
%</th:GaussMethod>
\end{theorem}

这三种操作各有一个限制。
不允许把一行乘以~\( 0 \)，因为显然那样
会改变解集。
类似地，不允许把一行的倍数加到该行自身，因为
把一行加上它的~\( -1 \) 倍相当于把该行
乘以~\( 0 \)。
我们也不允许把一行与它自身对换，
这是为了让第四章的某些结果更容易。
再说，这种操作毫无意义。

\begin{proof}
我们在这里讨论方程对换这种操作。
另外两种
情形类似，见\nearbyexercise{ex:ProveGaussMethod}。

%<*pf:GaussMethod0>
考虑一个线性方程组。
\begin{equation*}
  \begin{linsys}{4}
    a_{1,1}x_1  &+  &a_{1,2}x_2 &+  &\cdots  &+ &a_{1,n}x_n  &=  &d_1\hfill\hbox{}  \\
                &   &           &   &        &   &     &\vdotswithin{=}  \\
    a_{i,1}x_1  &+  &a_{i,2}x_2 &+  &\cdots  &+ &a_{i,n}x_n  &=  &d_i\hfill\hbox{}  \\
                &   &           &   &        &   &     &\vdotswithin{=}  \\
    a_{j,1}x_1  &+  &a_{j,2}x_2 &+  &\cdots  &+ &a_{j,n}x_n  &=  &d_j\hfill\hbox{}  \\
                &   &           &   &        &   &     &\vdotswithin{=} \\
    a_{m,1}x_1  &+  &a_{m,2}x_2 &+  &\cdots  &+ &a_{m,n}x_n  &=  &d_m  \\
  \end{linsys}
\end{equation*}
元组 \( (s_1,\ldots\,,s_n) \)
满足这个方程组，
当且仅当将这些值代入各变量，即用 $s$ 代替 $x$，得到一个真陈述的合取：
$a_{1,1}s_1+a_{1,2}s_2+\cdots+a_{1,n}s_n=d_1$
与\ldots\
$a_{i,1}s_1+a_{i,2}s_2+\cdots+a_{i,n}s_n=d_i$
与\ldots\ $a_{j,1}s_1+a_{j,2}s_2+\cdots+a_{j,n}s_n=d_j$
与\ldots\ $a_{m,1}s_1+a_{m,2}s_2+\cdots+a_{m,n}s_n=d_m$。
%</pf:GaussMethod0>

%<*pf:GaussMethod1>
在用“与”连接的一列陈述中，我们可以
重新排列这些陈述的次序。
因此
该要求成立当且仅当
$a_{1,1}s_1+a_{1,2}s_2+\cdots+a_{1,n}s_n=d_1$
与\ldots\ $a_{j,1}s_1+a_{j,2}s_2+\cdots+a_{j,n}s_n=d_j$
与\ldots\ $a_{i,1}s_1+a_{i,2}s_2+\cdots+a_{i,n}s_n=d_i$
与\ldots\ $a_{m,1}s_1+a_{m,2}s_2+\cdots+a_{m,n}s_n=d_m$。
这正是 \( (s_1,\ldots\,,s_n) \)
在对换行之后满足方程组的要求。
%</pf:GaussMethod1>
\end{proof}

\begin{definition} \label{df:GaussMethod}
%<*df:GaussMethod>
来自
\nearbytheorem{th:GaussMethod}
的这三种操作是
\definend{初等化简操作},\index{Gauss's Method!elementary operations}%
\index{elementary reduction operations}
或称\definend{行变换}\index{elementary row operations}，
或称\definend{高斯操作}。
它们是
\definend{对换}\index{elementary reduction operations! swapping}%
\index{swapping rows}、
\definend{倍乘一个标量}（或称
\definend{倍乘}\index{elementary reduction operations! rescaling}%
\index{rescaling rows}），以及
\definend{倍加}\index{elementary reduction operations! row combination}%
\index{combining rows}\index{adding rows}。
%</df:GaussMethod>
\end{definition}

在写出计算过程时，我们将
把“第 \(i\) 行”缩写为“\( \rho_i \)”
（这是希腊字母 rho，读音如同英文的“row”）。
例如，我们将用
\( k\rho_i+\rho_j \)
表示一个倍加操作，
其中发生变化的行写在后面。
为了减少书写，当多个加法步骤使用同一个 $\rho_i$ 时，我们常常
将它们合并，如下面的
例子所示。

\begin{example}
高斯消元法系统地应用行变换来求解方程组。
下面是一个典型的例子。
\begin{equation*}
  \begin{linsys}{3}
    x  &+  &y  &   &   &=  &0  \\
   2x  &-  &y  &+  &3z &=  &3  \\
    x  &-  &2y &-  &z  &=  &3
  \end{linsys}
\end{equation*}
我们先用第一行
消去第二行中的 $2x$ 和第三行中的 $x$。
为消去 $2x$，我们在心中将整个第一行乘以 $-2$，
把它加到
第二行上，并把结果写成新的第二行。
为消去第三行中的~$x$，我们将第一行乘以
$-1$，加到第三行上，并把结果写成
新的第三行。
\begin{align*}
  &\grstep[-\rho_1 +\rho_3]{-2\rho_1 +\rho_2}
  \begin{linsys}{3}
     x  &+  &y  &   &   &=  &0  \\
        &   &-3y&+  &3z &=  &3  \\
        &   &-3y&-  &z  &=  &3
  \end{linsys}
\end{align*}
% In this version of the system, the last two equations involve only two unknowns.
最后，我们把第二个方程组变换成第三个方程组，使得
最下面的方程只含一个未知数。
为此我们用
第二行消去第三行中的 $y$~项。
\begin{equation*}
  \grstep{-\rho_2 +\rho_3}
  \begin{linsys}{3}
     x  &+  &y  &   &   &=  &0  \\
        &   &-3y&+  &3z &=  &3  \\
        &   &   &   &-4z&=  &0
   \end{linsys}
\end{equation*}
现在求方程组的解就容易了。
第三行给出 \( z=0 \)。
将它回代\index{Gauss's Method!back-substitution}%
\index{back-substitution}
到第二行，得到 \( y=-1 \)。
再回代到第一行，得到 \( x=1 \)。
\end{example}

\begin{example}
对于本章开头来自物理的问题\index{Statics problem}，高斯消元法给出
\begin{equation*}
   \begin{linsys}{2}
     40h  &+  &15c  &=  &100      \\
     -50h &+  &25c  &=  &50
   \end{linsys}
   \grstep{5/4\rho_1 +\rho_2}
   \begin{linsys}{2}
      40h  &+  &15c       &=  &100      \\
           &   &(175/4)c  &=  &175
    \end{linsys}
\end{equation*}
于是 \( c=4 \)，回代给出 \( h=1 \)。
（化学问题我们稍后求解。）
\end{example}

\begin{example}
行化简
\begin{align*}
   \begin{linsys}{3}
        x  &+  &y  &+  &z  &=  &9  \\
       2x  &+  &4y &-  &3z &=  &1  \\
       3x  &+  &6y &-  &5z &=  &0
   \end{linsys}
   &\grstep[-3\rho_1 +\rho_3]{-2\rho_1 +\rho_2}
   \begin{linsys}{3}
      x  &+  &y  &+  &z  &=  &9  \\
         &   &2y &-  &5z &=  &-17\\
         &   &3y &-  &8z&=  &-27
    \end{linsys}                                    \\
   &\grstep{-(3/2)\rho_2+\rho_3}
   \begin{linsys}{3}
      x  &+  &y  &+  &z            &=  &9  \\
         &   &2y &-  &5z           &=  &-17\\
         &   &   &   &-(1/2)z      &=  &-(3/2)
    \end{linsys}
\end{align*}
表明 \( z=3 \)、\( y=-1 \)、\( x=7 \)。
\end{example}

如上面的例子所示，高斯消元法
的要点是利用初等化简
操作来为回代做好准备。

\begin{definition} \label{df:EchelonForm}
%<*df:EchelonForm>
在方程组的每一行中，
系数非零的第一个变量是该行的
\definend{首变量}\index{echelon form!leading variable}%
\index{leading!variable}。%
若每个首变量
都位于其上一行的首变量的右侧，
第一行的首变量除外，
且系数全为零的行都在底部，
则称方程组处于\definend{阶梯形}\index{echelon form}。
%</df:EchelonForm>
\end{definition}

\begin{example}
前面的三个例题只用了倍加操作。
下面这个线性方程组需要对换操作
才能化成阶梯形，因为
在第一次倍加之后
\begin{align*}
   \begin{linsys}{4}
                x  &-  &y  &   &   &   &   &=  &0  \\
               2x  &-  &2y &+  &z  &+  &2w &=  &4  \\
                   &   &y  &   &   &+  &w  &=  &0  \\
                   &   &   &   &2z &+  &w  &=  &5
   \end{linsys}
   &\grstep{-2\rho_1 +\rho_2}
   \begin{linsys}{4}
      x  &-  &y  &\spaceforemptycolumn   &   &   &   &=  &0  \\
         &   &   &   &z  &+  &2w &=  &4  \\
         &   &y  &   &   &+  &w  &=  &0  \\
         &   &   &   &2z &+  &w  &=  &5
   \end{linsys}
\end{align*}
第二个方程没有首变量 $y$。
我们将它与下方有首变量 $y$ 的一行对换。
\begin{align*}
   &\grstep{\rho_2 \leftrightarrow\rho_3}
   \begin{linsys}{4}
      x  &-  &y  &\spaceforemptycolumn   &   &   &   &=  &0  \\
         &   &y  &   &   &+  &w  &=  &0  \\
         &   &   &   &z  &+  &2w &=  &4  \\
         &   &   &   &2z &+  &w  &=  &5
    \end{linsys}
\end{align*}
（若第二个方程下方有多于一个合适的行，
则我们可以任选其一。）
此后，高斯消元法像之前那样继续进行。
\begin{align*}
   &\grstep{-2\rho_3 +\rho_4}
   \begin{linsys}{4}
      x  &-  &y  &\spaceforemptycolumn   &   &   &   &=  &0  \\
         &   &y  &   &   &+  &w  &=  &0  \\
         &   &   &   &z  &+  &2w &=  &4  \\
         &   &   &   &   &   &-3w&=  &-3
    \end{linsys}
\end{align*}
回代给出 \( w=1 \)、\( z=2 \)、\( y=-1 \) 与 \( x=-1 \)。
\end{example}

严格说来，求解线性方程组并不需要
行倍乘操作。
我们在这里引入它，是因为它很方便，而且本章稍后我们将把它用于
高斯消元法的一个变体，即
高斯–若尔当消元法。

到目前为止的所有方程组中，方程的个数都与未知数的个数相同。
它们全都有解，而且全都只有一个解。
在本小节结束时，我们来看看
还可能发生哪些别的情况。

\begin{example} \label{ex:MoreEqsThanUnks}
这个方程组
的方程个数多于变量的个数。
\begin{equation*}
    \begin{linsys}{2}
      x  &+  &3y  &=  &1  \\
     2x  &+  &y   &=  &-3 \\
     2x  &+  &2y  &=  &-2
    \end{linsys}
\end{equation*}
高斯消元法也能帮助我们理解这个方程组，因为
\begin{align*}
    &\grstep[-2\rho_1 +\rho_3]{-2\rho_1 +\rho_2}
    \begin{linsys}{2}
       x  &+  &3y  &=  &1  \\
          &   &-5y &=  &-5 \\
          &   &-4y &=  &-4
     \end{linsys}
\end{align*}
表明其中一个方程是多余的。
阶梯形
\begin{equation*}
    \grstep{-(4/5)\rho_2 +\rho_3}
    \begin{linsys}{2}
       x  &+  &3y  &=  &1  \\
          &   &-5y &=  &-5 \\
          &   &0   &=  &0
     \end{linsys}
\end{equation*}
给出 \( y=1 \) 与 \( x=-2 \)。
“\( 0=0 \)”反映了这种多余性。
\end{example}

对于变量个数多于方程个数的方程组，高斯消元法同样有用。
下一小节有很多例子。

线性方程组与上面所示例子的另一种不同之处
在于某些线性方程组没有唯一解。
这可能以两种方式发生。
第一种是方程组可能根本没有解。

\begin{example} \label{ex:MoreEqsThanUnksInconsis}
将上一个例题中的方程组与这个方程组对比。
\begin{equation*}
    \begin{linsys}{2}
      x  &+  &3y  &=  &1  \\
     2x  &+  &y   &=  &-3 \\
     2x  &+  &2y  &=  &0
    \end{linsys}
    \grstep[-2\rho_1 +\rho_3]{-2\rho_1 +\rho_2}
    \begin{linsys}{2}
       x  &+  &3y  &=  &1  \\
          &   &-5y &=  &-5 \\
          &   &-4y &=  &-2
     \end{linsys}
\end{equation*}
这里方程组是无解的：~没有任何数对 $(s_1,s_2)$ 能
同时满足全部三个方程。
阶梯形使这种无解性显而易见。
\begin{equation*}
  \grstep{-(4/5)\rho_2 +\rho_3}
  \begin{linsys}{2}
     x  &+  &3y  &=  &1  \\
        &   &-5y &=  &-5 \\
        &   &0   &=  &2
   \end{linsys}
\end{equation*}
解集为空。
\end{example}

\begin{example}
前面的方程组中方程个数多于未知数的个数，但造成无解的原因并不在于此\Dash
\nearbyexample{ex:MoreEqsThanUnks}
的方程个数多于未知数的个数，但它却是相容的。
反过来，方程个数多于未知数对于无解来说也不是必要的，正如我们看到这个方程个数与未知数个数相同的
无解方程组那样。
\begin{equation*}
  \begin{linsys}{2}
    x  &+  &2y  &=  &8  \\
   2x  &+  &4y  &=  &8
  \end{linsys}
  \grstep{-2\rho_1 + \rho_2}
  \begin{linsys}{2}
     x  &+  &2y  &=  &8  \\
        &   &0   &=  &-8
   \end{linsys}
\end{equation*}
其实，
无解与左右两边的相互作用有关；
在上面第一个方程组中，左边的第二个方程是第一个的两倍，但右边的第二个常数却不是第一个的两倍。
以后我们将进一步讨论方程组各部分之间的
依赖关系。
\end{example}

除了没有解之外，
线性方程组没有唯一解的另一种方式
是有许多解。

\begin{example}
在这个方程组中
\begin{equation*}
  \begin{linsys}{2}
    x  &+  &y   &=  &4  \\
   2x  &+  &2y  &=  &8
  \end{linsys}
\end{equation*}
满足第一个方程的任何数对也
满足第二个方程。
解集
% \( \{ (x,y)\suchthat x+y=4 \} \)
\( \set{ (x,y)\suchthat x+y=4} \)
是无穷的；一些例子
中的数对为~$(0,4)$、$(-1,5)$ 与 $(2.5,1.5)$。

在这里应用高斯消元法的结果与前面的例题形成对比，
因为我们没有得到矛盾方程。
\begin{equation*}
  \grstep{-2\rho_1 + \rho_2}
  \begin{linsys}{2}
    x  &+  &y   &=  &4  \\
       &   &0   &=  &0
   \end{linsys}
\end{equation*}
\end{example}

不要被那个例题迷惑：~一个 $0=0$ 方程
并不是方程组有许多解的信号。

\begin{example}  \label{ex:NoZerosInfManySols}
没有 \( 0=0 \) 方程并不能阻止一个方程组有
许多不同的解。
这个方程组处于阶梯形，
没有 $0=0$，却有无穷多解，
包括 $(0,1,-1)$、
$(0,1/2,-1/2)$、$(0,0,0)$ 与 $(0,-\pi,\pi)$
（任何第一个分量为 $0$ 且第二个分量为第三个的
相反数的三元组都是一个解）。
\begin{equation*}
  \begin{linsys}{3}
     x  &+  &y  &+  &z  &=  &0  \\
        &   &y  &+  &z  &=  &0
   \end{linsys}
\end{equation*}

\( 0=0 \) 的出现也不意味着方程组必定有
许多解。
\nearbyexample{ex:MoreEqsThanUnks} 表明了这一点。
下面这个方程组也是如此，它
根本没有解，尽管
其阶梯形中有一个 $0=0$ 行。
\begin{align*}
  \begin{linsys}{3}
    2x  &   &   &-   &2z  &=  &6  \\
        &   &y  &+   &z   &=  &1  \\
    2x  &+  &y  &-   &z   &=  &7  \\
        &   &3y &+   &3z  &=  &0
  \end{linsys}
  &\grstep{-\rho_1 +\rho_3}
  \begin{linsys}{3}
     2x  &\spaceforemptycolumn   &   &-   &2z  &=  &6  \\
         &   &y  &+   &z   &=  &1  \\
         &   &y  &+   &z   &=  &1  \\
         &   &3y &+   &3z  &=  &0
  \end{linsys}                                    \\
  &\grstep[-3\rho_2 +\rho_4]{-\rho_2 +\rho_3}
  \begin{linsys}{3}
     2x  &\spaceforemptycolumn    &   &-   &2z  &=  &6  \\
         &   &y  &+   &z   &=  &1  \\
         &   &   &    &0   &=  &0  \\
         &   &   &    &0   &=  &-3
   \end{linsys}
\end{align*}
\end{example}

总之，
高斯消元法用行变换
把方程组准备好以便进行回代。
如果任何一步出现矛盾方程，那么我们就可以停下来，
得出方程组无解的结论。
如果我们到达阶梯形而没有矛盾方程，
并且每个变量都是其所在行的
首变量，那么方程组有唯一解，我们可以通过
回代求出它。
最后，如果我们到达阶梯形而没有矛盾方程，
但解不唯一\Dash
也就是说，至少有一个变量不是首变量\Dash
那么方程组有许多解。

下一小节将探讨第三种情形。
我们将看到这样的方程组必定有无穷多解，
并且我们将描述
解集。


\medskip
\noindent\textbf{注。}\hspace*{.2em}
\textit{在这里的练习以及全书其余练习中，
你必须对每个答案给出理由。
例如，若一个问题问某个方程组是否有解，那么你必须
通过给出解来论证“有”的回答，并且必须
通过证明不存在解来论证“无”的回答。}
\begin{exercises}
  \recommended \item 
    用高斯消元法求出每个方程组的唯一解。
    \begin{exparts*}
      \partsitem 
        $\begin{linsys}[t]{2}
          2x  &+  &3y  &=  &13  \\
          x   &-  &y   &=  &-1
        \end{linsys}$
      \partsitem 
        $\begin{linsys}[t]{3}
          x   &  &  &-  &z  &=  &0  \\
          3x  &+ &y &   &   &=  &1  \\
          -x  &+ &y &+  &z  &=  &4
        \end{linsys}$
    \end{exparts*}
    \begin{answer}
      \begin{exparts}
        \partsitem 高斯消元法
          \begin{equation*}
            \grstep{-(1/2)\rho_1+\rho_2}
            \begin{linsys}{2}
               2x  &+  &3y      &=  &13  \\
                   &-  &(5/2)y  &=  &-15/2              
            \end{linsys}
          \end{equation*}
          得到解为$y=3$与$x=2$。
        \partsitem 这里高斯消元法
          \begin{equation*}
            \grstep[\rho_1+\rho_3]{-3\rho_1+\rho_2}
            \begin{linsys}{3}
              x   &  &  &-  &z  &=  &0  \\
                  &  &y &+  &3z &=  &1  \\
                  &  &y &   &   &=  &4
            \end{linsys}
            \grstep{-\rho_2+\rho_3}
            \begin{linsys}{3}
              x   &  &  &-  &z    &=  &0  \\
                  &  &y &+  &3z   &=  &1  \\
                  &  &  &   &-3z  &=  &3
            \end{linsys}
          \end{equation*}
          给出$x=-1$，$y=4$，$z=-1$。
      \end{exparts}
    \end{answer}
  \item
    每个方程组都已是阶梯形。
    对每一个，说明该方程组是唯一解、无解，还是无穷多解。
    \begin{exparts*}
      \partsitem 
        \(\begin{linsys}[t]{2}
           -3x &+ &2y   &= &0  \\          
               &  &-2y  &= &0            
        \end{linsys}\)
      \partsitem 
        \(\begin{linsys}[t]{3}
             x &+ &y &  &   &= &4  \\          
               &  &y &- &z  &= &0  \\           
        \end{linsys}\)
      \partsitem 
        \(\begin{linsys}[t]{3}
             x &+ &y &  &   &= &4  \\          
               &  &y &  &-z  &= &0  \\ 
               &  &  &  &0  &= &0  
        \end{linsys}\)
      \partsitem 
        \(\begin{linsys}[t]{2}
             x &+ &y    &= &4  \\          
               &  &0    &= &4            
        \end{linsys}\)
      \partsitem 
        \(\begin{linsys}[t]{3}
            3x &+ &6y &+  &z  &= &-0.5  \\          
               &  &   &   &-z &= &2.5  \\           
        \end{linsys}\)
      \partsitem 
        \(\begin{linsys}[t]{2}
             x &- &3y    &= &2  \\          
               &  &0    &= &0            
        \end{linsys}\)
      \partsitem 
        \(\begin{linsys}[t]{2}
             2x &+ &2y    &= &4  \\
                &  &y     &= &1 \\          
                 &  &0    &= &4            
        \end{linsys}\)
      \partsitem 
        \(\begin{linsys}[t]{2}
              2x &+ &y   &= &0  
        \end{linsys}\)
      \partsitem 
        \(\begin{linsys}[t]{2}
              x &- &y    &= &-1  \\
                &  &0     &= &0 \\          
                 &  &0    &= &4            
        \end{linsys}\)
      \partsitem 
        \(\begin{linsys}[t]{3}
              x &+ &y  &- &3z  &= &-1  \\
                &  &y  &- &z  &= &2  \\
                &  &  &  &z    &= &0  \\
                &  &  &  &0   &= &0  
        \end{linsys}\)
    \end{exparts*}
    \begin{answer}
      如果方程组含有矛盾方程，那么它无解。
      否则，如果有变量不引领任何一行，那么它有无穷多解。
      最后一种情形是，既没有矛盾方程，而每个变量又都引领某一行，
      这时它有唯一解。
      \begin{exparts}
        \partsitem 唯一解
        \partsitem 无穷多解
        \partsitem 无穷多解
        \partsitem 无解
        \partsitem 无穷多解
        \partsitem 无穷多解
        \partsitem 无解
        \partsitem 无穷多解
        \partsitem 无解
        \partsitem 唯一解
      \end{exparts}
    \end{answer}
  \recommended \item  
    用高斯消元法求解每个方程组，
    或者得出“无穷多解”或“无解”的结论。
    \begin{exparts*}
      \partsitem \(
               \begin{linsys}[t]{2}
                  2x  &+  &2y  &=  &5  \\
                   x  &-  &4y  &=  &0  
               \end{linsys}
             \)
      \partsitem \(
               \begin{linsys}[t]{2}
                  -x  &+  &y   &=  &1  \\
                   x  &+  &y   &=  &2  
               \end{linsys}
             \) 
      \partsitem  \(
               \begin{linsys}[t]{3}
                   x  &-  &3y  &+  &z  &=  &1  \\
                   x  &+  &y   &+  &2z &=  &14 
                \end{linsys}
             \) 
      \partsitem  \(
               \begin{linsys}[t]{2}
                  -x  &-  &y   &=  &1  \\
                 -3x  &-  &3y  &=  &2  
               \end{linsys}
             \) 
      \partsitem  \(
               \begin{linsys}[t]{3}
                      &   &4y  &+  &z  &=  &20 \\
                  2x  &-  &2y  &+  &z  &=  &0  \\
                   x  &   &    &+  &z  &=  &5  \\
                   x  &+  &y   &-  &z  &=  &10 
                \end{linsys}
             \)
      \partsitem \( \begin{linsys}[t]{4}
                 2x  &   &   &+  &z  &+  &w  &=  &5  \\
                     &   &y  &   &   &-  &w  &=  &-1 \\
                 3x  &   &   &-  &z  &-  &w  &=  &0  \\
                 4x  &+  &y  &+  &2z &+  &w  &=  &9  
               \end{linsys}
            \)
    \end{exparts*}
    \begin{answer} 
      \begin{exparts}
       \partsitem 高斯消元
        \begin{equation*}
          \grstep{-(1/2)\rho_1+\rho_2}
          \begin{linsys}{2}
             2x  &+  &2y  &=  &5  \\
                 &   &-5y &=  &-5/2  
          \end{linsys}
        \end{equation*}
        表明\( y=1/2 \)与\( x=2 \)是唯一解。
      \partsitem 高斯消元法
        \begin{equation*}
          \grstep{\rho_1+\rho_2}
          \begin{linsys}{2}
             -x  &+  &y   &=  &1  \\
                 &   &2y  &=  &3  
           \end{linsys}
        \end{equation*}
        给出\( y=3/2 \)与\( x=1/2 \)这个唯一解。
      \partsitem 行化简
        \begin{equation*}
            \grstep{-\rho_1+\rho_2}
            \begin{linsys}{3}
                x  &-  &3y  &+  &z  &=  &1  \\
                   &   &4y  &+  &z  &=  &13 
             \end{linsys}
        \end{equation*}
        由于变量$z$不是任何一行的首变量，
        这表明方程组有无穷多解。
      \partsitem 行化简
        \begin{equation*}
          \grstep{-3\rho_1+\rho_2}
          \begin{linsys}{2}
             -x  &-  &y   &=  &1  \\
                 &   &0   &=  &-1 
           \end{linsys}
        \end{equation*}
        表明方程组无解。
      \partsitem 高斯消元法
        \begin{align*}
            \grstep{\rho_1\leftrightarrow\rho_4}
            \begin{linsys}{3}
                x  &+  &y   &-  &z  &=  &10 \\
               2x  &-  &2y  &+  &z  &=  &0  \\
                x  &   &    &+  &z  &=  &5  \\
                   &   &4y  &+  &z  &=  &20 
             \end{linsys}
            &\grstep[-\rho_1+\rho_3]{-2\rho_1+\rho_2}
            \begin{linsys}{3}
                x  &+  &y   &-  &z  &=  &10 \\
                   &   &-4y &+  &3z &=  &-20\\
                   &   &-y  &+  &2z &=  &-5 \\
                   &   &4y  &+  &z  &=  &20 
             \end{linsys}                                      \\
            &\grstep[\rho_2+\rho_4]{-(1/4)\rho_2+\rho_3}
            \begin{linsys}{3}
                x  &+  &y   &-  &z      &=  &10 \\
                   &   &-4y &+  &3z     &=  &-20\\
                   &   &    &   &(5/4)z &=  &0  \\
                   &   &    &   &4z     &=  &0  
             \end{linsys}
        \end{align*}
        给出唯一解\( (x,y,z)=(5,5,0) \)。
      \partsitem 这里高斯消元法给出
         \begin{align*}
            &\grstep[-2\rho_1+\rho_4]{-(3/2)\rho_1+\rho_3}
            \begin{linsys}{4}
               2x  &\spaceforemptycolumn   &   &+  &z       &+  &w       &=  &5  \\
                   &   &y  &   &        &-  &w       &=  &-1 \\
                   &   &   &-  &(5/2)z  &-  &(5/2)w  &=  &-15/2  \\
                   &   &y  &   &        &-  &w       &=  &-1  
             \end{linsys}                                             \\ 
            &\grstep{-\rho_2+\rho_4}
            \begin{linsys}{4}
               2x  &\spaceforemptycolumn   &   &+  &z       &+  &w       &=  &5  \\
                   &   &y  &   &        &-  &w       &=  &-1 \\
                   &   &   &-  &(5/2)z  &-  &(5/2)w  &=  &-15/2  \\
                   &   &   &   &        &   &0       &=  &0 
             \end{linsys}
         \end{align*}
         这表明方程组有无穷多解。
      \end{exparts} 
    \end{answer}
  \item  
    求解每个方程组，
    或者得出“无穷多解”或“无解”的结论。
    使用高斯消元法。
    \begin{exparts*}
      \partsitem \(
               \begin{linsys}[t]{3}
                   x  &+  &y   &+  &z   &=  &5  \\
                   x  &-  &y   &   &    &=  &0  \\
                      &   &y   &+  &2z  &=  &7  
               \end{linsys}
             \)
      \partsitem \(
               \begin{linsys}[t]{3}
                   3x  &  &    &+  &z   &=  &7  \\
                    x  &-  &y   &+   &3z    &=  &4  \\
                    x   &+  &2y   &-  &5z  &=  &-1  
               \end{linsys}
             \)
      \partsitem \(
               \begin{linsys}[t]{3}
                   x   &+  &3y   &+  &z   &=  &0  \\
                   -x  &-  &y   &   &    &=  &2  \\
                   -x   &+  &y   &+  &2z  &=  &8  
               \end{linsys}
             \)
    \end{exparts*}
    \begin{answer} 
      \begin{exparts}
       \partsitem 
         % sage: M = matrix(RDF, [[1,1,1], [1,-1,0], [0,1,2]])
         % sage: v = vector(RDF, [5,0,7])
         % sage: M_prime = M.augment(v)
         % sage: gauss_method(M_prime)
         % [ 1.0  1.0  1.0  5.0]
         % [ 1.0 -1.0  0.0  0.0]
         % [ 0.0  1.0  2.0  7.0]
         %  take -1.0 times row 1 plus row 2
         % [ 1.0  1.0  1.0  5.0]
         % [ 0.0 -2.0 -1.0 -5.0]
         % [ 0.0  0.0  1.5  4.5]
         %  take 0.5 times row 2 plus row 3
         % [ 1.0  1.0  1.0  5.0]
         % [ 0.0 -2.0 -1.0 -5.0]
         % [ 0.0  0.0  1.5  4.5]
          高斯消元法
          \begin{align*}
           \begin{linsys}{3}
              x  &+  &y   &+  &z   &=  &5  \\
              x  &-  &y   &   &    &=  &0  \\
                 &   &y   &+  &2z  &=  &7  
           \end{linsys}
           &\grstep{-\rho_1+\rho_2} 
           \begin{linsys}{3}
              x  &+  &y   &+  &z   &=  &5  \\
              0  &   &-2y  &-   &z    &=  &-5  \\
                 &   &y   &+  &2z  &=  &7  
           \end{linsys}                                     \\
           &\grstep{(1/2)\rho_2+\rho_3} 
           \begin{linsys}{3}
              x  &+  &y   &+  &z   &=  &5  \\
              0  &   &-2y  &-   &1    &=  &-5  \\
                 &   &    &+  &(3/2)z  &=  &7/2  
           \end{linsys}
          \end{align*}
          再经回代得到$x=1$，$y=1$，$z=3$。
       \partsitem 
         % sage: M = matrix(QQ, [[3,0,1], [1,-1,3], [1,2,-5]])
         % sage: v = vector(RDF, [7,4,-1])
         % sage: v = vector(QQ, [7,4,-1])
         % sage: M_prime = M.augment(v)
         % sage: gauss_method(M_prime)
         % [ 3  0  1  7]
         % [ 1 -1  3  4]
         % [ 1  2 -5 -1]
         %  take -1/3 times row 1 plus row 2
         %  take -1/3 times row 1 plus row 3
         % [    3     0     1     7]
         % [    0    -1   8/3   5/3]
         % [    0     2 -16/3 -10/3]
         %  take 2 times row 2 plus row 3
         % [  3   0   1   7]
         % [  0  -1 8/3 5/3]
         % [  0   0   0   0]
          这里高斯消元法
          \begin{align*}
            \begin{linsys}{3}
                3x  &  &    &+  &z   &=  &7  \\
                 x  &-  &y   &+   &3z    &=  &4  \\
                 x   &+  &2y   &-  &5z  &=  &-1  
            \end{linsys}
            &\grstep[-(1/3)\rho_1+\rho_3]{-(1/3)\rho_1+\rho_2} 
            \begin{linsys}{3}
                3x  &  &    &+  &z        &=  &7  \\
                    &  &-y   &+ &(8/3)z   &=  &5/3  \\
                    &  &2y   &- &(16/3)z  &=  &-(10/3)  
            \end{linsys}                                      \\
            &\grstep{2\rho_2+\rho_3} 
            \begin{linsys}{3}
                3x  &  &    &+  &z        &=  &7  \\
                    &  &-y   &+ &(8/3)z   &=  &5/3  \\
                    &  &     &  &0        &=  &0  
            \end{linsys}
        \end{align*}
        发现变量~$z$不引领任何一行。
        方程组有无穷多解。
        \partsitem
          % sage: M = matrix(QQ, [[1,3,1], [-1,-1,0], [-1,1,2]])
          % sage: v = vector(QQ, [0,2,8])
          % sage: M_prime = M.augment(v)
          % sage: gauss_method(M_prime)
          % [ 1  3  1  0]
          % [-1 -1  0  2]
          % [-1  1  2  8]
          %  take 1 times row 1 plus row 2
          %  take 1 times row 1 plus row 3
          % [1 3 1 0]
          % [0 2 1 2]
          % [0 4 3 8]
          %  take -2 times row 2 plus row 3
          % [1 3 1 0]
          % [0 2 1 2]
          % [0 0 1 4]
          下列步骤
          \begin{equation*}
            \begin{linsys}{3}
               x   &+  &3y   &+  &z   &=  &0  \\
               -x  &-  &y   &   &    &=  &2  \\
               -x   &+  &y   &+  &2z  &=  &8  
            \end{linsys}
            \grstep[\rho_1+\rho_3]{\rho_1+\rho_2}
            \begin{linsys}{3}
               x   &+  &3y   &+  &z   &=  &0  \\
                   &   &2y   &+  &z   &=  &2  \\
                   &   &4y   &+  &3z  &=  &8  
            \end{linsys}
            \grstep{-2\rho_2+\rho_3}
            \begin{linsys}{3}
               x   &+  &3y   &+  &z   &=  &0  \\
                   &   &2y   &+  &z   &=  &2  \\
                   &   &     &   &z   &=  &4  
            \end{linsys}
          \end{equation*}
          给出$(x,y,z)=(-1,-1,4)$。
      \end{exparts} 
    \end{answer}
  \recommended \item 
    我们可以用高斯消元法以外的方法求解线性方程组。
    中学里常教的一种方法是：从某个方程中解出一个变量，
    再把所得的表达式代入其他方程。
    然后重复这一步骤，直到得到只含一个
    变量的方程。
    由此我们得到解中的第一个数，再通过
    回代得到其余的数。
    这个方法比高斯消元法耗时更长，因为它涉及
    更多算术运算，而且也更容易出错。
    为了说明它如何可能导致错误的结论，我们将使用方程组
    \begin{equation*}
      \begin{linsys}{2}
            x  &+  &3y  &=  &1  \\
            2x  &+  &y   &=  &-3 \\
            2x  &+  &2y  &=  &0  
      \end{linsys}
    \end{equation*}
    它来自\nearbyexample{ex:MoreEqsThanUnksInconsis}。
    \begin{exparts}
      \partsitem 从第一个方程解出$x$，
        并把该表达式代入第二个方程。
        求出所得的$y$。
      \partsitem 再次从第一个方程解出$x$，
        但这次把该表达式代入第三个方程。
        求出这个$y$。
    \end{exparts}
    使用这个方法的人必须额外采取什么步骤，
    才能避免错误地得出方程组有解的结论？
    \begin{answer}
      \begin{exparts}
        \partsitem 由$x=1-3y$得到$2(1-3y)+y=-3$，从而$y=1$。
        \partsitem 由$x=1-3y$得到$2(1-3y)+2y=0$，
           从而得出$y=1/2$的结论。
      \end{exparts}
      使用这个方法的人必须把每个可能的解代回所有方程加以检验。
    \end{answer}
  \recommended \item 
    当\( k \)取哪些值时，这个方程组
    无解、有无穷多解，或有唯一解？
    \begin{equation*}
       \begin{linsys}{2}
          x  &-  &y  &=  &1  \\
         3x  &-  &3y &=  &k  
       \end{linsys}
    \end{equation*}
    \begin{answer}
      作消元
      \begin{equation*}
        \grstep{-3\rho_1+\rho_2}
        \begin{linsys}{2}
          x  &-  &y  &=  &1\hfill  \\
             &   &0  &=  &-3+k\hfill  
        \end{linsys}
      \end{equation*}
      由此得出：若\( k\neq 3 \)，则该方程组无解；
      若\( k=3 \)，则它有无穷多解。
      它永远不会有唯一解。  
    \end{answer}
  \item 
    这个方程组不是线性的，因为其中出现的是$\sin\alpha$而不是$\alpha$
    \begin{equation*}
      \begin{linsys}{3}
         2\sin\alpha  &-  &\cos\beta  &+  &3\tan\gamma  &=  &3  \\
         4\sin\alpha  &+  &2\cos\beta &-  &2\tan\gamma  &=  &10  \\
         6\sin\alpha  &-  &3\cos\beta &+  &\tan\gamma   &=  &9  
      \end{linsys}
    \end{equation*}
    然而我们仍可以应用高斯消元法。
    请做一做。
    该方程组有解吗？
    \begin{answer}
      令\( x=\sin\alpha \)，\( y=\cos\beta \)，\( z=\tan\gamma \)：
      \begin{equation*}
        \begin{linsys}{3}
           2x  &-  &y  &+  &3z  &=  &3  \\
           4x  &+  &2y &-  &2z  &=  &10  \\
           6x  &-  &3y &+  &z   &=  &9  
        \end{linsys}
         \grstep[-3\rho_1+\rho_3]{-2\rho_1+\rho_2}
         \begin{linsys}{3}
           2x  &-  &y  &+  &3z  &=  &3  \\
               &   &4y &-  &8z  &=  &4   \\
               &   &   &   &-8z &=  &0  
         \end{linsys}
      \end{equation*}
      得到\( z=0 \)，\( y=1 \)，\( x=2 \)。
      注意，没有\( \alpha\in\R \)满足$\sin\alpha=2$。  
      \end{answer}
  \recommended \item 
    常数（各个$b$）必须满足什么条件，
    才能使下面每个方程组都有解？
    \textit{提示。}
    应用高斯消元法，看看右边会发生什么。
    \begin{exparts*}
      \partsitem  \(
        \begin{linsys}[t]{2}
           x  &-  &3y  &=  &b_1 \\
          3x  &+  &y   &=  &b_2 \\
           x  &+  &7y  &=  &b_3 \\
          2x  &+  &4y  &=  &b_4 
        \end{linsys}   \)
      \partsitem \(
        \begin{linsys}[t]{3}
           x_1  &+  &2x_2  &+  &3x_3  &=  &b_1  \\
          2x_1  &+  &5x_2  &+  &3x_3  &=  &b_2  \\
           x_1  &   &      &+  &8x_3  &=  &b_3  
        \end{linsys}   \)
    \end{exparts*}
    \begin{answer} 
      \begin{exparts}
       \partsitem 高斯消元法
         \begin{equation*}
           \grstep[-\rho_1+\rho_3 \\ -2\rho_1+\rho_4]{-3\rho_1+\rho_2}
           \begin{linsys}{2}
              x  &-  &3y  &=  &b_1\hfill \\
                 &   &10y &=  &-3b_1+b_2\hfill \\
                 &   &10y &=  &-b_1+b_3\hfill \\
                 &   &10y &=  &-2b_1+b_4\hfill 
            \end{linsys}         
           \grstep[-\rho_2+\rho_4]{-\rho_2+\rho_3}
           \begin{linsys}{2}
              x  &-  &3y  &=  &b_1\hfill \\
                 &   &10y &=  &-3b_1+b_2\hfill \\
                 &   &0   &=  &2b_1-b_2+b_3\hfill \\
                 &   &0   &=  &b_1-b_2+b_4\hfill 
            \end{linsys}
         \end{equation*}
         表明该方程组相容当且仅当
         \( b_3=-2b_1+b_2 \)与\( b_4=-b_1+b_2 \)同时成立。
       \partsitem 消元
         \begin{align*}
            &\grstep[-\rho_1+\rho_3]{-2\rho_1+\rho_2}
            \begin{linsys}{3}
              x_1  &+  &2x_2  &+  &3x_3  &=  &b_1\hfill  \\
                   &   &x_2   &-  &3x_3  &=  &-2b_1+b_2\hfill  \\
                   &   &-2x_2 &+  &5x_3  &=  &-b_1+b_3\hfill  
             \end{linsys}                                          \\
            &\grstep{2\rho_2+\rho_3}
            \begin{linsys}{3}
              x_1  &+  &2x_2  &+  &3x_3  &=  &b_1\hfill  \\
                   &   &x_2   &-  &3x_3  &=  &-2b_1+b_2\hfill  \\
                   &   &      &   &-x_3  &=  &-5b_1+2b_2+b_3\hfill  
             \end{linsys}
         \end{align*}
         表明\( b_1 \)、\( b_2 \)与\( b_3 \)各自可以取任何
         实数\Dash 该方程组总有唯一解。
      \end{exparts}   
      \end{answer}
  \item 
    判断真假：未知数多于方程个数的方程组
    至少有一个解。
    （与往常一样，要断言“真”必须给出证明，
    而要断言“假”则必须举出反例。）
    \begin{answer}
      下面这个未知数多于方程个数的方程组
      \begin{equation*}
        \begin{linsys}{3}
          x  &+  &y  &+  &z  &=  &0  \\
          x  &+  &y  &+  &z  &=  &1  
        \end{linsys}
      \end{equation*}
      无解。   
      \end{answer}
  \item 
    任何像本小节开头那样的化学\index{Chemistry problem}问题\Dash
    配平反应的问题\Dash 都必须有无穷多解吗？
    \begin{answer}
      是的。
      例如，同一个反应可以在两个不同的烧瓶中进行，
      这一事实表明任何一个解的两倍是另一个不同的解
      （如果物理上确实发生反应，那么至少必须
      有一个非零解）。
    \end{answer}
  \recommended \item 
    求系数
    \( a \)、\( b \)、\( c \)，使得\( f(x)=ax^2+bx+c \)的图像
    经过点\( (1,2) \)、\( (-1,6) \)和\( (2,3) \)。
    \begin{answer}
      由\( f(1)=2 \)、\( f(-1)=6 \)以及\( f(2)=3 \)我们得到
      一个线性方程组。
      \begin{equation*}
        \begin{linsys}{3}
          1a  &+  &1b  &+  &c  &=  &2  \\
          1a  &-  &1b  &+  &c  &=  &6  \\
          4a  &+  &2b  &+  &c  &=  &3  
         \end{linsys}
      \end{equation*}
      高斯消元法
      \begin{equation*}
         \grstep[-4\rho_1+\rho_3]{-\rho_1+\rho_2}
         \begin{linsys}{3}
            a  &+  &b  &+  &c  &=  &2  \\
               &   &-2b&   &   &=  &4  \\
               &   &-2b&-  &3c &=  &-5 
          \end{linsys}               
         \grstep{-\rho_2+\rho_3}
         \begin{linsys}{3}
            a  &+  &b  &+  &c  &=  &2  \\
               &   &-2b&   &   &=  &4  \\
               &   &   &   &-3c&=  &-9 
           \end{linsys}
      \end{equation*}
      表明解为\( f(x)=1x^2-2x+3 \)。  
      \end{answer}
  \item 在\nearbytheorem{th:GaussMethod}之后我们注意到，
   用~$0$乘某一行是不允许的，因为这可能改变解集。 
   举一个解集为~$S_0$的方程组的例子，使得
   用~$0$乘某一行之后，新方程组的解集为~$S_1$，
   并且$S_0$是$S_1$的真子集，即$S_0\neq S_1$。
   再举一个$S_0=S_1$的例子。 
    \begin{answer}
      这里$S_0=\set{(1,1)}$
      \begin{equation*}
          \begin{linsys}{2}
             x  &+  &y  &=  &2  \\
             x  &-  &y  &=  &0
           \end{linsys}               
          \grstep{0\rho_2}
          \begin{linsys}{2}
             x  &+  &y  &=  &2  \\
                &   &0  &=  &0
           \end{linsys}               
      \end{equation*}
      而$S_1$是一个真超集，因为它
      至少包含两个点：$(1,1)$和~$(2,0)$。
      在这个例子中解集没有改变。  
      \begin{equation*}
          \begin{linsys}{2}
             x  &+  &y  &=  &2  \\
            2x  &+  &2y &=  &4
           \end{linsys}               
          \grstep{0\rho_2}
          \begin{linsys}{2}
             x  &+  &y  &=  &2  \\
                &   &0  &=  &0
           \end{linsys}               
      \end{equation*}
    \end{answer}
  \item 
    高斯消元法的工作方式是对方程组中的方程作组合，
    得到新的方程。
    \begin{exparts}
      \partsitem 我们能否通过一系列
        高斯消元步骤，从这个方程组中的方程导出方程\( 3x-2y=5 \)？
        \begin{equation*}
          \begin{linsys}{2}
             x  &+  &y  &=  &1  \\
            4x  &-  &y  &=  &6
          \end{linsys}
        \end{equation*}
      \partsitem 我们能否通过一系列
        高斯消元步骤，从这个方程组中的方程导出方程\( 5x-3y=2 \)？
        \begin{equation*}
          \begin{linsys}{2}
            2x  &+  &2y &=  &5  \\
            3x  &+  &y  &=  &4
          \end{linsys}
        \end{equation*}
      \partsitem 我们能否通过一系列
        高斯消元步骤，从该方程组中的方程导出\( 6x-9y+5z=-2 \)？
        \begin{equation*}
          \begin{linsys}{3}
            2x  &+  &y  &-  &z  &=  &4  \\
            6x  &-  &3y &+  &z  &=  &5
          \end{linsys}
        \end{equation*}
    \end{exparts}
    \begin{answer} 
       \begin{exparts} 
        \partsitem 能。观察可知，所给方程来自
          \( -\rho_1+\rho_2 \)。
        \partsitem 不能。
          数对\( (1,1) \)满足所给方程。
          然而，这个数对
          不满足该方程组中的第一个方程。
        \partsitem 能。
          要判断所给行是否为\( c_1\rho_1+c_2\rho_2 \)，求解
          联系$x$、$y$、
          $z$的系数与常数的方程组：
          \begin{equation*}
            \begin{linsys}{2}
               2c_1  &+  &6c_2  &=  &6  \\
                c_1  &-  &3c_2  &=  &-9 \\
               -c_1  &+  &c_2   &=  &5  \\
               4c_1  &+  &5c_2  &=  &-2 
             \end{linsys}
          \end{equation*}
          得到$c_1=-3$与$c_2=2$，所以所给行是
          \( -3\rho_1+2\rho_2 \)。
      \end{exparts}  
     \end{answer}
  \item 
    证明：设\( a,b,c,d,e \)是实数
    且\( a\neq 0 \)，若线性方程
    \begin{equation*}
       ax+by=c
    \end{equation*}
    与方程
    \begin{equation*}
       ax+dy=e
    \end{equation*}
    有相同的解集，则它们是同一个方程。
    若\( a=0 \)呢？
    \begin{answer}
      若\( a\neq 0 \)，则第一个方程的解集是
      下面这个。
      \begin{equation*}
        \set{(x,y)\in\Re^2\suchthat 
             \text{$x=(c-by)/a=(c/a)-(b/a)\cdot y$}}
        \tag{$*$} 
      \end{equation*}
      于是，给定~$y$就能算出相应的~$x$。
      取$y=0$得到解$(c/a,0)$；由于第二个
      方程$ax+dy=e$被假定有相同的解集，
      代入其中得到$a(c/a)+d\cdot 0=e$，所以$c=e$。
      在($*$)中取$y=1$得到$a((c-b)/a)+d\cdot 1=e$，
      从而$b=d$。
      因此它们是同一个方程。

      当\( a=0 \)时，两个方程可以不同却仍
      有相同的解集，例如
      \( 0x+3y=6 \)与\( 0x+6y=12 \)。   
     \end{answer}
  \item 
    证明：若\( ad-bc\neq 0 \)，则
    \begin{equation*}
      \begin{linsys}{2}
        ax  &+  &by  &=  &j  \\
        cx  &+  &dy  &=  &k  
      \end{linsys}
    \end{equation*}
    有唯一解。
    \begin{answer}
      我们分三种情形：$a\neq 0$；$a=0$而
      $c\neq 0$；以及$a=0$与$c=0$同时成立。

      第一种情形，我们假设\( a\neq 0 \)。
      那么消元
      \begin{equation*}
        \grstep{-(c/a)\rho_1+\rho_2}
        \begin{linsys}{2}
          ax  &+  &by                  &=  &j \hfill\hbox{} \\
              &   &(-(cb/a)+d)y  &=  &-(cj/a)+k \hfill  
         \end{linsys}
      \end{equation*}
      表明该方程组有唯一解当且仅当
      \( -(cb/a)+d\neq 0   \)；记住\( a\neq 0 \)，因
      此回代给出唯一的\( x \)
      （顺带指出，\( j \)与\( k \)对“存在唯一解”这一
      结论没有任何影响，不过若存在唯一解，
      则它们会影响这个解的值）。
      但\( -(cb/a)+d = (ad-bc)/a \)，而一个分数等于\( 0 \)
      当且仅当它的分子等于\( 0 \)。
      于是，在第一种情形，存在唯一解当且仅当
      $ad-bc\neq 0$。

      在第二种情形，若\( a=0 \)但\( c\neq 0 \)，则我们对换
      \begin{equation*}
        \begin{linsys}{2}
          cx  &+  &dy  &=  &k  \\
              &   &by  &=  &j  
        \end{linsys}
      \end{equation*}
      从而得出：该方程组有唯一解当且仅当
      \( b\neq 0 \)
      （在回代中我们利用情形假设\( c\neq 0 \)来得到唯一的
      \( x \)）。
      但\Dash 在\( a=0 \)而\( c\neq 0 \)时\Dash
      条件“\( b\neq 0 \)”
      等价于条件“\( ad-bc\neq 0 \)”。
      第二种情形到此完成。

      最后，第三种情形，
      若\( a \)与\( c \)都是\( 0 \)，则方程组
      \begin{equation*}
        \begin{linsys}{2}
          0x  &+  &by  &=  &j  \\
          0x  &+  &dy  &=  &k  
        \end{linsys}
      \end{equation*}
      可能无解（如果第二个方程不是第一个方程的倍式），
      也可能有无穷多解（如果第二个方程是第一个方程的倍式，
      那么对每个满足两个方程的\( y \)，任何数对\( (x,y) \)都行），
      但它永远不会有唯一解。
      注意\( a=0 \)与\( c=0 \)给出\( ad-bc=0 \)。  
    \end{answer}
  \recommended \item 
    在方程组
    \begin{equation*}
      \begin{linsys}{2}
         ax  &+  &by  &=  &c  \\
         dx  &+  &ey  &=  &f  
      \end{linsys}
    \end{equation*}
    中，每个方程描述了\( xy \)平面中的一条直线。
    用几何推理证明有三种可能：
    有唯一解，无解，
    以及有无穷多解。
    \begin{answer}
      回忆一下，若两条直线有两个不同的公共点，则
      它们是同一条直线。
      这是因为两点确定一条直线，所以这两个点
      同时确定了这两条直线，
      因而它们是同一条直线。

      于是，两条直线可以有一个公共点（给出唯一解）、
      没有公共点（给出无解），
      或者至少有两个公共点（这使它们成为同一条直线）。  
    \end{answer}
  \item \label{ex:ProveGaussMethod}
    补全\nearbytheorem{th:GaussMethod}的证明。
    \begin{answer}
     对于用非零
     实数$k$乘$\rho_i$的消元操作，我们有：\( (s_1,\ldots,s_n) \)满足
     方程组
     \begin{equation*}
       \begin{linsys}{4}
         a_{1,1}x_1  &+  &a_{1,2}x_2 &+  &\cdots  &+  &a_{1,n}x_n  &=  &d_1  \\
                   &   &          &  &        &   &           &\vdotswithin{=}   \\
        ka_{i,1}x_1  &+  &ka_{i,2}x_2 &+  &\cdots  &+  &ka_{i,n}x_n &=  &kd_i  \\
                   &   &           &   &        &   &            &\vdotswithin{=}   \\
         a_{m,1}x_1  &+  &a_{m,2}x_2 &+  &\cdots  &+  &a_{m,n}x_n  &=  &d_m  
       \end{linsys}
     \end{equation*}
     当且仅当
     % \begin{align*}
     %    a_{1,1}s_1+a_{1,2}s_2+\cdots+a_{1,n}s_n
     %    &=d_1                                              \\
     %    &\alignedvdots                                     \\
     %    \text{and\ } ka_{i,1}s_1+ka_{i,2}s_2+\cdots+ka_{i,n}s_n
     %    &=kd_i                                              \\
     %    &\alignedvdots                                      \\
     %    \text{and\ } a_{m,1}s_1+a_{m,2}s_2+\cdots+a_{m,n}s_n
     %    &=d_m
     % \end{align*}
     $a_{1,1}s_1+a_{1,2}s_2+\cdots+a_{1,n}s_n=d_1$ 
     且\ldots{} $ka_{i,1}s_1+ka_{i,2}s_2+\cdots+ka_{i,n}s_n=kd_i$
     且\ldots{} 
     $a_{m,1}s_1+a_{m,2}s_2+\cdots+a_{m,n}s_n=d_m$，
     这由“满足”的定义得出。
     因为\( k\neq 0 \)，上式成立当且仅当
     % \begin{align*}
     %    a_{1,1}s_1+a_{1,2}s_2+\cdots+a_{1,n}s_n
     %    &=d_1                                              \\
     %    &\alignedvdots                                     \\
     %    \text{and\ } a_{i,1}s_1+a_{i,2}s_2+\cdots+a_{i,n}s_n
     %    &=d_i                                              \\
     %    &\alignedvdots                                      \\
     %    \text{and\ } a_{m,1}s_1+a_{m,2}s_2+\cdots+a_{m,n}s_n
     %    &=d_m
     % \end{align*}
     $a_{1,1}s_1+a_{1,2}s_2+\cdots+a_{1,n}s_n=d_1$
     且\ldots{}
     $a_{i,1}s_1+a_{i,2}s_2+\cdots+a_{i,n}s_n=d_i$
     且\ldots{}
     $a_{m,1}s_1+a_{m,2}s_2+\cdots+a_{m,n}s_n=d_m$
     （这只是把第$i$个方程两边直接消去$k$），
     这说明\( (s_1,\ldots,s_n) \)解方程组
     \begin{equation*}
       \begin{linsys}{4}
         a_{1,1}x_1  &+  &a_{1,2}x_2 &+  &\cdots  &+  &a_{1,n}x_n  &=  &d_1  \\
                     &   &           &   &        &   &            &\vdotswithin{=}   \\
         a_{i,1}x_1  &+  &a_{i,2}x_2 &+  &\cdots  &+  &a_{i,n}x_n  &=  &d_i  \\
                     &   &           &   &        &   &            &\vdotswithin{=}   \\
         a_{m,1}x_1  &+  &a_{m,2}x_2 &+  &\cdots  &+  &a_{m,n}x_n  &=
              &d_m  
         \end{linsys}
     \end{equation*}
     即为所需。

     对于组合操作$k\rho_i+\rho_j$，元组
     \( (s_1,\ldots,s_n) \)满足
     \begin{equation*}
       \begin{linsys}{4}
         a_{1,1}x_1             &+  &\cdots  &+  &a_{1,n}x_n  &=  &d_1\hfill\hbox{} \\
                                &   &        &   &            &\vdotswithin{=}   \\
         a_{i,1}x_1             &+  &\cdots  &+  &a_{i,n}x_n  &=  &d_i\hfill\hbox{} \\
                                &   &        &   &            &\vdotswithin{=}   \\
         (ka_{i,1}+a_{j,1})x_1  &+  &\cdots  &+  &(ka_{i,n}+a_{j,n})x_n
               &=  &kd_i+d_j \hfill \\
                                &   &        &   &            &\vdotswithin{=}   \\
         a_{m,1}x_1             &+   &\cdots  &+  &a_{m,n}x_n  &=  
          &d_m\hfill\hbox{} 
        \end{linsys}
     \end{equation*}
     当且仅当
     $a_{1,1}s_1+\cdots+a_{1,n}s_n=d_1$
     且\ldots{}
     $a_{i,1}s_1+\cdots+a_{i,n}s_n=d_i$
     且\ldots{} 
     $(ka_{i,1}+a_{j,1})s_1+\cdots+(ka_{i,n}+a_{j,n})s_n=kd_i+d_j$
     且\ldots{}
     $a_{m,1}s_1+a_{m,2}s_2+\cdots+a_{m,n}s_n=d_m$，
     同样由“满足”的定义得出。
     从方程~\( j \)中减去\( k \)倍的方程~\( i \)。
     （这正是需要\( i\neq j \)的地方；若\( i=j \)，则上面
     两个\( d_i \)不相等。）
     前面的复合命题成立当且仅当
     $a_{1,1}s_1+\cdots+a_{1,n}s_n=d_1$
     且\ldots{}
     $a_{i,1}s_1+\cdots+a_{i,n}s_n=d_i$
     且\ldots{} $
     (ka_{i,1}+a_{j,1})s_1+\cdots+(ka_{i,n}+a_{j,n})s_n
         -(ka_{i,1}s_1+\cdots+ka_{i,n}s_n)
         =kd_i+d_j-kd_i$
      且\ldots{} $a_{m,1}s_1+\cdots+a_{m,n}s_n=d_m$，
     消去之后，这说明\( (s_1,\ldots,s_n) \)解方程组
     \begin{equation*}
       \begin{linsys}{4}
         a_{1,1}x_1  &+   &\cdots  &+  &a_{1,n}x_n  &=  &d_1  \\
                     &    &        &   &            &\vdotswithin{=}   \\
         a_{i,1}x_1  &+   &\cdots  &+  &a_{i,n}x_n  &=  &d_i  \\
                     &    &        &   &            &\vdotswithin{=}   \\
         a_{j,1}x_1  &+  &\cdots  &+  &a_{j,n}x_n  &=  &d_j  \\
                     &   &        &   &            &\vdotswithin{=}   \\
         a_{m,1}x_1  &+  &\cdots  &+  &a_{m,n}x_n  &=
              &d_m\hfill\hbox{}
       \end{linsys}
     \end{equation*}  
     即为所需。
   \end{answer}
  \item 
    存在解集为整个\( \Re^2 \)的
    两未知数线性方程组吗？
    \begin{answer}
      存在，下面这个只有一个方程的方程组：
      \begin{equation*}
         0x+0y=0
      \end{equation*}
      每个\( (x,y)\in\Re^2 \)都满足它。  
    \end{answer}
  \recommended \item 
    高斯消元法中用到的操作有多余的吗？
    也就是说，我们能否用其他操作的组合来实现某个操作？
    \begin{answer}
      有。
      下面这串操作对换了行\( i \)与行\( j \)
      \begin{equation*}
         \grstep{\rho_i+\rho_j}
         \repeatedgrstep{-\rho_j+\rho_i}
         \repeatedgrstep{\rho_i+\rho_j}
         \repeatedgrstep{-1\rho_i}
      \end{equation*}  
      所以在另外两种操作存在的情况下，对换操作是多余的。
     \end{answer}
  \item 
    证明高斯消元法的每种操作都是可逆的。
    也就是说，证明若两个方程组通过一个行变换相联系
    $S_1\rightarrow S_2$，则存在一个行变换变回去
    $S_2\rightarrow S_1$。
    \begin{answer}
      行对换$\rho_i\leftrightarrow\rho_j$可以用
      对换回去$\rho_j\leftrightarrow\rho_i$来逆转。
      $k\rho_i$这一步是用\( k\neq 0  \)乘
      某一行的两边，
      可以用除以~$(1/k)\rho_i$来逆转。

      倍加的情形是非平凡的。
      操作$k\rho_i+\rho_j$
      得到如下的第$j$行。
      \begin{equation*}
        k\cdot a_{i,1}+a_{j,1}+\cdots+k\cdot a_{i,n}+a_{j,n}=k\cdot d_i+d_j
      \end{equation*}
      由于$i\neq j$的限制，第$i$行不变。
      因为第$i$行不变，操作$-k\rho_i+\rho_j$
      把第$j$行变回原状。

      （注意$k\rho_i+\rho_j$中\( i=j \)的条件
       是必需的，否则可能出现
       \begin{equation*}
         \begin{linsys}{2}
           3x  &+  &2y  &=  &7  
         \end{linsys}
         \grstep{2\rho_1+\rho_1}
         \begin{linsys}{2}
           9x  &+  &6y  &=  &21 
          \end{linsys}                 
         \grstep{-2\rho_1+\rho_1}
         \begin{linsys}{2}
          -9x  &-  &6y  &=  &-21 
         \end{linsys}
       \end{equation*}
       于是结论不成立。）
    \end{answer}
  \puzzle \item  
    \cite{Anton}
    一个盒子装着1美分、5美分和10美分的硬币，
    共十三枚，总价值为\( 83 \)美分。
    盒子中每种硬币各有多少枚？
    （这些是美国硬币；
    1美分硬币值$1$~美分，5美分硬币值$5$~美分，
    10美分硬币值$10$~美分。）
    \begin{answer}
      设\( p \)、\( n \)、\( d \)分别是
      1美分、5美分和10美分硬币的个数。
      当变量取实数值时，方程组
      \begin{equation*}
         \begin{linsys}{3}
              p  &+ &n   &+  &d   &=  &13   \\
              p  &+ &5n  &+  &10d &=  &83    
         \end{linsys}
         \grstep{-\rho_1+\rho_2}
         \begin{linsys}{3}
              p  &+ &n   &+  &d   &=  &13   \\
                 &  &4n  &+  &9d  &=  &70    
          \end{linsys}
      \end{equation*}
      的解不止一个；事实上，它有无穷多个解。
      然而，\( p \)、\( n \)、
      \( d \)为非负整数的解只有有限多个。
      逐一检查\( d=0 \)，\ldots，\( d=8 \)可知
      \( (p,n,d)=(3,4,6) \)
      是唯一使用自然数的解。  
    \end{answer}
  \puzzle \item 
    \cite{ContestProb1955no38}
    给定四个正整数。
    任取其中三个数，求它们的算术平均，
    并把这个结果加到第四个整数上。
    这样得到数29、23、21和17。
    原来的整数之一是：
    \begin{exparts*}
      \partsitem 19
      \partsitem 21
      \partsitem 23
      \partsitem 29
      \partsitem 17
    \end{exparts*}
    \begin{answer}
      求解方程组
      \begin{equation*}
        \begin{linsys}{2}
        (1/3)(a+b+c)  &+  &d  &=  &29  \\
        (1/3)(b+c+d)  &+  &a  &=  &23  \\
        (1/3)(c+d+a)  &+  &b  &=  &21  \\
        (1/3)(d+a+b)  &+  &c  &=  &17
        \end{linsys}
      \end{equation*}
      我们得到$a=12$，$b=9$，$c=3$，$d=21$。
      因此第二项，21，是正确答案。
     \end{answer}
  \puzzle \item  
      \cite{Monthly35p47}
      笑一笑这个：\( \mbox{AHAHA}+\mbox{TEHE}=\mbox{TEHAW} \)。
      它来自把一个简单的加法算式中的每个数字换成代码字母，
      要求辨认出这些字母并证明解唯一。
      \begin{answer}
        \answerasgiven
        比较这个加法竖式的个位列与百位列可知，
        十位列必有进位。
        十位列于是告诉我们\( A<H \)，所以
        个位列与百位列都不会有进位。
        五个列于是给出下面五个方程。
        \begin{align*}
          A+E  &=  W  \\
          2H   &=  A+10  \\
          H    &=  W+1  \\
          H+T  &=  E+10  \\
          A+1  &=  T
        \end{align*}
        这五个未知数的五个线性方程若联立求解，
        给出唯一解：\( A=4 \)，\( T=5 \)，\( H=7 \)，
        \( W=6 \)与\( E=2 \)，于是原来的加法算式
        是\( 47474+5272=52746 \)。  
      \end{answer}
  \puzzle \item 
     \cite{Wohascum2}
     有20名成员的Wohascum县委员会
     最近需要选举一位主席。
     有三名候选人（$A$、$B$和$C$）；每张选票上
     要按偏好次序列出这三名
     候选人，不允许弃权。
     结果发现，11名成员（多数）偏好$A$胜过$B$
     （于是其余9名偏好$B$胜过$A$）。
     类似地，发现12名成员偏好$C$胜过$A$。
     鉴于这些结果，有人建议$B$退出，以便
     在$A$与$C$之间进行决选。
     然而$B$表示抗议，随后又发现14名成员偏好
     $B$胜过$C$！
     委员会至今仍未从由此产生的混乱中恢复过来。
     已知$A$、$B$、$C$的每种可能次序都至少出现在
     一张选票上，问有多少成员把$B$列为第一选择？
     \begin{answer}
       \answerasgiven{} 
       \textit{其中一些补充材料取自\cite{joriki}。}
       有八名委员投票给$B$。
       为说明这一点，我们将利用已知信息来研究有多少投票者
       选择了$A$、$B$、$C$的每种次序。

       六种偏好次序是$ABC$、$ACB$、$BAC$、$BCA$、$CAB$、
       $CBA$；假设它们分别获得$a$、$b$、$c$、$d$、$e$、$f$票。
       我们知道
       \begin{equation*}
         \begin{linsys}{3}
           a  &+  &b  &+  &e  &=  &11  \\
           d  &+  &e  &+  &f  &=  &12  \\
           a  &+  &c  &+  &d  &=  &14
         \end{linsys}
       \end{equation*}
       它们分别来自偏好$A$胜过$B$的人数、偏好
       $C$胜过$A$的人数，以及偏好$B$胜过$C$的人数。
       因为总共投了20票，
       所以$a+b+\cdots+f=20$。
       从两倍的前一个方程中减去上面三个方程之和，
       得到$b+c+f=3$。
       我们已规定每种偏好次序至少获得
       一票，所以这意味着$b=c=f=1$。 

       由上面三个方程，完整解为
       $a=6$，$b=1$，$c=1$，$d=7$，$e=4$以及~$f=1$，
       这可以用高斯消元法求出。
       把$B$列为第一选择的委员人数
       因此是$c+d=1+7=8$。

       \par\noindent {\em 评注。}
       即使我们只知道{\em 至少\/}有14名委员偏好$B$胜过$C$，
       这个问题的答案也会一样。

       委员们的偏好看似悖论的性质
       （$A$优于$B$，$B$优于$C$，而$C$又
       优于$A$）是“非传递支配”的一个例子，
       这种情形在汇总个人选择时很常见。
     \end{answer}
  \puzzle \item   
     \cite{Monthly63p93}
    “这个由\( n \)个线性方程组成的、含
     \( n \)个未知数的方程组，”大数学家说，“有一个奇特的
     性质。”

     “天哪！”可怜的菜鸟说，“是什么？”

     “注意，”大数学家说，“常数项构成
     等差数列。”

     “你一解释我就全明白了！”可怜的菜鸟说。
     “你是说像\( 6x+9y=12 \)和\( 15x+18y=21 \)这样？”

     “正是，”大数学家说着，掏出了他的巴松管。
     “的确，这个方程组有唯一解。
     你能找出来吗？”

     “天哪！”可怜的菜鸟喊道，“我被难住了。”

     你呢？
     \begin{answer}
       \answerasgiven
       \textit{我们还没有用到“相依”一词；
       它在这里的意思是高斯
       消元法表明解不唯一。}
       若\( n\geq 3 \)，则方程组相依且解不
       唯一。
       因此\( n<3 \)。
       但“方程组”一词意味着\( n>1 \)。
       因此\( n=2 \)。
       若方程为
       \begin{equation*}
         \begin{linsys}{2}
              ax  &+ &(a+d)y  &=  &a+2d  \\
         (a+3d)x  &+ &(a+4d)y &=  &a+5d  
         \end{linsys}
       \end{equation*}
       则\( x=-1 \)，\( y=2 \)。  
    \end{answer}
\end{exercises}







\subsection{解集的描述}
具有唯一解的线性方程组，其解集只有一个元素。
无解的线性方程组，其解集是空集。
在这两种情形下，解集都很容易描述。
只有当解集包含许多元素时，描述解集才是一个挑战。

\begin{example}
这个方程组有很多解，因为化为阶梯形后
\begin{align*}
  \begin{linsys}{3}
    2x  &   &   &+  &z  &=  &3 \\
     x  &-  &y  &-  &z  &=  &1 \\
    3x  &-  &y  &   &   &=  &4 
  \end{linsys}
  &\grstep[-(3/2)\rho_1 +\rho_3]{-(1/2)\rho_1+\rho_2}
  \begin{linsys}{3}
     2x  &   &   &+  &z      &=  &3    \\
         &   &-y &-  &(3/2)z &=  &-1/2 \\
         &   &-y &-  &(3/2)z &=  &-1/2 
   \end{linsys}                                   \\
  &\grstep{-\rho_2+\rho_3}
  \begin{linsys}{3}
     2x  &   &   &+  &z      &=  &3    \\
         &   &-y &-  &(3/2)z &=  &-1/2 \\
         &   &   &   &0      &=  &0    
   \end{linsys}
\end{align*}
并非所有变量都是首变量。
\nearbytheorem{th:GaussMethod}表明，一个$(x,y,z)$
满足第一个方程组当且仅当它满足第三个方程组。
于是我们可以将解集
$\set{(x,y,z)\suchthat\text{$2x+z=3$ 且 $x-y-z=1$ 且 $3x-y=4$}}$
描述成下面的形式。
\begin{equation*}
  \set{(x,y,z)\suchthat\text{$2x+z=3$ 且 $-y-3z/2=-1/2$}}
  \tag{$*$}
\end{equation*}
这个描述更好，因为它用两个方程代替了三个方程，
但它还不是最优的，
因为变量之间仍有一些难以理解的相互作用。

为了改进它，用不引领任何方程的变量$z$来描述
起引领作用的变量$x$和$y$。
第二个方程给出
$y=(1/2)-(3/2)z$，
第一个方程给出
$x=(3/2)-(1/2)z$。
于是我们可以把解集描述成下面这个三元组的集合。
\begin{equation*}
  \set{((3/2)-(1/2)z,\,(1/2)-(3/2)z,\,z)\suchthat z\in\Re}
  \tag{$**$}
\end{equation*}
\end{example}

与($*$)相比，
($**$)的优点
是$z$可以取任何实数。
这使得判断哪些元组在解集之中的工作
容易得多。
例如，取$z=2$就表明$(1/2,-5/2,2)$是一个解。

\begin{definition} \label{df:FreeVars}
%<*df:FreeVars>
在一个阶梯形的线性方程组中，不是首变量的变量
称为
\definend{自由变量}。\index{echelon form!free variable}\index{free variable}
%</df:FreeVars>
\end{definition}

\begin{example}   \label{ex:Parametrize2}
对一个线性方程组进行消元，结束时可能有不止一个变量是自由的。
对这个方程组用高斯消元法
\begin{align*}
   \begin{linsys}{4}
               x  &+  &y   &+  &z   &-  &w   &=  &1  \\
                  &   &y   &-  &z   &+  &w   &=  &-1 \\
              3x  &   &    &+  &6z  &-  &6w  &=  &6  \\
                  &   &-y  &+  &z   &-  &w   &=  &1  
   \end{linsys}
  &\grstep{-3\rho_1 +\rho_3}
  \begin{linsys}{4}
     x  &+  &y   &+  &z   &-  &w   &=  &1  \\
        &   &y   &-  &z   &+  &w   &=  &-1 \\
        &   &-3y &+  &3z  &-  &3w  &=  &3  \\
        &   &-y  &+  &z   &-  &w   &=  &1  
  \end{linsys}                                      \\
  &\grstep[\rho_2 +\rho_4]{3\rho_2 +\rho_3}
  \begin{linsys}{4}
     x  &+  &y   &+  &z   &-  &w   &=  &1  \\
        &   &y   &-  &z   &+  &w   &=  &-1 \\
        &   &    &   &    &   &0   &=  &0  \\
        &   &    &   &    &   &0   &=  &0  
   \end{linsys}
\end{align*}
留下\( x \)和\( y \)为首变量，而\( z \)和\( w \)都是自由变量。
为了得到我们更喜欢的描述，我们从底部往上做。
我们先把首变量$y$用
$z$和$w$表示，即$y=-1+z-w$。
再往上看最上面的方程，
代入$y$得到
$x+(-1+z-w)+z-w=1$，解出$x$得$x=2-2z+2w$。
解集
\begin{equation*}
   \set{(2-2z+2w,-1+z-w,z,w)\suchthat z,w\in\Re}
  \tag{$**$}
\end{equation*}
把首变量用自由变量表示了出来。
\end{example}

\begin{example}    \label{ex:Parametrize1}
首变量的列表可能会跳过某些列。
经过这样的消元
\begin{align*}
  \begin{linsys}{4}
              2x  &-  &2y  &   &    &   &    &=  &0  \\
                  &   &    &   &z   &+  &3w  &=  &2  \\
              3x  &-  &3y  &   &    &   &    &=  &0  \\
               x  &-  &y   &+  &2z  &+  &6w  &=  &4  
  \end{linsys}
  &\grstep[-(1/2)\rho_1+ \rho_4]{-(3/2)\rho_1 +\rho_3}
  \begin{linsys}{4}
     2x  &-  &2y  &\spaceforemptycolumn   &    &   &    &=  &0  \\
         &   &    &   &z   &+  &3w  &=  &2  \\
         &   &    &   &    &   &0   &=  &0  \\
         &   &    &   &2z  &+  &6w  &=  &4  
   \end{linsys}                                    \\
  &\grstep{-2\rho_2 +\rho_4}
  \begin{linsys}{4}
     2x  &-  &2y  &\spaceforemptycolumn   &    &   &    &=  &0  \\
         &   &    &   &z   &+  &3w  &=  &2  \\
         &   &    &   &    &   &0   &=  &0  \\
         &   &    &   &    &   &0   &=  &0  
   \end{linsys}
\end{align*}
$x$和$z$是首变量，而不是$x$和$y$。
自由变量是$y$和$w$，因此我们可以把解集描述为
$\set{ (y,y,2-3w,w)\suchthat y,w\in\Re }$。
例如，\( (1,1,2,0) \)满足该方程组\Dash取
$y=1$和$w=0$。
四元组\( (1,0,5,4) \)不是解，
因为它的第一个分量不等于第二个分量。
\end{example}

%<*df:Parameter>
我们用来描述一族解的变量是一个\definend{参数}。\index{parameter}  
%</df:Parameter>
我们说前面的例题中的解集
是用$y$和$w$
来\definend{参数化\/}\index{parametrized} 
的。

“参数”和“自由变量”这两个术语的含义并不相同。
在前面的例题中
$y$和$w$是自由变量，因为在阶梯形方程组中它们
不引领任何方程。
它们之所以是参数，是因为
我们用它们来描述解的集合。
假如我们改而
把第二个方程改写成$w=2/3-(1/3)z$，那么
自由变量仍然是$y$和$w$，但参数
将是$y$和$z$。

在本书余下的部分中
我们将通过把线性方程组化为
阶梯形、然后用自由变量进行参数化来求解线性方程组。

\begin{example}
这是另一个有无穷多解的方程组。
\begin{align*}
   \begin{linsys}{4}
               x  &+  &2y  &   &   &   &   &=  &1  \\
              2x  &   &    &+  &z  &   &   &=  &2  \\
              3x  &+  &2y  &+  &z  &-  &w  &=  &4  
   \end{linsys}
  &\grstep[-3\rho_1 +\rho_3]{-2\rho_1+\rho_2}
  \begin{linsys}{4}
     x  &+  &2y  &   &   &   &   &=  &1  \\
        &   &-4y &+  &z  &   &   &=  &0  \\
        &   &-4y &+  &z  &-  &w  &=  &1  
   \end{linsys}                                    \\
  &\grstep{-\rho_2+\rho_3}
  \begin{linsys}{4}
     x  &+  &2y  &   &   &   &   &=  &1  \\
        &   &-4y &+  &z  &   &   &=  &0  \\
        &   &    &   &   &   &-w &=  &1  
   \end{linsys}
\end{align*}
首变量是\( x \)、\( y \)和\( w \)。
变量\( z \)是自由的。
注意，虽然有无穷多
解，但$w$的值并不变化，而是恒为$w=-1$。
为了参数化，把\( w \)用\( z \)写成\( w=-1+0z \)。
于是\( y=(1/4)z \)。
把\( y \)代入第一个
方程得到\( x=1-(1/2)z \)。
解集是$\set{(1-(1/2)z,(1/4)z,z,-1)\suchthat z\in\Re}$。
\end{example}

对解集进行参数化表明，带自由变量的方程组
有无穷多解。
例如，上面的$z$取遍无穷多个实数值，
每个值对应一个不同的解。

我们以建立一套简化的记法来结束本小节，
这套记法用于线性方程组
及其解集。

\begin{definition}\label{df:matrix}
%<*df:matrix>
一个\( \nbym{m}{n} \) \definend{矩阵}\index{matrix}
是数的矩形阵列，
有\( m \)~\definend{行}\index{matrix!row}\index{row} 
和\( n \)~\definend{列}\index{matrix!column}\index{column}。
矩阵中的每个数都是一个
\definend{元素}\index{matrix!entry}\index{entry, matrix}。
%</df:matrix>
\end{definition}

我们通常用一个大写罗马字母来表示矩阵。
例如，
\begin{equation*}
  A=
  \begin{mat}[r]
    1  &2.2  &5  \\
    3  &4    &-7
  \end{mat}
\end{equation*}
它有$2$~行和$3$~列，因此
是一个\( \nbym{2}{3} \)矩阵。
把它读作“二乘三”；
行的个数总是先说。
（矩阵围有括号，
这样当
两个矩阵相邻时，
我们能分辨出一个在哪里结束、另一个从哪里开始。）
我们用对应的小写字母命名矩阵的元素，
于是上面阵列中第二行第一列的
元素是\( a_{2,1}=3 \)。
注意下标的次序很重要：
$a_{1,2}\neq a_{2,1}$，因为\( a_{1,2}=2.2 \)。
我们把
所有\( \nbym{m}{n} \)
矩阵组成的集合记作\( \matspace_{\nbym{m}{n}} \)。\index{matrices|set of}

我们用矩阵做高斯消元法，其方式与
我们对方程组所用的方式基本相同：
矩阵一行的
\definend{首主元}\index{echelon form!leading entry}%
\index{leading!entry} %
是它的第一个非零元素（如果有的话），我们进行行变换来达到
\definend{矩阵阶梯形}，\index{echelon form!matrix}%
\index{matrix!echelon form} %
即下面各行的首主元位于其上方各
行首主元的右侧。
我们喜欢矩阵记法，因为它减轻了
抄写工作的负担，即抄写变量以及
书写$+$号和$=$号的负担。

\begin{example}
我们可以将这个线性方程组
\begin{equation*}
  \begin{linsys}{3}
    x  &+  &2y  &   &    &=  &4   \\
       &   &y   &-  &z   &=  &0   \\
    x  &   &    &+  &2z  &=  &4   
  \end{linsys}
\end{equation*}
缩写成这个矩阵。
\begin{equation*}
    \begin{amat}[r]{3}
      1  &2  &0  &4  \\
      0  &1  &-1 &0  \\
      1  &0  &2  &4
    \end{amat}
\end{equation*}
竖线提醒读者方程组左边的
系数与右边的常数之间的区别。
带竖线的矩阵是
\definend{增广\/}\index{matrix!augmented}\index{augmented matrix}矩阵。
\begin{equation*}
    \begin{amat}[r]{3}
      1  &2  &0  &4  \\
      0  &1  &-1 &0  \\
      1  &0  &2  &4
    \end{amat}
  \grstep{-\rho_1 +\rho_3}
  \begin{amat}[r]{3}
       1  &2  &0  &4  \\
       0  &1  &-1 &0  \\
       0  &-2 &2  &0
     \end{amat}                        
  \grstep{2\rho_2 +\rho_3}
  \begin{amat}[r]{3}
       1  &2  &0  &4  \\
       0  &1  &-1 &0  \\
       0  &0  &0  &0
     \end{amat}
\end{equation*}
第二行表示$y-z=0$，第一行表示
$x+2y=4$，因此解集是
\( \set{(4-2z,z,z)\suchthat z\in\Re} \)。
\end{example}

矩阵记法也
使解集的描述更加清晰。
\nearbyexample{ex:Parametrize2}的
$\set{(2-2z+2w,-1+z-w,z,w)\suchthat z,w\in\Re}$ 
很难读。
我们将把它改写，使所有常数放在一起，\( z \)的所有系数放在一起，以及\( w \)的所有系数
放在一起。
我们把它们竖着写成一列的矩阵。
\begin{equation*}
  \set{\colvec[r]{2 \\ -1 \\ 0 \\ 0}
       +\colvec[r]{-2 \\ 1 \\ 1 \\ 0}\cdot z
       +\colvec[r]{2 \\ -1 \\ 0 \\ 1}\cdot w
       \suchthat z,w\in\Re}
\end{equation*}
例如，最上面一行说明\( x=2-2z+2w \)，
第二行说明\( y= -1+z-w \)。
（下一节给出一个几何解释，它将帮助
我们想象解集。）

\begin{definition}\label{df:vector}
%<*df:vector>
一个
\definend{列向量}，\index{column!vector}\index{vector!column}
通常简称\definend{向量}，\index{vector} 
是只有一列的矩阵。
只有一行的矩阵是
\definend{行向量}\index{row!vector}\index{vector!row}。
向量的元素有时称为
\definend{分量}\index{component of a vector}\index{vector!component}。
分量全为零的列向量或行向量是一个
\definend{零向量}。\index{zero vector}\index{vector!zero}
%</df:vector>
\end{definition}

向量是用大写罗马字母表示矩阵这一惯例的一个例外。
我们使用带上箭头的小写罗马字母或希腊字母：
\( \vec{a} \)、\( \vec{b} \)、\ldots\, 或
\( \vec{\alpha} \)、\( \vec{\beta} \)、\ldots\
（粗体也很常见：
{\boldmath \( a \)}或{\boldmath \( \alpha \)}）。
例如，这是一个第三个分量为\( 7 \)的列向量。
\begin{equation*}
  \vec{v}=
  \colvec[r]{ 1  \\  3  \\ 7}
\end{equation*}
零向量记作\( \zero \)。
零向量有很多种\Dash高度为一的零向量、高度为二的零向量，等等\Dash但
尽管如此，我们仍常说“这个”零向量，
并期望由上下文可以清楚地知道其大小。

\begin{definition}
带有未知数\( x_1,\ldots\,,x_n \)的线性方程
\( a_1x_1+\,\cdots\,+a_nx_n=d \)
被
\begin{equation*}
  \vec{s}=\colvec{s_1 \\ \vdotswithin{s_1} \\ s_n}
\end{equation*}
\definend{满足}\index{vector!satisfies an equation}%
\index{linear equation!satisfied by a vector}，是指
\( a_1s_1+\,\cdots\,+a_ns_n=d \)。
若一个向量满足线性方程组中的每一个方程，则称它满足该线性方程组。
\end{definition}

我们所使用的解集描述方式
涉及把向量相加，
以及用实数去乘它们。
在给出展示这种描述方式的例子之前
我们首先需要定义这两种运算。

\begin{definition} \label{df:VectorSum}
%<*df:VectorSum>
\( \vec{u} \)与\( \vec{v} \)的\definend{向量和}\index{vector!sum}\index{sum!vector}\index{addition of vectors} 
是由各分量之和组成的向量。
\begin{equation*}
  \vec{u}+\vec{v}=
  \colvec{u_1 \\ \vdotswithin{u_1} \\ u_n}
   +
  \colvec{v_1 \\ \vdotswithin{v_1} \\ v_n}
   =
  \colvec{u_1+v_1 \\ \vdotswithin{u_1+v_1} \\ u_n+v_n}
\end{equation*}
%</df:VectorSum>
\end{definition}

注意，要使加法有定义，
各向量必须有相同个数的元素。
这种逐元素的加法对任意一对矩阵都适用，而不只是对向量，
只要它们有相同的行数和
列数。\index{matrix!sum}\index{sum!matrix}

\begin{definition} \label{df:VectorScalarMultiplication}
%<*df:VectorScalarMultiplication>
实数
\( r \)与向量\( \vec{v} \)的\definend{标量乘法\/}\index{vector!scalar multiple}%
\index{scalar multiple!vector}是由各倍数组成的向量。
\begin{equation*}
  r\cdot\vec{v}=
  r\cdot\colvec{v_1 \\ \vdotswithin{v_1} \\ v_n}
  =
  \colvec{rv_1 \\ \vdotswithin{rv_1} \\ rv_n}
\end{equation*}
%</df:VectorScalarMultiplication>
\end{definition}

与加法运算一样，逐元素的标量乘法
运算也超出向量的范围而适用于任意
矩阵。\index{matrix!scalar multiplication}\index{scalar multiplication!matrix}

标量乘法可以写成\( r\cdot\vec{v} \)或
\( \vec{v}\cdot r \)，有时甚至省略“$\cdot$”符号：~$r\vec{v}$。
（不要把标量乘法
称为“标量积”，因为那个名称指的是另一种运算。）

\begin{example}
\begin{equation*}
  \colvec[r]{2 \\ 3 \\ 1}
   +
  \colvec[r]{3 \\ -1 \\ 4}
  =
  \colvec{2+3 \\ 3-1 \\ 1+4}
   =
  \colvec[r]{5 \\ 2 \\ 5}
  \qquad
  7\cdot\colvec[r]{1 \\ 4 \\ -1 \\ -3}
  =
  \colvec[r]{7 \\ 28 \\ -7 \\ -21}
\end{equation*}
\end{example}

注意，加法与标量乘法的定义在重合之处是一致的；例如，\( \vec{v} +\vec{v} = 2\vec{v} \)。

有了这些定义，我们就可以用矩阵与向量记法
既求解方程组又表示解了。

\begin{example} \label{ex:ManyParamsInfManySolsSystem}
这个方程组
\begin{equation*}
   \begin{linsys}{5}
      2x  &+  &y  &  &  &-  &w  &   &   &=  &4  \\
          &   &y  &  &  &+  &w  &+  &u  &=  &4  \\
       x  &   &   &- &z &+  &2w &   &   &=  &0  
   \end{linsys}
\end{equation*}
按如下方式化简。
\begin{align*}
  \begin{amat}[r]{5}
    2  &1  &0  &-1  &0  &4  \\
    0  &1  &0  &1   &1  &4  \\
    1  &0  &-1 &2   &0  &0
  \end{amat}
  &\grstep{-(1/2)\rho_1+\rho_3}
  \begin{amat}[r]{5}
    2  &1     &0  &-1    &0  &4  \\
    0  &1     &0  &1     &1  &4  \\
    0  &-1/2  &-1 &5/2   &0  &-2
  \end{amat}                                 \\
  &\grstep{(1/2)\rho_2+\rho_3}
  \begin{amat}[r]{5}
    2  &1     &0  &-1    &0    &4  \\
    0  &1     &0  &1     &1    &4  \\
    0  &0     &-1 &3     &1/2  &0
  \end{amat}
\end{align*}
解集是
\( \set{(w+(1/2)u,4-w-u,3w+(1/2)u,w,u)\suchthat w,u\in\Re} \)。
我们把它写成向量形式。
\begin{equation*}
  \set{\colvec{x \\ y \\ z \\ w \\ u}=
       \colvec[r]{0 \\ 4 \\ 0 \\ 0 \\ 0}+
       \colvec[r]{1 \\ -1 \\ 3 \\ 1 \\ 0}w+
       \colvec[r]{1/2 \\ -1 \\ 1/2 \\ 0 \\ 1}u
       \suchthat w,u\in\Re}
\end{equation*}
注意向量记法把
每个参数的系数区分得多么清楚。
例如，向量形式的第三行清楚地表明，如果\( u \)固定，那么\( z \)增长的速度是\( w \)的三倍。
同样清楚表明的是，令\( w \)和\( u \)都为零
得到
\begin{equation*}
  \colvec{x \\ y \\ z \\ w \\ u}
  =\colvec[r]{0 \\ 4 \\ 0 \\ 0 \\ 0}
\end{equation*}
是线性方程组的一个特解。
\end{example}

\begin{example}
同样地，方程组
\begin{equation*}
   \begin{linsys}{3}
     x  &-  &y  &+  &z  &=  &1  \\
    3x  &   &   &+  &z  &=  &3  \\
    5x  &-  &2y &+  &3z &=  &5  
  \end{linsys}
\end{equation*}
化简
\begin{align*}
  \begin{amat}[r]{3}
    1  &-1  &1  &1  \\
    3  &0   &1  &3  \\
    5  &-2  &3  &5
  \end{amat}
  &\grstep[-5\rho_1+\rho_3]{-3\rho_1+\rho_2}
  \begin{amat}[r]{3}
    1  &-1  &1  &1  \\
    0  &3   &-2 &0  \\
    0  &3   &-2 &0
  \end{amat}                                    \\
  &\grstep{-\rho_2+\rho_3}
  \begin{amat}[r]{3}
    1  &-1  &1  &1  \\
    0  &3   &-2 &0  \\
    0  &0   &0  &0
  \end{amat}
\end{align*}
后给出一个单参数的解集。
\begin{equation*}
  \set{\colvec[r]{1 \\ 0 \\ 0}
       +\colvec[r]{-1/3 \\ 2/3 \\ 1}z
       \suchthat z\in\Re}
\end{equation*}
与前面的例题中一样，与参数无关的向量
\begin{equation*}
   \colvec[r]{1 \\ 0 \\ 0}
\end{equation*}
是该方程组的一个特解。
\end{example}

在练习之前，我们先考虑一下已经完成了什么，
以及本章余下部分要做什么。
到目前为止，我们做的是高斯消元法的机械步骤。
我们还没有停下来考虑过出现的任何问题，
只是证明了\nearbytheorem{th:GaussMethod}\Dash该定理
通过表明这个方法给出
正确的答案而为方法提供了依据。

例如，我们总能像上面那样描述解集吗，
即用一个特解向量
加上若干其他向量的不受限制的线性组合？
我们已经注意到，这样描述的解集
有无穷多个成员，
因此回答这个问题
会告诉我们解集的大小。
下一小节将表明答案是“肯定的”。
本章第二节随后
将利用这个答案来描述解集的几何。

另一些问题来自这样的观察：我们做高斯消元法可以
不止一种方式（例如，对换行时，我们可能有多行
可供选择对换）。
\nearbytheorem{th:GaussMethod}说，无论我们怎么做，都必定得到同一个解集，但是
如果用两种方式做高斯消元法，
在每个阶梯形方程组中我们必定得到相同个数的自由变量吗？
那些必定是同样的变量吗？也就是说，用一种方式求解一个问题得到$y$和$w$为自由变量，而用另一种方式求解却得到$y$和$z$为自由变量，这是不可能的吗？
本章第三节回答“是的”：
从任何初始线性方程组出发，
所有导出的阶梯形版本
都有相同的自由变量。

因此，到本章结束时，我们不仅会在
高斯消元法的实践上有扎实的基础，
也会在理论上
有扎实的基础。
我们将确切地知道在一次消元中什么会发生、什么不会发生。

\begin{exercises}
  \recommended \item  
    求矩阵的指定元素，
    （若它有定义）。
    \begin{equation*}
      A=\begin{mat}[r]
        1  &3  &1  \\
        2  &-1 &4
      \end{mat}
    \end{equation*}
    \begin{exparts*}
      \partsitem \( a_{2,1} \)
      \partsitem \( a_{1,2} \)
      \partsitem \( a_{2,2} \)
      \partsitem \( a_{3,1} \)
    \end{exparts*}
    \begin{answer}
      \begin{exparts*}
        \partsitem \( 2 \)
        \partsitem \( 3 \)
        \partsitem \(-1 \)
        \partsitem 无定义。
      \end{exparts*}  
    \end{answer}
  \recommended \item 
    给出每个矩阵的尺寸。
    \begin{exparts*}
      \partsitem \(
        \begin{mat}[r]
          1  &0  &4  \\
          2  &1  &5
        \end{mat}  \)
      \partsitem \(
        \begin{mat}[r]
          1  &1  \\
         -1  &1  \\
          3  &-1
        \end{mat}  \)
      \partsitem \(
        \begin{mat}[r]
          5  &10 \\
         10  &5
        \end{mat}  \)
    \end{exparts*}
    \begin{answer}
      \begin{exparts*}
        \partsitem \( \nbym{2}{3} \)
        \partsitem \( \nbym{3}{2} \)
        \partsitem \( \nbym{2}{2} \)
      \end{exparts*}  
    \end{answer}
  \recommended \item 
    进行指定的向量运算（若它有定义）。
    \begin{exparts*}
      \partsitem \( \colvec[r]{2 \\ 1 \\ 1}
               +\colvec[r]{3 \\ 0 \\ 4} \)
      \partsitem \( 5\colvec[r]{4 \\ -1} \)
      \partsitem \( \colvec[r]{1 \\ 5 \\ 1}
               -\colvec[r]{3 \\ 1 \\ 1} \)
      \partsitem \( 7\colvec[r]{2 \\ 1}
               +9\colvec[r]{3 \\ 5} \)
      \partsitem \( \colvec[r]{1 \\ 2}
               +\colvec[r]{1 \\ 2 \\ 3} \)
      \partsitem \( 6\colvec[r]{3 \\ 1 \\ 1}
               -4\colvec[r]{2 \\ 0 \\ 3}
               +2\colvec[r]{1 \\ 1 \\ 5} \)
    \end{exparts*}
    \begin{answer}
      \begin{exparts*}
        \partsitem \( \colvec[r]{5 \\ 1 \\ 5} \)
        \partsitem \( \colvec[r]{20 \\ -5} \)
        \partsitem \( \colvec[r]{-2 \\ 4 \\ 0} \)
        \partsitem \( \colvec[r]{41 \\ 52} \)
        \partsitem 无定义。
        \partsitem \( \colvec[r]{12 \\ 8 \\ 4} \)
      \end{exparts*}  
     \end{answer}
  \recommended \item  \label{exer:SolveInMatrixNotation}
    用矩阵记法求解每个方程组。
    用向量表示解。
    \begin{exparts*}
      \partsitem \( \begin{linsys}[t]{2}
                  3x  &+  &6y  &=  &18  \\
                   x  &+  &2y  &=  &6   
                   \end{linsys}  \)
      \partsitem \( \begin{linsys}[t]{2}
                   x  &+  &y   &=  &1  \\
                   x  &-  &y   &=  &-1   
                    \end{linsys}  \)
      \partsitem \( \begin{linsys}[t]{3}
                   x_1  &   &     &+  &x_3   &=  &4  \\
                   x_1  &-  &x_2  &+  &2x_3  &=  &5  \\
                  4x_1  &-  &x_2  &+  &5x_3  &=  &17  
                   \end{linsys}  \)
      \partsitem \( \begin{linsys}[t]{3}
                   2a   &+  &b    &-  &c     &=  &2  \\
                   2a   &   &     &+  &c     &=  &3  \\
                    a   &-  &b    &   &      &=  &0   
                    \end{linsys}  \)
      \partsitem \( \begin{linsys}[t]{4}
                     x  &+  &2y   &-   &z   &    &    &=  &3  \\
                    2x  &+  &y    &    &    &+   &w   &=  &4  \\
                     x  &-  &y    &+   &z   &+   &w   &=  &1  
                    \end{linsys}  \)
      \partsitem \( \begin{linsys}[t]{4}
                     x  &   &     &+   &z   &+   &w   &=  &4  \\
                    2x  &+  &y    &    &    &-   &w   &=  &2  \\
                    3x  &+  &y    &+   &z   &    &    &=  &7  
                     \end{linsys}  \)
    \end{exparts*}
    \begin{answer}
      \begin{exparts}
        \partsitem 这一消元
          \begin{equation*}
            \begin{amat}[r]{2}
              3  &6  &18 \\
              1  &2  &6
            \end{amat}
            \grstep{(-1/3)\rho_1+\rho_2}
            \begin{amat}[r]{2}
              3  &6  &18 \\
              0  &0  &0
            \end{amat}
          \end{equation*}
          留下\( x \)为首变量而\( y \)为自由变量。
          取\( y \)为参数，得到\( x=6-2y \)以及这个解
          集。
          \begin{equation*}
            \set{\colvec[r]{6 \\ 0}+\colvec[r]{-2 \\ 1}y
              \suchthat y\in\Re}
          \end{equation*}
        \partsitem 一次消元
          \begin{equation*}
            \begin{amat}[r]{2}
              1  &1  &1  \\
              1  &-1 &-1
            \end{amat}
            \grstep{-\rho_1+\rho_2}
            \begin{amat}[r]{2}
              1  &1  &1  \\
              0  &-2 &-2
            \end{amat}
          \end{equation*}
          给出唯一解\( y=1 \)、\( x=0 \)。
          解集如下。
          \begin{equation*}
            \set{\colvec[r]{0 \\ 1} }
          \end{equation*}
        \partsitem 高斯消元法
          \begin{equation*}
            \begin{amat}[r]{3}
              1  &0  &1  &4  \\
              1  &-1 &2  &5  \\
              4  &-1 &5  &17
            \end{amat}
            \grstep[-4\rho_1+\rho_3]{-\rho_1+\rho_2}
            \begin{amat}[r]{3}
              1  &0  &1  &4  \\
              0  &-1 &1  &1  \\
              0  &-1 &1  &1
            \end{amat}     
            \grstep{-\rho_2+\rho_3}
            \begin{amat}[r]{3}
              1  &0  &1  &4  \\
              0  &-1 &1  &1  \\
              0  &0  &0  &0
            \end{amat}
          \end{equation*}
          留下\( x_1 \)和\( x_2 \)为首变量，而\( x_3 \)为自由变量。
          解集如下。
          \begin{equation*}
            \set{\colvec[r]{4 \\ -1 \\ 0}+\colvec[r]{-1 \\ 1 \\ 1}x_3
              \suchthat x_3\in\Re}
          \end{equation*}
        \partsitem 这一消元
          \begin{align*}
            \begin{amat}[r]{3}
              2  &1  &-1 &2  \\
              2  &0  &1  &3  \\
              1  &-1 &0  &0
            \end{amat}
            &\grstep[-(1/2)\rho_1+\rho_3]{-\rho_1+\rho_2}
            \begin{amat}[r]{3}
              2  &1    &-1   &2  \\
              0  &-1   &2    &1  \\
              0  &-3/2 &1/2  &-1
            \end{amat}                                          \\
            &\grstep{(-3/2)\rho_2+\rho_3}
            \begin{amat}[r]{3}
              2  &1  &-1   &2  \\
              0  &-1 &2    &1  \\
              0  &0  &-5/2 &-5/2
            \end{amat}
          \end{align*}
          表明解集是单元素集。
          \begin{equation*}
            \set{\colvec[r]{1 \\ 1 \\ 1}}
          \end{equation*}
        \partsitem 这一消元很容易
          \begin{align*}
            \begin{amat}[r]{4}
              1  &2  &-1 &0  &3 \\
              2  &1  &0  &1  &4 \\
              1  &-1 &1  &1  &1
            \end{amat}
            &\grstep[-\rho_1+\rho_3]{-2\rho_1+\rho_2}
            \begin{amat}[r]{4}
              1  &2  &-1 &0  &3  \\
              0  &-3 &2  &1  &-2 \\
              0  &-3 &2  &1  &-2
            \end{amat}                                       \\
            &\grstep{-\rho_2+\rho_3}
            \begin{amat}[r]{4}
              1  &2  &-1 &0  &3  \\
              0  &-3 &2  &1  &-2 \\
              0  &0  &0  &0  &0
            \end{amat}
          \end{align*}
          结束时\( x \)和$y$为首变量，而\( z \)和\( w \)为
          自由变量。
          解出\( y \)得到\( y=(2+2z+w)/3 \)，代入表明
          \( x+2(2+2z+w)/3-z=3 \)，于是\( x=(5/3)-(1/3)z-(2/3)w \)，
          由此得到这个解集。
          \begin{equation*}
            \set{\colvec[r]{5/3 \\ 2/3 \\ 0 \\ 0}
                 +\colvec[r]{-1/3 \\ 2/3 \\ 1 \\ 0}z
                 +\colvec[r]{-2/3 \\ 1/3 \\ 0 \\ 1}w
                 \suchthat z,w\in\Re}
          \end{equation*}
        \partsitem 消元
          \begin{align*}
            \begin{amat}[r]{4}
              1  &0  &1  &1  &4 \\
              2  &1  &0  &-1 &2 \\
              3  &1  &1  &0  &7
            \end{amat}
            &\grstep[-3\rho_1+\rho_3]{-2\rho_1+\rho_2}
            \begin{amat}[r]{4}
              1  &0  &1  &1  &4 \\
              0  &1  &-2 &-3 &-6\\
              0  &1  &-2 &-3 &-5
            \end{amat}                                       \\
            &\grstep{-\rho_2+\rho_3}
            \begin{amat}[r]{4}
              1  &0  &1  &1  &4 \\
              0  &1  &-2 &-3 &-6\\
              0  &0  &0  &0  &1
            \end{amat}
          \end{align*}
          表明无解\Dash解集为空。
      \end{exparts}  
     \end{answer}
  \item \label{exer:SlvMatNot}
    用矩阵记法求解每个方程组。
    用向量记法给出每个解集。
    \begin{exparts*}
      \partsitem \( \begin{linsys}[t]{3}
                  2x  &+  &y  &-  &z  &=  &1  \\
                  4x  &-  &y  &   &   &=  &3  
                \end{linsys}  \)
      \partsitem \( \begin{linsys}[t]{4}
                   x  &   &   &-  &z  &   &   &=  &1  \\
                      &   &y  &+  &2z &-  &w  &=  &3  \\
                   x  &+  &2y &+  &3z &-  &w  &=  &7  
               \end{linsys}  \)
      \partsitem \( \begin{linsys}[t]{4}
                   x  &-  &y  &+  &z  &   &   &=  &0  \\
                      &   &y  &   &   &+  &w  &=  &0  \\
                  3x  &-  &2y &+  &3z &+  &w  &=  &0  \\
                      &   &-y &   &   &-  &w  &=  &0  
               \end{linsys}  \)
      \partsitem \( \begin{linsys}[t]{5}
                   a  &+  &2b &+  &3c &+  &d  &-  &e  &=  &1  \\
                  3a  &-  &b  &+  &c  &+  &d  &+  &e  &=  &3  
               \end{linsys}  \)
    \end{exparts*}
    \begin{answer}
      \begin{exparts}
      \partsitem 这一消元
        \begin{equation*}
          \begin{amat}[r]{3}
            2  &1  &-1  &1  \\
            4  &-1 &0   &3
          \end{amat}
          \grstep{-2\rho_1+\rho_2}
          \begin{amat}[r]{3}
            2  &1  &-1  &1  \\
            0  &-3 &2   &1
          \end{amat}
        \end{equation*}
        结束时\( x \)和\( y \)为首变量，而\( z \)为自由变量。
        解出\( y \)得到\( y=(1-2z)/(-3) \)，然后代入
        \( 2x+(1-2z)/(-3)-z=1 \)表明\( x=((4/3)+(1/3)z)/2 \)。
        因此解集如下。
        \begin{equation*}
          \set{\colvec[r]{2/3 \\ -1/3 \\ 0}
               +\colvec[r]{1/6 \\ 2/3 \\ 1}z
              \suchthat z\in\Re}
        \end{equation*}
      \partsitem 这样应用高斯消元法
        \begin{align*}
          \begin{amat}[r]{4}
            1  &0  &-1  &0  &1 \\
            0  &1  &2   &-1 &3 \\
            1  &2  &3   &-1 &7
          \end{amat}
          &\grstep{-\rho_1+\rho_3}
          \begin{amat}[r]{4}
            1  &0  &-1  &0  &1 \\
            0  &1  &2   &-1 &3 \\
            0  &2  &4   &-1 &6
          \end{amat}                                   \\
          &\grstep{-2\rho_2+\rho_3}
          \begin{amat}[r]{4}
            1  &0  &-1  &0  &1 \\
            0  &1  &2   &-1 &3 \\
            0  &0  &0   &1  &0
          \end{amat}
        \end{align*}
        留下\( x \)、\( y \)和\( w \)为首变量。
        解集如下。
        \begin{equation*}
          \set{\colvec[r]{1 \\ 3 \\ 0 \\ 0}
               +\colvec[r]{1 \\ -2 \\ 1 \\ 0}z
              \suchthat z\in\Re}
        \end{equation*}
      \partsitem 这次行化简
        \begin{align*}
          \begin{amat}[r]{4}
            1  &-1 &1   &0  &0 \\
            0  &1  &0   &1  &0 \\
            3  &-2 &3   &1  &0 \\
            0  &-1 &0   &-1 &0
          \end{amat}
          &\grstep{-3\rho_1+\rho_3}
          \begin{amat}[r]{4}
            1  &-1 &1   &0  &0 \\
            0  &1  &0   &1  &0 \\
            0  &1  &0   &1  &0 \\
            0  &-1 &0   &-1 &0
          \end{amat}                                       \\
          &\grstep[\rho_2+\rho_4]{-\rho_2+\rho_3}
          \begin{amat}[r]{4}
            1  &-1 &1   &0  &0 \\
            0  &1  &0   &1  &0 \\
            0  &0  &0   &0  &0 \\
            0  &0  &0   &0  &0
          \end{amat}
        \end{align*}
        结束时\( z \)和\( w \)为自由变量。
        我们得到这个解集。
        \begin{equation*}
          \set{\colvec[r]{0 \\ 0 \\ 0 \\ 0}
               +\colvec[r]{-1 \\ 0 \\ 1 \\ 0}z
               +\colvec[r]{-1 \\ -1 \\ 0 \\ 1}w
              \suchthat z,w\in\Re}
        \end{equation*}
      \partsitem 按这种方式做高斯消元法
        \begin{equation*}
          \begin{amat}[r]{5}
            1  &2  &3   &1  &-1 &1  \\
            3  &-1 &1   &1  &1  &3
          \end{amat}
          \grstep{-3\rho_1+\rho_2}
          \begin{amat}[r]{5}
            1  &2  &3   &1  &-1 &1  \\
            0  &-7 &-8  &-2 &4  &0
          \end{amat}
        \end{equation*}
        结束时\( c \)、\( d \)和\( e \)为自由变量。
        解出\( b \)表明\( b=(8c+2d-4e)/(-7) \)，然后 
        代入
        \( a+2(8c+2d-4e)/(-7)+3c+1d-1e=1 \)表明
        \( a=1-(5/7)c-(3/7)d-(1/7)e \)，于是我们得到解集。
        \begin{equation*}
          \set{\colvec[r]{1 \\ 0 \\ 0 \\ 0 \\ 0}
               +\colvec[r]{-5/7 \\ -8/7 \\ 1 \\ 0 \\ 0}c
               +\colvec[r]{-3/7 \\ -2/7 \\ 0 \\ 1 \\ 0}d
               +\colvec[r]{-1/7 \\ 4/7 \\ 0 \\ 0 \\ 1}e
              \suchthat c,d,e\in\Re}
        \end{equation*}
    \end{exparts}  
   \end{answer}
  \item 
    用矩阵记法求解每个方程组。
    用向量表示解集。
    \begin{exparts*}
      \partsitem 
        $\begin{linsys}{3}
          3x &+ &2y &+ &z &= &1 \\
          x  &- &y  &+ &z &= &2 \\
          5x &+ &5y &+ &z &= &0
        \end{linsys}$
      \partsitem
        $\begin{linsys}{4}
          x  &+ &y  &- &2z &= &0 \\
          x  &- &y  &  &   &= &-3 \\
          3x &- &y  &- &2z &= &-6  \\
             &  &2y &- &2z &= &3  
        \end{linsys}$
      \partsitem
        $\begin{linsys}{5}
          2x  &- &y  &- &z &+ &w &= &4 \\
           x  &+ &y  &+ &z &  &  &= &-1 
        \end{linsys}$
      \partsitem
         $\begin{linsys}{3}
          x  &+ &y  &- &2z &= &0 \\
          x  &- &y  &  &   &= &-3 \\
          3x &- &y  &- &2z &= &0    
        \end{linsys}$ 
    \end{exparts*}
    \begin{answer}
      \begin{exparts}
        \partsitem 这一消元
          \begin{align*}
            \begin{amat}{3}
              3 &2  &1  &1 \\
              1 &-1 &1  &2 \\
              5 &5  &1  &0
            \end{amat}
            &\grstep[-(5/3)\rho_1+\rho_3]{-(1/3)\rho_1+\rho_2}
            \begin{amat}{3}
              3 &2    &1     &1 \\
              0 &-5/3 &2/3   &5/3 \\
              0 &5/3  &-2/3  &-5/3
            \end{amat}                                     \\
            &\grstep{\rho_2+\rho_3}
            \begin{amat}{3}
              3 &2    &1     &1 \\
              0 &-5/3 &2/3   &5/3 \\
              0 &0    &0     &0
            \end{amat}
          \end{align*}
          给出这个解集。
          \begin{equation*}
            \set{\colvec{x \\ y \\ z}=\colvec{1 \\ -1 \\ 0}
                                       +\colvec{-3/5 \\ 2/5 \\ 1}z
                             \suchthat z\in\Re}
          \end{equation*}
        \partsitem
          消元如下。
          \begin{align*}
            \begin{amat}{3}
              1 &1   &-2 &0 \\
              1 &-1  &0  &3  \\
              3 &-1  &-2 &-6 \\
              0 &2   &-2 &3
            \end{amat}
            &\grstep[-3\rho_1+\rho_3]{-\rho_1+\rho_2}
            \begin{amat}{3}
              1 &1   &-2 &0 \\
              0 &-2  &2  &-3  \\
              0 &-4  &4  &-6 \\
              0 &2   &-2 &3
            \end{amat}                                  \\
            &\grstep[\rho_2+\rho_4]{-2\rho_2+\rho_3}
            \begin{amat}{3}
              1 &1   &-2 &0 \\
              0 &-2  &2  &-3  \\
              0 &0   &0  &0 \\
              0 &0   &0  &0
            \end{amat}
          \end{align*}
          解集如下。
          \begin{equation*}
            \set{\colvec{-3/2 \\ 3/2 \\ 0}
                 +\colvec{1 \\ 1 \\ 1}z
                 \suchthat z\in\Re}
          \end{equation*}
      \partsitem
         高斯消元法
         \begin{equation*}
           \begin{amat}{4}
             2 &-1 &-1 &1 &4 \\
             1 &1  &1  &0 &-1
           \end{amat}
           \grstep{-(1/2)\rho_1+\rho_2}
           \begin{amat}{4}
             2 &-1   &-1   &1    &4 \\
             0 &3/2  &3/2  &-1/2 &-3
           \end{amat}
         \end{equation*}
         给出解集。
         \begin{equation*}
           \set{\colvec{1 \\ -2 \\ 0 \\ 0}
                  +\colvec{0 \\ -1 \\ 1 \\ 0}z
                   \colvec{-1/3 \\ 1/3 \\ 0 \\ 1}w
                 \suchthat z,w\in\Re}
         \end{equation*}
       \partsitem
          消元如下。
          \begin{equation*}
            \begin{amat}{3}
              1 &1  &-2 &0  \\
              1 &-1 &0  &-3 \\
              3 &-1 &-2 &0
           \end{amat}
           \grstep[-3\rho_1+\rho_3]{-\rho_1+\rho_2}
           \begin{amat}{3}
             1 &1  &-2 &0  \\
             0 &-2 &2  &-3 \\
             0 &-4  &4  &0
          \end{amat}
          \grstep{-2\rho_2+\rho_3}
          \begin{amat}{3}
            1 &1  &-2 &0  \\
            0 &-2 &2  &-3 \\
            0 &0  &0  &6
          \end{amat}
        \end{equation*}
        解集为空集~$\set{}$。
      \end{exparts}
    \end{answer}
  \recommended \item 
    该向量在此集合中。
    参数取什么值会产生这个向量？
    \begin{exparts}
      \partsitem $\colvec[r]{5 \\ -5}$、
        $\set{\colvec[r]{1 \\ -1}k\suchthat k\in\Re}$
      \partsitem $\colvec[r]{-1 \\ 2 \\ 1}$、
        $\set{\colvec[r]{-2 \\ 1 \\ 0}i
           +\colvec[r]{3 \\ 0 \\ 1}j\suchthat i,j\in\Re}$
      \partsitem $\colvec[r]{0 \\ -4 \\ 2}$、
        $\set{\colvec[r]{1 \\ 1 \\ 0}m
               +\colvec[r]{2 \\ 0 \\ 1}n\suchthat m,n\in\Re}$
    \end{exparts}
    \begin{answer}
      对每一小题，我们通过考察分量的等式得到一个
      线性方程组。
      \begin{exparts}
       \partsitem $k=5$
       \partsitem 第二个分量表明$i=2$，第三个
       分量表明$j=1$。
       \partsitem $m=-4$，$n=2$
      \end{exparts} 
    \end{answer}
  \item 
    判断该向量是否在集合中。
    \begin{exparts}
      \partsitem $\colvec[r]{3 \\ -1}$、
        $\set{\colvec[r]{-6 \\ 2}k\suchthat k\in\Re}$
      \partsitem $\colvec[r]{5 \\ 4}$、
        $\set{\colvec[r]{5 \\ -4}j\suchthat j\in\Re}$
      \partsitem $\colvec[r]{2 \\ 1 \\ -1}$、
        $\set{\colvec[r]{0 \\ 3 \\ -7}
             +\colvec[r]{1 \\ -1 \\ 3}r\suchthat r\in\Re}$
      \partsitem $\colvec[r]{1 \\ 0 \\ 1}$、
        $\set{\colvec[r]{2 \\ 0 \\ 1}j
            +\colvec[r]{-3 \\ -1 \\ 1}k\suchthat j,k\in\Re}$
    \end{exparts}
    \begin{answer}
      对每一小题，我们通过考察分量的等式得到一个
      线性方程组。
      \begin{exparts}
        \partsitem 是；取$k=-1/2$。
        \partsitem 否；由方程$5=5\cdot j$和
            $4=-4\cdot j$组成的方程组无解。
        \partsitem 是；取$r=2$。
        \partsitem 否。
           第二个分量给出$k=0$。
           于是第三个分量给出$j=1$。
           但第一个分量对不上。
      \end{exparts}
     \end{answer}
  \item \cite{Cleary}
    一个有1200英亩土地的农民在考虑种植三种不同的庄稼：
    玉米、大豆和燕麦。   
    农民想用全部~$1200$英亩土地。  
    玉米种子每英亩花费\$$20$，而大豆种子和燕麦种子分别
    每英亩花费\$$50$和\$$12$。  
    农民有\$$40\,000$可用于购买种子，并打算
    把它全部花掉。
    \begin{exparts}  
      \item 利用上面的信息建立两个含三个未知数的线性方程
        并求解它。
     \item 该方程组的解就是农民可以做的选择。  
        写出两个合理的解。
     \item 假设秋天庄稼成熟时，农民
        每英亩玉米可带来\$$100$的收入，
        每英亩大豆\$$300$、每英亩燕麦\$$80$。  
        前一题的两个解中哪一个原本会带来
        更大的收入？ 
    \end{exparts}
    \begin{answer}
      \begin{exparts}
        \item 设$c$为玉米的英亩数，$s$为
          大豆的英亩数，$a$为燕麦的英亩数。
          \begin{equation*}
            \begin{linsys}{3}
              c   &+   &s   &+   &a   &=   &1200 \\ 
            20c   &+   &50s &+   &12a &=   &40\,000  
            \end{linsys}
            \grstep{-20\rho_1+\rho_2}
            \begin{linsys}{3}
              c   &+   &s   &+   &a   &=   &1200 \\ 
                  &    &30s &-   &8a  &=   &16\,000  
            \end{linsys}
          \end{equation*}
          为描述解集，我们可以用$a$参数化。
          \begin{equation*}
            \set{\colvec{c \\ s \\ a}
                 =\colvec{20\,000/30 \\ 16\,000/30 \\ 0}
                  +\colvec{-38/30 \\ 8/30 \\ 1}a
                 \suchthat a\in\Re}
          \end{equation*}
        \item 这里有许多可能的答案。 
          例如我们可以取$a=0$得到$c=20\,000/30\approx 666.66$和
          $s=16000/30\approx 533.33$。
          另一个例子是取$a=20\,000/38\approx 526.32$，给出
          $c=0$和$s=7360/38\approx 193.68$。
        \item 把前一题的答案代入
          $100c+300s+80a$。
      \end{exparts}
    \end{answer}
  \item 
    参数化这个单方程方程组的解集。
    \begin{equation*}
      x_1+x_2+\cdots+x_n=0
    \end{equation*}
    \begin{answer}
      这个方程组有一个方程。
      首变量是\( x_1 \)，其余变量是自由变量。
      \begin{equation*}
        \set{\colvec{-1 \\ 1 \\ \vdotswithin{-1} \\ 0}x_2
             +\cdots+
             \colvec{-1 \\ 0 \\ \vdotswithin{-1} \\ 1}x_n
             \suchthat x_2,\ldots,x_n\in\Re}
      \end{equation*}  
     \end{answer}
  \recommended \item 
    \begin{exparts}
    \partsitem 对左边应用高斯消元法来求解
      \begin{equation*}
        \begin{linsys}{4}
          x  &+  &2y  &    &    &-   &w   &=   &a   \\
         2x  &   &    &+   &z   &    &    &=   &b   \\
          x  &+  &y   &    &    &+   &2w  &=   &c   
        \end{linsys}
      \end{equation*}
      把\( x \)、\( y \)、\( z \)和\( w \)用
      常数$a$、$b$和$c$表示。
    \partsitem 用前一题的答案求解这个方程组。
      \begin{equation*}
        \begin{linsys}{4}
          x  &+  &2y  &    &    &-   &w   &=   &3   \\
         2x  &   &    &+   &z   &    &    &=   &1   \\
          x  &+  &y   &    &    &+   &2w  &=   &-2
        \end{linsys}
      \end{equation*}
    \end{exparts}
    \begin{answer}
      \begin{exparts}
        \partsitem 这里高斯消元法给出
          \begin{align*}
            \begin{amat}{4}
              1  &2  &0  &-1  &a  \\
              2  &0  &1  &0   &b  \\
              1  &1  &0  &2   &c
            \end{amat}
            &\grstep[-\rho_1+\rho_3]{-2\rho_1+\rho_2}
            \begin{amat}{4}
              1  &2  &0  &-1  &a  \\
              0  &-4 &1  &2   &-2a+b  \\
              0  &-1 &0  &3   &-a+c
            \end{amat}                                  \\
            &\grstep{-(1/4)\rho_2+\rho_3}
            \begin{amat}{4}
              1  &2  &0    &-1  &a  \\
              0  &-4 &1    &2   &-2a+b  \\
              0  &0  &-1/4 &5/2 &-(1/2)a-(1/4)b+c
            \end{amat}
          \end{align*}
          留下\( w \)为自由变量。
          求解：\( z=2a+b-4c+10w \)，
          且\( -4y=-2a+b-(2a+b-4c+10w)-2w \)于是
          \( y=a-c+3w \)，并且
          \( x=a-2(a-c+3w)+w=-a+2c-5w. \)
          因此解集如下。
          \begin{equation*}
             \set{\colvec{-a+2c \\ a-c \\ 2a+b-4c \\ 0}
                  +\colvec{-5 \\ 3 \\ 10 \\ 1}w
                  \suchthat w\in\Re}
          \end{equation*}
        \partsitem 代入\( a=3 \)、\( b=1 \)和\( c=-2 \)。
          \begin{equation*}
             \set{\colvec[r]{-7 \\ 5 \\ 15 \\ 0}
                  +\colvec[r]{-5 \\ 3 \\ 10 \\ 1}w
                  \suchthat w\in\Re}
          \end{equation*}
      \end{exparts}  
     \end{answer}
  \item 
    在矩阵元素的记法“\( a_{i,j} \)”中，
    为什么需要逗号？
    \begin{answer}
       省略逗号，比如写成\( a_{123} \)，
       是有歧义的，因为它可能指$a_{1,23}$或$a_{12,3}$。  
    \end{answer}
  \recommended \item 
    给出\( \nbyn{4} \)矩阵，其
    \( i,j \)元素为
    \begin{exparts*}
      \partsitem \( i+j \)；
      \partsitem \( -1 \)的\( i+j \)次幂。
    \end{exparts*}
    \begin{answer}
      \begin{exparts*}
        \partsitem \(
           \begin{mat}[r]
             2  &3  &4  &5  \\
             3  &4  &5  &6  \\
             4  &5  &6  &7  \\
             5  &6  &7  &8
           \end{mat} \)
        \partsitem \(
           \begin{mat}[r]
             1  &-1  &1   &-1  \\
            -1  &1   &-1  &1  \\
             1  &-1  &1   &-1  \\
            -1  &1   &-1  &1
           \end{mat} \)
      \end{exparts*}  
    \end{answer}
  \item  
    对任意矩阵\( A \)，
    \( A \)的\definend{转置}\index{transpose}
    \index{matrix!transpose}
    记作
    \( \trans{A} \)，是以\( A \)的行作为列的矩阵。
    求下面每一个的转置。
    \begin{exparts*}
      \partsitem \( \begin{mat}[r]
                  1  &2  &3  \\
                  4  &5  &6
               \end{mat}  \)
      \partsitem \( \begin{mat}[r]
                  2  &-3 \\
                  1  &1
               \end{mat}  \)
      \partsitem \( \begin{mat}[r]
                  5  &10 \\
                 10  &5
               \end{mat}  \)
      \partsitem \( \colvec[r]{1 \\ 1 \\ 0} \)
    \end{exparts*}
    \begin{answer}
      \begin{exparts*}
        \partsitem \( \begin{mat}[r]
                   1  &4  \\
                   2  &5  \\
                   3  &6
                 \end{mat}  \)
        \partsitem \( \begin{mat}[r]
                   2  &1  \\
                  -3  &1
                 \end{mat}  \)
        \partsitem \( \begin{mat}[r]
                   5  &10 \\
                  10  &5
                 \end{mat}  \)
        \partsitem \( \rowvec{1 &1 &0}  \)
      \end{exparts*}  
     \end{answer}
  \recommended \item 
    \begin{exparts}
      \partsitem 描述所有满足\( f(1)=2 \)且\( f(-1)=6 \)的
        函数\( f(x)=ax^2+bx+c \)。
      \partsitem 描述所有满足\( f(1)=2 \)的
        函数\( f(x)=ax^2+bx+c \)。
    \end{exparts}
    \begin{answer}
      \begin{exparts}
        \partsitem 代入\( x=1 \)和\( x=-1 \)得到
          \begin{equation*}
            \begin{linsys}{3}
              a  &+  &b   &+  &c  &=  &2  \\
              a  &-  &b   &+  &c  &=  &6  
            \end{linsys}
            \grstep{-\rho_1+\rho_2}
            \begin{linsys}{3}
              a  &+  &b   &+  &c  &=  &2  \\
                 &   &-2b &   &   &=  &4  
              \end{linsys}
          \end{equation*}
          所以函数的集合是
          \( \set{f(x)=(4-c)x^2-2x+c\suchthat c\in\Re} \)。
        \partsitem 代入\( x=1 \)得到
          \begin{equation*}
            \begin{linsys}{3}
              a  &+  &b   &+  &c  &=  &2  
            \end{linsys}
          \end{equation*}
          所以函数的集合是
          \( \set{f(x)=(2-b-c)x^2+bx+c\suchthat b,c\in\Re} \)。
      \end{exparts}  
    \end{answer}
  \item 证明：平面\( \Re^2 \)中任意五个点都位于一条共同的二次曲线上，也就是说，它们都满足某个形如
    \( ax^2+by^2+cxy+dx+ey+f=0 \)的方程，其中\( a,\,\ldots\,,f \)
    中有些不为零。
    \begin{answer}
      把五对$(x,y)$代入，我们得到一个有
      五个方程、六个未知数$a$、\ldots、$f$的方程组。
      由于未知数比方程多，如果方程之间没有不相容，
      那么就有无穷多解
      （至少有一个变量最终会成为自由变量）。

      但不相容不可能存在，因为$a=0$、\ldots、$f=0$是
      一个解（我们只是用这个零解来表明方程组
      是相容的\Dash前面一段表明
      存在非零解）。 
    \end{answer}
  \item 
    编出一个四方程/四未知数的方程组，使其具有
    \begin{exparts}
      \partsitem 单参数的解集；
      \partsitem 双参数的解集；
      \partsitem 三参数的解集。
    \end{exparts}
    \begin{answer}
      \begin{exparts}
      \partsitem 这是一个\Dash第四个方程是多余的
        但仍然可以。
        \begin{equation*}
          \begin{linsys}{4}
             x  &+  &y  &-  &z  &+  &w  &=  &0  \\
                &   &y  &-  &z  &   &   &=  &0  \\
                &   &   &   &2z &+  &2w &=  &0  \\
                &   &   &   &z  &+  &w  &=  &0
          \end{linsys}
        \end{equation*}
      \partsitem 这是一个。
        \begin{equation*}
          \begin{linsys}{4}
             x  &+  &y  &-  &z  &+  &w  &=  &0  \\
                &   &   &   &   &   &w  &=  &0  \\
                &   &   &   &   &   &w  &=  &0  \\
                &   &   &   &   &   &w  &=  &0
          \end{linsys}
        \end{equation*}
      \partsitem 这是一个。
        \begin{equation*}
          \begin{linsys}{4}
             x  &+  &y  &-  &z  &+  &w  &=  &0  \\
             x  &+  &y  &-  &z  &+  &w  &=  &0  \\
             x  &+  &y  &-  &z  &+  &w  &=  &0  \\
             x  &+  &y  &-  &z  &+  &w  &=  &0 
          \end{linsys}
        \end{equation*}
    \end{exparts}  
   \end{answer}
  \puzzle \item 
    \cite{Shepelev}
    这个谜题来自俄罗斯网站
    \texttt{http://www.arbuz.uz/}，它有很多种解法，
    但我的解法用到线性代数，而且非常朴素。有一颗行星上居住着arbuzoid
    （从俄语翻译过来即“西瓜人”）。
    这些生物有三种颜色：红、绿、蓝。
    有$13$~个红色的
    arbuzoid、$15$~个蓝色的、$17$~个绿色的。
    当
    两只颜色不同的arbuzoid相遇时，它们
    都变成第三种颜色。

    问题是，有没有可能出现它们
    全部变成同一种颜色的情况？
     \begin{answer}
        \answerasgiven
        我的解法是把arbuzoid的
        数量定义为$3$维向量，并把所有可能的
        基本转变也表示成这样的向量：
        \begin{center}
          \begin{tabular}{rr}
            R: &$13$  \\
            G: &$15$  \\
            B: &$17$
          \end{tabular}
          \qquad
          操作：
          $\colvec[r]{-1 \\ -1 \\ 2}$、
          $\colvec[r]{-1 \\ 2 \\ -1}$、
          以及
          $\colvec[r]{2 \\ -1 \\ -1}$
        \end{center}
        现在，只需检查下面某个线性方程组的解
        是否
        存在：
        \begin{equation*}
          \colvec[r]{13 \\ 15 \\ 17}
          +x\colvec[r]{-1 \\ -1  \\ 2}
          +y\colvec[r]{-1 \\ 2 \\ -1}
          +z\colvec[r]{2 \\ -1 \\ -1}
          =\colvec[r]{0 \\ 0 \\ 45}
          \qquad
          \text{（或 $\colvec[r]{0 \\ 45 \\ 0}$ 或 $\colvec[r]{45 \\ 0 \\ 0}$）}
        \end{equation*}
        求解
        \begin{equation*}
          \begin{amat}[r]{3}
           -1  &-1 &2  &-13  \\
           -1  &2  &-1 &-15  \\
            2  &-1 &-1 &28
          \end{amat}
          \grstep[2\rho_1+\rho_3]{-\rho_1+\rho_2}
          \repeatedgrstep{\rho_2+\rho_3}
          \begin{amat}[r]{3}
           -1  &-1 &2  &-13  \\
            0  &3  &-3 &-2  \\
            0  &0  &0  &0
          \end{amat}
        \end{equation*}
        给出$y+2/3=z$，所以如果变换次数$z$是一个整数
        那么$y$就不是。
        另外两个方程组给出类似的结论，因此无解。
     \end{answer}
  \puzzle \item 
    \cite{USSROlympiad174}
    \begin{exparts}
      \partsitem 求解方程组。
        \begin{equation*}
          \begin{linsys}{2}
            ax  &+  &y  &=  &a^2  \\
             x  &+  &ay &=  &1
         \end{linsys}
        \end{equation*}
        对哪些$a$值，方程组无解，而对哪些$a$值，方程组有无穷多解？
      \partsitem 对方程组回答上面的问题。
        \begin{equation*}
          \begin{linsys}{2}
            ax  &+  &y  &=  &a^3  \\    
             x  &+  &ay &=  &1
          \end{linsys}
        \end{equation*}
    \end{exparts}
    \begin{answer}
       \answerasgiven
       \begin{exparts}
        \partsitem 对方程组作形式求解得到
          \begin{equation*}
            x=\frac{a^3-1}{a^2-1}  
            \qquad
            y=\frac{-a^2+a}{a^2-1}.
          \end{equation*}
          如果$a+1\neq 0$且$a-1\neq 0$，那么方程组有唯一的
          解
          \begin{equation*}
            x=\frac{a^2+a+1}{a+1}
            \qquad
            y=\frac{-a}{a+1}.
          \end{equation*}
          如果$a=-1$，或者$a=+1$，那么这些公式没有意义；在
          第一种情形，我们得到方程组
          \begin{equation*}
            \left\{ 
            \begin{linsys}{2}
              -x &+  &y  &=  &1 \\
               x &-  &y  &=  &1
            \end{linsys}\right.
          \end{equation*}
          这是一个矛盾方程组。
          在第二种情形，我们有
          \begin{equation*}
            \left\{
            \begin{linsys}{2}
               x &+  &y  &=  &1 \\
               x &+  &y  &=  &1
            \end{linsys}\right.
          \end{equation*}
          它有无穷多个解（例如，$x$任取，
          $y=1-x$）。
        \partsitem 求解方程组得到
          \begin{equation*}
            x=\frac{a^4-1}{a^2-1}
            \qquad
            y=\frac{-a^3+a}{a^2-1}.
          \end{equation*}
          这里，若$a^2-1\neq 0$，则方程组有唯一解
          $x=a^2+1$，$y=-a$。
          对于$a=-1$和$a=1$，我们得到方程组
          \begin{equation*}
            \left\{
            \begin{linsys}{2}
              -x &+  &y  &=  &-1 \\
               x &-  &y  &=  &1
            \end{linsys}\right.
            \qquad
            \left\{
            \begin{linsys}{2}
               x &+  &y  &=  &1 \\
               x &+  &y  &=  &1
            \end{linsys}\right.
         \end{equation*}
         两者都有无穷多个解。
      \end{exparts}
    \end{answer}
  \puzzle \item 
    \cite{MathMag52p48}
    在空气中，一个表面镀金的球重\( 7588 \)
    克。
    已知它可能含有铝、铜、银或铅中的一种或多种金属。
    在标准条件下依次在水中、苯中、
    酒精中和甘油中称量时，其重量分别为\( 6588 \)、\( 6688 \)、
    \( 6778 \)和\( 6328 \)克。
    如果指定的各种物质的比重取如下数值，那么它含有上述金属各多少（如果有的话）？
    \begin{center}
       \begin{tabular}{lrclr}
         铝  &\( 2.7 \)
            &\makebox[3em]{\mbox{}\hfill\mbox{}} &酒精 &0.81 \\
         铜    &\( 8.9 \)  &     &苯   &\( 0.90 \) \\
         金      &\( 19.3 \) &     &甘油 &\( 1.26 \) \\
         铅      &\( 11.3 \) &     &水     &\( 1.00 \) \\
         银    &\( 10.8 \)
       \end{tabular}
    \end{center}
    \begin{answer}
      \answerasgiven
      设\( u \)、\( v \)、\( x \)、\( y \)、\( z \)分别为球中所含
      Al、Cu、Pb、Ag和Au的体积（单位
      \( {\rm cm}^3 \)），我们假设球不是空心的。
      由于在水中（比重\( 1.00 \)）的失重
      为
      \( 1000 \)克，球的体积为\( 1000\mbox{ cm}^3 \)。
      于是这些数据（其中一些是多余的，尽管是相容的）只导出
      \( 2 \)个独立的方程，一个联系体积，另一个联系重量。
      \begin{equation*}
        \begin{linsys}{5}
           u  &+  &v    &+  &x     &+  &y     &+  &z     &=  &1000  \\
        2.7u  &+  &8.9v &+  &11.3x &+  &10.5y &+  &19.3z &=  &7558
        \end{linsys}
      \end{equation*}
      显然，球必须含有一些铝，以使其平均比重
      低于其他所有金属的比重。
      问题的这一部分没有唯一的结果，因为三种金属的量
      可以任意选取，只要所选取的量
      不会使任何金属的量为负。

      如果球只含铝和金，那么有
      \( 294.5\mbox{ cm}^3 \)的金和\( 705.5\mbox{ cm}^3 \)的铝。
      另一种可能是Cu、Au、Pb和
      Ag各\( 124.7\mbox{ cm}^3 \)，以及\( 501.2\mbox{ cm}^3  \)的Al。   
    \end{answer}
\end{exercises}

























\subsection{\texorpdfstring{通解$=$特解$+$齐次解}{通解=特解+齐次解}}
前一小节中对解集的种种描述都符合同一个模式：
它们都是一个向量——该方程组的一个特解——
加上另外一些向量的无限制组合。
\nearbyexample{ex:ManyParamsInfManySolsSystem} 给出的解集
例示了这一点。
\begin{equation*}
  \set{
   \underbracket[.7pt]{
     \colvec[r]{0 \\ 4 \\ 0 \\ 0 \\ 0}}_{\text{\shortstack{\rule{0pt}{2ex}特解}}}+
   \underbracket[.7pt]{w\colvec[r]{1 \\ -1 \\ 3 \\ 1 \\ 0}+
       u\colvec[r]{1/2 \\ -1 \\ 1/2 \\ 0 \\ 1}}_{\text{\shortstack{\rule{0pt}{2ex}无限制\\
                                                                组合}}}
       \suchthat w,u\in\Re}
\end{equation*}
组合之所以是无限制的，是指 $w$ 和 $u$ 可以取任何实数\Dash 没有
类似“要求满足 $2w-u=0\mspace{1mu}$”这样的条件来限制
我们可以使用哪些数对 $w,u$。

前面的例题展示了一个符合该模式的无穷解集。
另外两类解集也同样符合该模式。
单元素解集符合该模式，是因为它有一个特解，
而无限制组合那一部分是平凡的。
也就是说，它不是两个向量的组合，也不是一个向量的组合，
而是零个向量的组合。
（按照约定，空向量集的和是零向量。）
空解集符合该模式，是因为此时没有特解，
从而不存在这种形式的和。

\begin{theorem} \label{th:GenEqPartPlusHomo}
%<*th:GenEqPartPlusHomo>
任一线性方程组的
解集都有形式
\begin{equation*}
   \set{\vec{p}+c_1\vec{\beta}_1+\,\cdots\,+c_k\vec{\beta}_k
     \suchthat c_1,\,\ldots\,,c_k\in\Re}
\end{equation*}
其中 \( \vec{p} \) 是任意一个特解，
而向量 $\vec{\beta}_1$, \ldots, $\vec{\beta}_k$ 的个数
等于该方程组经高斯消元后自由变量的个数。
%</th:GenEqPartPlusHomo>
\end{theorem}

解的描述分为两部分：
特解 $\vec{p}$，
以及诸 $\vec{\beta}$ 的无限制线性组合。
我们将用与之相应的两个引理来证明该定理。

我们首先关注无限制的组合。
为此，我们考虑以零向量
作为特解的方程组，
这样便能将 $\vec{p}+c_1\vec{\beta}_1+\dots+c_k\vec{\beta}_k$
缩短为 $c_1\vec{\beta}_1+\dots+c_k\vec{\beta}_k$。

\begin{definition} \label{df:HomogeneousEquation}
%<*df:HomogeneousEquation>
如果一个线性方程带有常数零，从而可以写成
$a_1x_1+a_2x_2+\,\cdots\,+a_nx_n=0$ 的形式，
则称它是\definend{齐次的}\index{homogeneous equation}%
\index{linear equation!homogeneous}。
%</df:HomogeneousEquation>
\end{definition}

\begin{example}  \label{ex:FirstExHomoSys}
对任一线性方程组，比如
\begin{equation*}
  \begin{linsys}{2}
    3x  &+  &4y  &=  3  \\
    2x  &-  &y   &=  1
  \end{linsys}
\end{equation*}
我们将右端设为零，从而得到与之相应的
一个齐次方程组。
\begin{equation*}
  \begin{linsys}{2}
    3x  &+  &4y  &=  0  \\
    2x  &-  &y   &=  0
  \end{linsys}
\end{equation*}
把原方程组的化简
\begin{equation*}
  \begin{linsys}{2}
    3x  &+  &4y  &=  3  \\
    2x  &-  &y   &=  1
  \end{linsys}
  \grstep{-(2/3)\rho_1+\rho_2}
  \begin{linsys}{2}
    3x  &+  &4y        &=  3  \\
        &   &-(11/3)y   &=  -1
   \end{linsys}
\end{equation*}
与相应的齐次方程组的化简加以比较。
\begin{equation*}
  \begin{linsys}{2}
    3x  &+  &4y  &=  0  \\
    2x  &-  &y   &=  0
  \end{linsys}
  \grstep{-(2/3)\rho_1+\rho_2}
  \begin{linsys}{2}
    3x  &+  &4y        &=  0  \\
        &   &-(11/3)y   &=  0
   \end{linsys}
\end{equation*}
显然这两个化简的进行方式完全相同。
于是我们可以通过研究如何化简相应的齐次方程组，
来代替研究如何化简一个线性方程组。
\end{example}

研究相应的齐次方程组比研究原方程组
有一个很大的优点。
非齐次方程组可能无解。
但齐次方程组必定相容，因为它总是至少有
一个解：零向量。


\begin{example} \label{ex:HomoZeroOnlySol}
有些齐次方程组以零向量作为它们唯一的解。
\begin{equation*}
   \begin{linsys}{3}
     3x  &+  &2y  &+  &z  &=  &0  \\
     6x  &+  &4y  &   &   &=  &0  \\
         &   &y   &+  &z  &=  &0
   \end{linsys}
   \grstep{-2\rho_1 +\rho_2}
   \begin{linsys}{3}
      3x  &+  &2y  &+  &z  &=  &0  \\
          &   &    &   &-2z&=  &0  \\
          &   &y   &+  &z  &=  &0
    \end{linsys}
   \grstep{\rho_2 \leftrightarrow\rho_3}
   \begin{linsys}{3}
      3x  &+  &2y  &+  &z  &=  &0  \\
          &   &y   &+  &z  &=  &0  \\
          &   &    &   &-2z&=  &0
    \end{linsys}
\end{equation*}
\end{example}

\begin{example} \label{ex:SolnChemProb}
有些齐次方程组有许多解。
其中一个就是第一节第一个小节首页的化学问题\index{Chemistry problem}。
\begin{align*}
  \begin{linsys}{4}
              7x  &   &   &-  &7z  &   &   &=  &0  \\
              8x  &+  &y  &-  &5z  &-  &2w &=  &0  \\
                  &   &y  &-  &3z  &   &   &=  &0  \\
                  &   &3y &-  &6z  &-  &w  &=  &0
  \end{linsys}
  &\grstep{-(8/7)\rho_1+\rho_2}
  \begin{linsys}{4}
                7x &   &   &-  &7z  &   &   &=  &0  \\
                   &   &y  &+  &3z  &-  &2w &=  &0  \\
                   &   &y  &-  &3z  &   &   &=  &0  \\
                   &   &3y &-  &6z  &-  &w  &=  &0
   \end{linsys}                                        \\
  &\grstep[-3\rho_2+\rho_4]{-\rho_2+\rho_3}
  \begin{linsys}{4}
                7x &   &   &-  &7z  &   &   &=  &0  \\
                   &   &y  &+  &3z  &-  &2w &=  &0  \\
                   &   &   &   &-6z &+  &2w &=  &0  \\
                   &   &   &   &-15z&+  &5w &=  &0
   \end{linsys}                                        \\
  &\grstep{-(5/2)\rho_3+\rho_4}
  \begin{linsys}{4}
                7x &   &   &-  &7z  &   &   &=  &0  \\
                   &   &y  &+  &3z  &-  &2w &=  &0  \\
                   &   &   &   &-6z &+  &2w &=  &0  \\
                   &   &   &   &    &   &0  &=  &0
   \end{linsys}
\end{align*}
解集
\begin{equation*}
  \set{\colvec[r]{1/3 \\ 1 \\ 1/3 \\ 1}w \suchthat w\in\Re}
\end{equation*}
除了零向量之外还有许多向量
（如果把 \( w \) 取为分子的个数，那么只有当 \( w \) 是 $3$ 的非负倍数时，
解才有意义）。
\end{example}

\begin{lemma} \label{le:HomoSltnSpanVecs}
%<*le:HomoSltnSpanVecs>
对任一齐次线性方程组，存在向量 $\vec{\beta}_1$, \ldots, $\vec{\beta}_k$，
使得该方程组的解集为
\begin{equation*}
  \set{c_1\vec{\beta}_1+\cdots+c_k\vec{\beta}_k \suchthat c_1,\ldots,c_k\in\Re}
\end{equation*}
其中 $k$ 是该方程组的阶梯形形式中
自由变量的个数。
%</le:HomoSltnSpanVecs>
\end{lemma}

在证明之前，我们先讨论关于它的几点。
第一，该证明验证的是：给定一个阶梯形的齐次方程组，
我们可以用回代把所有首变量
都用自由变量表示出来。
要明白为什么这足以建立该引理，我们给出一个例子
（我们也将借这次逐步演算来说明
证明中的记号）。

%<*ex:HomoSystemGivesSpanningSet0>
考虑这个阶梯形的齐次方程组。
\begin{equation*}
  \begin{linsys}{5}
     x  &+  &y   &+  &2z  &+  &u  &+  &v  &=  &0  \\
        &   &y   &+  &z  &+  &u  &-  &v  &=  &0  \\
        &   &    &   &   &   &u  &+  &v  &=  &0
  \end{linsys}
\end{equation*}
它有五个变量：$x_1=x$, \ldots{}, $x_5=v$。
做回代时，我们从最底下的行，即方程~$m=3$ 开始。
利用它的首项 $a_{m,\ell_m}x_m=a_{3,4}x_4=1\cdot u$，
我们解出它的首变量 $x_{\ell_3}=x_4=u$
（记号中的“$\ell$”代表“leading”，即“首”）。
我们得到 $u=-v$。
%</ex:HomoSystemGivesSpanningSet0>

现在回代进入迭代阶段。
%<*ex:HomoSystemGivesSpanningSet1>
我们将用变量~$t$ 来数当前行位于底部之上多少行，
于是取 $t=1$，
考察行 $m-t=3-1=2$。
来到这一行时，我们已经把行~$3$ 的首变量
用自由变量表示了出来。
将 $x_{\ell_3}=x_4=u$ 用它关于
自由变量的表达式代入，得到 $y+z+(-v)-v=0$。
然后解出这一行的首变量，得
$x_{\ell_{m-t}}=x_2=y=-z+2v$。

接下来是位于底部之上 $t=2$ 的那一行，即行 $m-t=1$。
至此我们已经把行 $3$ 和行~$2$ 的首变量
用自由变量表示了出来。
代入 $x_{\ell_3}=u$ 与 $x_{\ell_2}=y$，得到
$x+(-z+2v)+2z+(-v)+v=0$。
然后解出首变量关于自由变量的表达式
$x_{\ell_1}=x=-z-2v$。
%</ex:HomoSystemGivesSpanningSet2>

%<*ex:HomoSystemGivesSpanningSet3>
这个例子的要点在于：现在我们已经完成了。
只要把解写成向量记法
\begin{equation*}
  \colvec{x \\ y \\ z \\ u \\ v}
  =\colvec{-1 \\ -1 \\ 1 \\ 0 \\ 0}z
   +\colvec{-2 \\ 2 \\ 0 \\ -1 \\ 1}v
  \qquad \text{其中 $z,v\in\Re$}
\end{equation*}
就能看出，引理陈述中的 $\vec{\beta}_1$ 与 $\vec{\beta}_2$
就是与自由变量 $z$ 和~$v$ 相关联的向量。
%</ex:HomoSystemGivesSpanningSet3>

关于证明的第二点也可以由那个例子看出。
回代沿方程组自下而上推进，用来自较低行的结论
处理下一行。
这提示了证明策略：数学归纳法。\index{mathematical induction}\index{proof techniques!induction}\index{induction}%

归纳法是一种重要而又不显然的证明技术，本书中
将会多次出现。
我们的归纳证明分两步进行：基础步骤
和归纳步骤。
在基础步骤中，我们验证该命题对某个最初的情形成立，
这里即：对阶梯形线性方程组，我们能写出
最底下方程的首变量关于自由变量的表达式。
在归纳步骤中，我们必须建立一个蕴含关系，
即：如果该命题对所有前面的情形都成立，
那么可以推出该命题对当前情形
也成立。
具体到这里，归纳假设是：如果对
下面的各行我们能把首变量
用自由变量表示出来，那么对再上面一行
我们也能把首变量用这些方式表示出来。

这两个步骤合在一起便对所有的行证明了该命题：
由基础步骤，它对最底下的方程成立，
而由归纳步骤，它对最底下方程成立这一事实
表明它对再上面一个方程也成立。
然后，再用一次
归纳步骤便推出它对再上面第三个方程成立，
如此等等。\appendrefs{mathematical induction}%

\begin{proof}
%<*pf:HomoSltnSpanVecs0>
用高斯消元法化成阶梯形。
其中可能出现一些 \( 0=0 \) 方程；我们忽略它们
（如果
方程组只由 \( 0=0 \)~方程构成，那么该引理显然成立，
因为没有首变量）。
% If there are any variables that appear only with with a coefficient of
% $0$ then they are free, and will never lead a row, so we will
% ignore that possibility.
但因为方程组是齐次的，所以没有矛盾方程。

我们将用归纳法验证
每个首变量都能用
自由变量表示出来。
这将结束证明，因为
我们可以把自由变量用作参数，
而诸 $\vec{\beta}$ 就是由那些
自由变量的系数构成的向量。
%</pf:HomoSltnSpanVecs0>

%<*pf:HomoSltnSpanVecs1>
对于基础步骤，考虑最底下的方程，
\begin{equation*}
  a_{m,\ell_m}x_{\ell_m}+a_{m,1+\ell_{m}}x_{1+\ell_{m}}+\cdots+a_{m,n}x_n=0
   % \tag{$*$}
\end{equation*}
其中 \( a_{m,\ell_m}\neq 0 \)。
%</pf:HomoSltnSpanVecs1>
%<*pf:HomoSltnSpanVecs2>
这是最底下一行，所以
首变量之后 % $x_{\ell_{m}+1}$, \ldots{}
的任何变量
必定都是自由变量。
把这些变量移到右端并除以 $a_{m,\ell_m}$
\begin{equation*}
  x_{\ell_m}
  =(-a_{m,1+\ell_{m}}/a_{m,\ell_m})x_{1+\ell_{m}}+\cdots+(-a_{m,n}/a_{m,\ell_m})x_n
\end{equation*}
便将首变量用自由变量表示了出来。
%</pf:HomoSltnSpanVecs2>
（一个技术性细节：~如果最底下的方程中
在~$x_{\ell_m}$ 的右侧没有变量，
那么 \( x_{\ell_m}=0 \)。
这也满足我们正在验证的命题，因为，
正如本小节开头所提及的，
它把 \( x_{\ell_m} \) 写成了
若干个自由变量的和，具体说是零个自由变量的和，
这里采用的惯例是：
平凡的和总计为~$0$。）

%<*pf:HomoSltnSpanVecs3>
对于归纳步骤，假设对第 \( m \) 个方程，
以及\text{第 \( (m-1) \) 个}
方程，等等，直到并包括\mbox{第 \( m-(t-1) \) 个}~方程，
我们都能
把该方程的首变量用自由变量表示出来。
也就是说，
假设“我们能将一行的首变量
用自由变量表示出来”这一命题
对最底下的 $t$~行成立，其中 $1\leq t<m$。
我们必须验证：由此该命题
对上面紧邻的那个方程，
即第 \( (m-t) \) 个方程也成立。

对 \( x_{\ell_m} \), \ldots, \( x_{\ell_{m-(t-1)}} \) 中的每一个，
代入它关于自由变量的表达式。
所得结果的首项为
$a_{m-t,\ell_{m-t}}x_{\ell_{m-t}}$，其中 \( a_{m-t,\ell_{m-t}}\neq 0 \)，
而左端的其余部分
是自由变量的组合。
把自由变量移到
右端并除以
\( a_{m-t,\ell_{m-t}} \)，最后就把
\( x_{\ell_{m-t}} \) 表示成了自由变量的表达式。

基础步骤和归纳步骤都已完成，所以由
数学归纳法原理，该命题为真。
%</pf:HomoSltnSpanVecs3>
\end{proof}

这表明，正如在引理与它的证明之间所讨论的，
我们可以用自由变量
来参数化解集。
我们说向量集
$\set{c_1\vec{\beta}_1+\cdots+c_k\vec{\beta}_k \suchthat c_1,\ldots,c_k\in\Re}$
是由集合
\( \set{\smash{\vec{\beta}_1},\ldots,\smash{\vec{\beta}_k}} \)
所\definend{生成的}\index{generated by}\index{generated}，
或者说是~\definend{由它张成的}\index{spanned by}\index{span}。

% The proof mentions a tricky point.
% We follow the convention that the sum of an empty set of vectors is the
% zero vector.
% In particular, we needs this where
% a homogeneous system has a unique solution because it
% takes the \( c \)'s to be the free variables
% and if there is a unique solution then there are no free variables.

为了完成对 \nearbytheorem{th:GenEqPartPlusHomo} 的证明，
下一个引理考虑解集描述中的
特解部分。

\begin{lemma}  \label{th:GenEqPartHomo}
%<*th:GenEqPartHomo>
对一个线性方程组以及它的任意一个特解 $\vec{p}\/$，
其解集等于
% \begin{equation*}
$
  \set{\vec{p}+\vec{h} \suchthat \text{ \( \vec{h} \) 满足
                                相应的齐次方程组}     }$。
%\end{equation*}
%</th:GenEqPartHomo>
\end{lemma}

% So fixing any particular solution gives the above
% description of the solution set.

\begin{proof}
我们将证明两个方向的集合包含：
方程组的任何解都在上述集合中，而且
上述集合中的任何向量都是该
方程组的解。\appendrefs{set equality}

%<*pf:GenEqPartHomo0>
先证第一个方向的包含：如果一个向量满足该方程组，
那么它在上述集合中。
设 \( \vec{s} \) 满足该方程组。
那么 \( \vec{s}-\vec{p} \) 满足相应的
齐次方程组，因为对每个方程下标 \( i \)，
\begin{multline*}
  a_{i,1}(s_1-p_1)+\cdots+a_{i,n}(s_n-p_n) \\
  =(a_{i,1}s_1+\cdots+a_{i,n}s_n)
  -(a_{i,1}p_1+\cdots+a_{i,n}p_n)
  =d_i-d_i
  =0
\end{multline*}
其中 \( p_j \) 与 \( s_j \) 分别是
\( \vec{p} \) 与 \( \vec{s} \) 的第 \( j \) 个分量。
把 \( \vec{s}-\vec{p} \) 记作 \( \vec{h} \)，
便将 \( \vec{s} \) 表成了所需的 \( \vec{p}+\vec{h} \) 形式。
%</pf:GenEqPartHomo0>

%<*pf:GenEqPartHomo1>
再证另一个方向的包含：取形如 $\vec{p}+\vec{h}$ 的向量，
其中 \( \vec{p} \) 满足该方程组，\( \vec{h} \) 满足
相应的齐次方程组，注意到 $\vec{p}+\vec{h}$
满足所给的方程组，因为对任一方程下标~$i$，
\begin{multline*}
  a_{i,1}(p_1+h_1)+\cdots+a_{i,n}(p_n+h_n)  \\
  =(a_{i,1}p_1+\cdots+a_{i,n}p_n)
   +(a_{i,1}h_1+\cdots+a_{i,n}h_n)
  =d_i+0
  =d_i
\end{multline*}
其中如前所述，\( p_j \) 与 \( h_j \) 分别是
\( \vec{p} \) 与 \( \vec{h} \) 的第 \( j \) 个分量。
%</pf:GenEqPartHomo1>
\end{proof}

这两个引理合起来便证明了 \nearbytheorem{th:GenEqPartPlusHomo}。
请用这句口诀记住该定理：
“\( \text{通解} = \text{特解} + \text{齐次解} \)”。

\begin{example} \label{ex:IllusGenEqPartHomo}
下面这个方程组例示了 \nearbytheorem{th:GenEqPartPlusHomo}。
\begin{equation*}
  \begin{linsys}{3}
    x  &+  &2y  &-  &z  &=  &1  \\
    2x &+  &4y  &   &   &=  &2  \\
       &   &y   &-  &3z &=  &0
  \end{linsys}
\end{equation*}
高斯消元法
\begin{equation*}
  \grstep{-2\rho_1+\rho_2}
  \begin{linsys}{3}
    x  &+  &2y  &-  &z  &=  &1  \\
       &   &    &   &2z &=  &0  \\
       &   &y   &-  &3z &=  &0
  \end{linsys}
  \grstep{\rho_2\leftrightarrow\rho_3}
  \begin{linsys}{3}
      x  &+  &2y  &-  &z  &=  &1  \\
         &   &y   &-  &3z &=  &0  \\
         &   &    &   &2z &=  &0
   \end{linsys}
\end{equation*}
表明通解是一个单元素集。
\begin{equation*}
  \set{\colvec[r]{1 \\ 0 \\ 0} }
\end{equation*}
这个单个的向量显然是一个特解。
相应的齐次方程组经同样的行变换化简
\begin{equation*}
  \begin{linsys}{3}
    x  &+  &2y  &-  &z  &=  &0  \\
    2x &+  &4y  &   &   &=  &0  \\
       &   &y   &-  &3z &=  &0
  \end{linsys}
  \grstep{-2\rho_1+\rho_2}
  \repeatedgrstep{\rho_2\swap\rho_3}
  \begin{linsys}{3}
      x  &+  &2y  &-  &z  &=  &0  \\
         &   &y   &-  &3z &=  &0  \\
         &   &    &   &2z &=  &0
   \end{linsys}
\end{equation*}
也给出一个单元素集。
\begin{equation*}
  \set{\colvec[r]{0 \\ 0 \\ 0} }
\end{equation*}
于是，正如本小节开头所讨论的，
在这种只有一个解的情形，通解就是取特解
再加上相应的齐次方程组的唯一解。
\end{example}

\begin{example}
本小节开头还讨论过，通解集为空的情形
也符合 $\text{通解}=\text{特解}+\text{齐次解}$ 的模式。
这个方程组例示了这一点。
\begin{equation*}
  \begin{linsys}{4}
    x  &   &  &+  &z  &+ &w  &=  &-1  \\
   2x  &-  &y &   &   &+ &w  &=  &3   \\
    x  &+  &y &+  &3z &+ &2w &=  &1
  \end{linsys}
  \grstep[-\rho_1+\rho_3]{-2\rho_1+\rho_2}
  \begin{linsys}{4}
    x  &   &  &+  &z  &+ &w  &=  &-1  \\
       &   &-y&-  &2z &- &w  &=  &5   \\
       &   &y &+  &2z &+ &w  &=  &2
   \end{linsys}
\end{equation*}
它无解，因为最后两个方程
相互冲突。
但相应的齐次方程组确实有解，正如所有
齐次方程组一样。
\begin{equation*}
  \begin{linsys}{4}
    x  &   &  &+  &z  &+ &w  &=  &0   \\
   2x  &-  &y &   &   &+ &w  &=  &0   \\
    x  &+  &y &+  &3z &+ &2w &=  &0
  \end{linsys}
  \grstep[-\rho_1+\rho_3]{-2\rho_1+\rho_2}
  \grstep{\rho_2+\rho_3}
  \begin{linsys}{4}
    x  &   &  &+  &z  &+ &w  &=  &0   \\
       &   &-y&-  &2z &- &w  &=  &0   \\
       &   &  &   &   &  &0  &=  &0
  \end{linsys}
\end{equation*}
事实上，这个解集是无穷的。
\begin{equation*}
  \set{\colvec[r]{-1 \\ -2 \\ 1 \\ 0}z+\colvec[r]{-1 \\ -1 \\ 0 \\ 1}w
         \suchthat z,w\in\Re}
\end{equation*}
尽管如此，由于原方程组没有特解，它的
通解集为空\Dash 不存在形如
$\vec{p}+\vec{h}$ 的向量，因为没有 $\vec{p}\:$。
\end{example}

\begin{corollary} \label{co:ThreeKindsSolutionSets}
%<*co:ThreeKindsSolutionSets>
线性方程组的解集要么为空，要么只有一个元素，
要么有无穷多个元素。
%</co:ThreeKindsSolutionSets>
\end{corollary}

\begin{proof}
%<*pf:ThreeKindsSolutionSets0>
这三种情形发生的例子我们都已见过，
所以只需证明没有其他的可能。

首先注意，一个至少有一个非 \( \zero \) 解 $\vec{v}$ 的齐次方程组
有无穷多个解。
这是因为~$\vec{v}$ 的任何标量倍数也满足该齐次
方程组，而且 $\vec{v}$ 的标量倍数组成的集合中有无穷多个向量：
如果 $s,t\in\Re$ 不相等，那么 $s\vec{v}\neq t\vec{v}$，
因为 $s\vec{v}-t\vec{v}=(s-t)\vec{v}$ 非 $\zero$：
$\vec{v}$ 的任何非 $0$ 分量在经非 $0$ 因子 $s-t$ 倍乘后，
都会给出非 $0$ 的值。
%</pf:ThreeKindsSolutionSets0>

%<*pf:ThreeKindsSolutionSets1>
现在应用 \nearbylemma{th:GenEqPartHomo} 得出：解集
\begin{equation*}
  \set{\vec{p}+\vec{h}\suchthat
    \text{\( \vec{h} \) 满足相应的齐次方程组}}
\end{equation*}
要么为空（若不存在特解 \( \vec{p} \)），
要么只有一个元素（若存在 \( \vec{p} \) 且齐次方程组
有唯一解 \( \zero \)），要么是无穷的（若存在
\( \vec{p} \) 且齐次方程组有非 $\zero$ 解，
从而由上一段可知它有无穷多个解）。
%</pf:ThreeKindsSolutionSets1>
\end{proof}

下表概括了影响通解大小的各个因素。

\smallskip
%<*table:KindsSolutionSets>
\begin{center} % \small
\begin{tabular}{r@{}c}
  &\hspace*{2.5em}\begin{tabular}{c}
      \textit{齐次方程组的} \\[-.5ex]
      \textit{解的个数}
    \end{tabular}                            \\[1.7ex]
  \begin{tabular}{r@{\hspace*{.5em}}}
     \ \\[.6ex]
     \textit{特解} \\[-.55ex]
     \textit{是否存在}
  \end{tabular}
  &\begin{tabular}{r|c@{\hspace*{1em}}c} % \cline{2-3}
     \multicolumn{1}{c}{\ }
         &\textit{一个}    &\textit{无穷多个}                    \\
     \cline{2-3}
     \textit{有}
        &\rule{0ex}{16pt}\begin{tabular}{@{}c@{}} 唯一 \\[-.5ex] 解 \end{tabular}
        &\begin{tabular}{@{}c@{}} 无穷多个 \\[-.5ex] 解 \end{tabular}
         \\[2ex] % \hline
     \textit{无}
        &\begin{tabular}{@{}c@{}} 无 \\[-.5ex] 解 \end{tabular}
        &\begin{tabular}{@{}c@{}} 无 \\[-.5ex] 解 \end{tabular}
         \\ % \hline
   \end{tabular}
\end{tabular}
\end{center}
%</table:KindsSolutionSets>
\smallskip

表中上方的那一维比较简单。
对一个线性方程组施行高斯消元法时，
忽略右端的常数，只关注
左端的系数，
那么最终要么每个变量都领头某一行，要么
我们发现某个变量不领头任何行，也就是说，
我们发现某个变量是自由变量。
（我们通过考虑相应的齐次方程组来将
“忽略右端的常数”形式化。）

一个值得注意的特殊情形是方程个数与未知数个数相同的方程组。
这样的方程组将有解，且解唯一，
当且仅当它
化简为一个每个变量都领头所在行的阶梯形方程组
（因为变量的个数与行的个数相同），
而这又当且仅当
相应的齐次方程组有唯一解。

\begin{definition} \label{df:Nonsingular}
%<*df:Nonsingular>
如果一个方阵是某个有唯一解的齐次方程组的系数矩阵，
则称它是\definend{非奇异的}\index{nonsingular!matrix}
\index{matrix!nonsingular}。
否则，即当它是某个有无穷多解的齐次方程组的系数矩阵时，
称它是\definend{奇异的}\index{singular!matrix}\index{matrix!singular}。
%</df:Nonsingular>
\end{definition}

\begin{example}
下面两个矩阵中，第一个是非奇异的而第二个是奇异的
\begin{equation*}
  \begin{mat}[r]
    1  &2  \\
    3  &4
  \end{mat}
  \qquad
  \begin{mat}[r]
    1  &2  \\
    3  &6
  \end{mat}
\end{equation*}
因为下面两个齐次方程组中第一个有唯一解，
而第二个有无穷多解。
\begin{equation*}
  \begin{linsys}[b]{2}
    x &+  &2y  &=  &0  \\
   3x &+  &4y  &=  &0
  \end{linsys}
  \qquad
  \begin{linsys}[b]{2}
    x &+  &2y  &=  &0  \\
   3x &+  &6y  &=  &0
  \end{linsys}
\end{equation*}
我们在定义中作出这一区分，是因为方程个数与变量个数相同的方程组
只有两种表现方式，取决于它的系数矩阵
是非奇异的还是奇异的。
当系数矩阵非奇异时，无论右端常数是什么，方程组
都有唯一解：~例如，高斯消元法表明方程组
\begin{equation*}
  \begin{linsys}{2}
    x  &+  &2y  &=  &a \\
    3x &+  &4y  &=  &b
  \end{linsys}
\end{equation*}
有唯一解 $x=b-2a$ 和  $y=(3a-b)/2$。
另一方面，当系数矩阵奇异时，方程组绝不会有唯一解\Dash 它
要么无解，要么有无穷多解，如下面这两个方程组。
\begin{equation*}
  \begin{linsys}[b]{2}
    x  &+  &2y  &=   &1   \\
   3x  &+  &6y  &=   &2
  \end{linsys}
  \qquad
  \begin{linsys}[b]{2}
    x  &+  &2y  &=   &1   \\
   3x  &+  &6y  &=   &3
  \end{linsys}
\end{equation*}
\end{example}

定义中使用“奇异”一词是因为它的意思是
“偏离一般的预期”。
人们常常想当然地以为变量个数与方程个数相同的方程组
会有唯一解。
因此，我们可以认为这个词隐含着
“麻烦的”，或者至少“不理想的”意味。
（“奇异”一词恰好用来指那些绝不会恰有一个解的方程组，
这有讽刺意味，但它是标准术语。）

\begin{example}
来自 \nearbyexample{ex:FirstExHomoSys}、
\nearbyexample{ex:HomoZeroOnlySol}
和 \nearbyexample{ex:IllusGenEqPartHomo} 的方程组，
其相应的齐次方程组都有唯一解。
因此这些矩阵都是非奇异的。
\begin{equation*}
  \begin{mat}[r]
    3  &4  \\
    2  &-1
  \end{mat}
  \qquad
  \begin{mat}[r]
    3  &2   &1  \\
    6  &-4  &0  \\
    0  &1   &1
  \end{mat}
  \qquad
  \begin{mat}[r]
    1  &2  &-1 \\
    2  &4  &0  \\
    0  &1  &-3
  \end{mat}
\end{equation*}
而 \nearbyexample{ex:SolnChemProb} 中的化学问题
是一个有不止一个解的齐次方程组，所以它的矩阵
是奇异的。
\begin{equation*}
  \begin{mat}[r]
    7  &0  &-7 &0  \\
    8  &1  &-5 &-2 \\
    0  &1  &-3 &0  \\
    0  &3  &-6 &-1
  \end{mat}
\end{equation*}
\end{example}

上面的表格有两个维度。
我们已经考虑了上面那一个：~只考虑方程组的左端，
我们就能判断一个给定的线性方程组落入哪一列；
右端的常数在这件事中不起作用。

表格的另一个维度，
即判断是否存在特解，就更棘手了。
考虑这两个左端相同但右端不同的方程组。
\begin{equation*}
  \begin{linsys}[b]{2}
    3x &+ &2y &= &5  \\
    3x &+ &2y &= &5
  \end{linsys}
  \qquad
  \begin{linsys}[b]{2}
    3x &+ &2y &= &5  \\
    3x &+ &2y &= &4
  \end{linsys}
\end{equation*}
第一个有解而第二个没有，所以在这里右端的常数
决定了方程组是否有解。
我们可能猜想线性方程组的左端
决定解的个数，而右端决定是否有解，
但这个猜测并不正确。
比较这两个右端相同但左端不同的方程组。
\begin{equation*}
  \begin{linsys}[b]{2}
    3x &+ &2y &= &5  \\
    4x &+ &2y &= &4
  \end{linsys}
  \qquad
  \begin{linsys}[b]{2}
    3x &+ &2y &= &5  \\
    3x &+ &2y &= &4
  \end{linsys}
\end{equation*}
第一个有解而第二个没有。
因此，仅凭方程组右端的常数
并不能决定解是否存在。
确切地说，这取决于左端与
右端之间的某种相互作用。

为了对这种相互作用有一些直观认识，
考虑这个方程组，其中一个系数未加指定，
记为变量~$c$。
\begin{equation*}
  \begin{linsys}{3}
    x  &+  &2y  &+  &3z  &=  &1  \\
    x  &+  &y   &+  &z   &=  &1  \\
   cx  &+  &3y  &+  &4z  &=  &0
  \end{linsys}
\end{equation*}
如果 \( c=2 \)，那么该方程组无解，因为左端
第三行是前两行的和，而右端不是。
如果 \( c\neq 2 \)，那么该方程组有唯一解（取 \( c=1 \) 试一试）。
要使方程组有解：如果左端系数矩阵的某一行
是其他几行的线性组合，
那么右端该行的常数必须是
同样那几行的常数的同样组合。

关于这种相互作用的更多直观认识来自对线性组合的研究。
这将是第二章的重点；在那之前，
我们要在本章其余部分完成对高斯消元法本身的研究。
\begin{exercises}
  \item 求解该方程组。
    再求解相应的齐次方程组。
    \begin{equation*}
      \begin{linsys}{4}
        x  &+ &y  &- &2z &= &0 \\
        x  &- &y  &  &   &= &-3 \\
        3x &- &y  &- &2z &= &-6  \\
           &  &2y &- &2z &= &3
      \end{linsys}
   \end{equation*}
   \begin{answer}
     下面的化简求解了该方程组。
     \begin{align*}
        \begin{amat}{3}
          1 &1   &-2 &0 \\
          1 &-1  &0  &3  \\
          3 &-1  &-2 &-6 \\
          0 &2   &-2 &3
        \end{amat}
        &\grstep[-3\rho_1+\rho_3]{-\rho_1+\rho_2}
        \begin{amat}{3}
          1 &1   &-2 &0 \\
          0 &-2  &2  &-3  \\
          0 &-4  &4  &-6 \\
          0 &2   &-2 &3
        \end{amat}                                 \\
        &\grstep[\rho_2+\rho_4]{-2\rho_2+\rho_3}
        \begin{amat}{3}
          1 &1   &-2 &0 \\
          0 &-2  &2  &-3  \\
          0 &0   &0  &0 \\
          0 &0   &0  &0
        \end{amat}
      \end{align*}
      解集如下。
      \begin{equation*}
        \set{\colvec{-3/2 \\ 3/2 \\ 0}
             +\colvec{1 \\ 1 \\ 1}z
             \suchthat z\in\Re}
      \end{equation*}
      类似地，我们可以化简相应的齐次方程组
      \begin{align*}
        \begin{amat}{3}
          1 &1   &-2 &0 \\
          1 &-1  &0  &0  \\
          3 &-1  &-2 &0 \\
          0 &2   &-2 &0
        \end{amat}
        &\grstep[-3\rho_1+\rho_3]{-\rho_1+\rho_2}
        \begin{amat}{3}
          1 &1   &-2 &0 \\
          0 &-2  &2  &0  \\
          0 &-4  &4  &0 \\
          0 &2   &-2 &0
        \end{amat}                                         \\
        &\grstep[\rho_2+\rho_4]{-2\rho_2+\rho_3}
        \begin{amat}{3}
          1 &1   &-2 &0 \\
          0 &-2  &2  &0  \\
          0 &0   &0  &0 \\
          0 &0   &0  &0
        \end{amat}
      \end{align*}
      得到它的解集。
      \begin{equation*}
        \set{
             \colvec{1 \\ 1 \\ 1}z
             \suchthat z\in\Re}
      \end{equation*}
    \end{answer}
  \recommended \item
    求解每个方程组。
    用向量表示解集。
    指出一个特解以及齐次方程组的解集。
    % (These systems also appear in
    % Exercise~2.\nearbyexercise{exer:SolveInMatrixNotation}.)
    \begin{exparts*}
      \partsitem \( \begin{linsys}[t]{2}
                  3x  &+  &6y  &=  &18  \\
                   x  &+  &2y  &=  &6
                    \end{linsys}  \)
      \partsitem \( \begin{linsys}[t]{2}
                   x  &+  &y   &=  &1  \\
                   x  &-  &y   &=  &-1
                    \end{linsys}  \)
      \partsitem \( \begin{linsys}[t]{3}
                   x_1  &   &     &+  &x_3   &=  &4  \\
                   x_1  &-  &x_2  &+  &2x_3  &=  &5  \\
                  4x_1  &-  &x_2  &+  &5x_3  &=  &17
                    \end{linsys}  \)
      \partsitem \( \begin{linsys}[t]{3}
                   2a   &+  &b    &-  &c     &=  &2  \\
                   2a   &   &     &+  &c     &=  &3  \\
                    a   &-  &b    &   &      &=  &0
                    \end{linsys}  \)
      \partsitem \( \begin{linsys}[t]{4}
                     x  &+  &2y   &-   &z   &    &    &=  &3  \\
                    2x  &+  &y    &    &    &+   &w   &=  &4  \\
                     x  &-  &y    &+   &z   &+   &w   &=  &1
                    \end{linsys}  \)
      \partsitem \( \begin{linsys}[t]{4}
                     x  &   &     &+   &z   &+   &w   &=  &4  \\
                    2x  &+  &y    &    &    &-   &w   &=  &2  \\
                    3x  &+  &y    &+   &z   &    &    &=  &7
                    \end{linsys}  \)
    \end{exparts*}
    \begin{answer}
      这几题的具体计算参见前一小节的答案。

      \begin{exparts}
        \partsitem
          解集如下。
          \begin{equation*}
            S=\set{\colvec[r]{6 \\ 0}+\colvec[r]{-2 \\ 1}y
              \suchthat y\in\Re}
          \end{equation*}
          下面是特解以及相应的齐次方程组的解集。
          \begin{equation*}
            \colvec[r]{6 \\ 0}
              \quad\text{与}\quad
            \set{\colvec[r]{-2 \\ 1}y
              \suchthat y\in\Re}
          \end{equation*}
          \textit{注。}
          这里可能有两点混淆。
          首先，上面给出的集合 $S$ 等于这个集合
          \begin{equation*}
            T=\set{\colvec[r]{4 \\ 1}+\colvec[r]{-2 \\ 1}y
              \suchthat y\in\Re}
          \end{equation*}
          因为这两个集合含有相同的元素。
          对于“什么是特解？”这一问题，下面这些都是正确的答案：
          \begin{equation*}
            \colvec{6 \\ 0},\quad
            \colvec{4 \\ 1},\quad
            \colvec{2 \\ 2},\quad
            \colvec{1 \\ 2.5}
          \end{equation*}
          第二点混淆是，我们在集合中使用的字母无关紧要。
          这个集合也等于 $S$。
          \begin{equation*}
            U=\set{\colvec[r]{6 \\ 0}+\colvec[r]{-2 \\ 1}u
                   \suchthat u\in\Re}
          \end{equation*}
        \partsitem
          解集如下。
          \begin{equation*}
            \set{\colvec[r]{0 \\ 1} }
          \end{equation*}
          下面是一个特解，
          以及相应的齐次方程组的解集。
          \begin{equation*}
            \colvec[r]{0 \\ 1}
              \qquad
            \set{\colvec[r]{0 \\ 0} }
          \end{equation*}
        \partsitem
          解集是无穷的。
          \begin{equation*}
            \set{\colvec[r]{4 \\ -1 \\ 0}+\colvec[r]{-1 \\ 1 \\ 1}x_3
              \suchthat x_3\in\Re}
          \end{equation*}
          下面是一个特解以及相应的齐次方程组的解集。
          \begin{equation*}
            \colvec[r]{4 \\ -1 \\ 0}
              \qquad
            \set{\colvec[r]{-1 \\ 1 \\ 1}x_3
              \suchthat x_3\in\Re}
          \end{equation*}
        \partsitem
          解集是单元素集。
          \begin{equation*}
            \set{\colvec[r]{1 \\ 1 \\ 1}}
          \end{equation*}
          下面是一个特解以及相应的齐次方程组的解集。
          \begin{equation*}
            \colvec[r]{1 \\ 1 \\ 1}
              \qquad
            \set{\colvec[r]{0 \\ 0 \\ 0}}
          \end{equation*}
        \partsitem
          解集是无穷的。
          \begin{equation*}
            \set{\colvec[r]{5/3 \\ 2/3 \\ 0 \\ 0}
                 +\colvec[r]{-1/3 \\ 2/3 \\ 1 \\ 0}z
                 +\colvec[r]{-2/3 \\ 1/3 \\ 0 \\ 1}w
                 \suchthat z,w\in\Re}
          \end{equation*}
          下面是一个特解以及相应的齐次方程组的解集。
          \begin{equation*}
            \colvec[r]{5/3 \\ 2/3 \\ 0 \\ 0}
              \qquad
            \set{\colvec[r]{-1/3 \\ 2/3 \\ 1 \\ 0}z
                 +\colvec[r]{-2/3 \\ 1/3 \\ 0 \\ 1}w
                 \suchthat z,w\in\Re}
          \end{equation*}
        \partsitem 该方程组的解集为空。
          因此，不存在特解。
          相应的齐次方程组的解集如下。
          \begin{equation*}
            \set{\colvec[r]{-1 \\ 2 \\ 1 \\ 0}z
                 +\colvec[r]{-1 \\ 3 \\ 0 \\ 1}w
                 \suchthat z,w\in\Re}
          \end{equation*}
      \end{exparts}
    \end{answer}
  \item
    求解每个方程组，用向量记法给出解集。
    指出一个特解以及齐次方程组的解。
    \begin{exparts*}
      \partsitem \( \begin{linsys}[t]{3}
                  2x  &+  &y  &-  &z  &=  &1  \\
                  4x  &-  &y  &   &   &=  &3
                  \end{linsys}  \)
      \partsitem \( \begin{linsys}[t]{4}
                   x  &   &   &-  &z  &   &   &=  &1  \\
                      &   &y  &+  &2z &-  &w  &=  &3  \\
                   x  &+  &2y &+  &3z &-  &w  &=  &7
                   \end{linsys}  \)
      \partsitem \( \begin{linsys}[t]{4}
                   x  &-  &y  &+  &z  &   &   &=  &0  \\
                      &   &y  &   &   &+  &w  &=  &0  \\
                  3x  &-  &2y &+  &3z &+  &w  &=  &0  \\
                      &   &-y &   &   &-  &w  &=  &0
                  \end{linsys}  \)
      \partsitem \( \begin{linsys}[t]{5}
                   a  &+  &2b &+  &3c &+  &d  &-  &e  &=  &1  \\
                  3a  &-  &b  &+  &c  &+  &d  &+  &e  &=  &3
                  \end{linsys}  \)
    \end{exparts*}
    \begin{answer}
    前一小节的答案给出了行变换。
    这里的每个答案只列出解集、特解和齐次解。
    \begin{exparts}
      \partsitem
        解集如下。
        \begin{equation*}
          \set{\colvec[r]{2/3 \\ -1/3 \\ 0}
               +\colvec[r]{1/6 \\ 2/3 \\ 1}z
              \suchthat z\in\Re}
        \end{equation*}
        下面是一个特解以及相应的齐次方程组的解集。
        \begin{equation*}
          \colvec[r]{2/3 \\ -1/3 \\ 0}
            \qquad
        \set{\colvec[r]{1/6 \\ 2/3 \\ 1}z
            \suchthat z\in\Re}
        \end{equation*}
      \partsitem
        解集是无穷的。
        \begin{equation*}
          \set{\colvec[r]{1 \\ 3 \\ 0 \\ 0}
               +\colvec[r]{1 \\-2 \\ 1 \\ 0}z
              \suchthat z\in\Re}
        \end{equation*}
        下面是
        一个特解以及相应的齐次方程组的解集。
        \begin{equation*}
          \colvec[r]{1 \\ 3 \\ 0 \\ 0}
            \qquad
          \set{\colvec[r]{1 \\ -2 \\ 1 \\ 0}z
              \suchthat z\in\Re}
        \end{equation*}
      \partsitem
        解集如下。
        \begin{equation*}
          \set{\colvec[r]{0 \\ 0 \\ 0 \\ 0}
               +\colvec[r]{-1 \\ 0 \\ 1 \\ 0}z
               +\colvec[r]{-1 \\ -1 \\ 0 \\ 1}w
              \suchthat z,w\in\Re}
        \end{equation*}
        下面是
        一个特解以及相应的齐次方程组的解集。
        \begin{equation*}
          \colvec[r]{0 \\ 0 \\ 0 \\ 0}
            \qquad
          \set{\colvec[r]{-1 \\ 0 \\ 1 \\ 0}z
               +\colvec[r]{-1 \\ -1 \\ 0 \\ 1}w
              \suchthat z,w\in\Re}
        \end{equation*}
      \partsitem
        解集如下。
        \begin{equation*}
          \set{\colvec[r]{1 \\ 0 \\ 0 \\ 0 \\ 0}
               +\colvec[r]{-5/7 \\ -8/7 \\ 1 \\ 0 \\ 0}c
               +\colvec[r]{-3/7 \\ -2/7 \\ 0 \\ 1 \\ 0}d
               +\colvec[r]{-1/7 \\ 4/7 \\ 0 \\ 0 \\ 1}e
              \suchthat c,d,e\in\Re}
        \end{equation*}
        下面是
        一个特解以及相应的齐次方程组的解集。
        \begin{equation*}
          \colvec[r]{1 \\ 0 \\ 0 \\ 0 \\ 0}
            \qquad
          \set{\colvec[r]{-5/7 \\ -8/7 \\ 1 \\ 0 \\ 0}c
               +\colvec[r]{-3/7 \\ -2/7 \\ 0 \\ 1 \\ 0}d
               +\colvec[r]{-1/7 \\ 4/7 \\ 0 \\ 0 \\ 1}e
              \suchthat c,d,e\in\Re}
        \end{equation*}
    \end{exparts}
   \end{answer}
  \recommended \item
    对于方程组
    \begin{equation*}
      \begin{linsys}{4}
       2x  &-  &y  &   &    &-  &w  &=  &3  \\
           &   &y  &+  &z   &+  &2w &=  &2  \\
        x  &-  &2y &-  &z   &   &   &=  &-1
      \end{linsys}
    \end{equation*}
    下列向量中哪些可以用作某个通解中的特解部分？
    \begin{exparts*}
      \partsitem   \( \colvec[r]{0 \\ -3 \\ 5 \\ 0} \)
      \partsitem   \( \colvec[r]{2 \\ 1 \\ 1 \\ 0} \)
      \partsitem   \( \colvec[r]{-1 \\ -4 \\ 8 \\ -1} \)
    \end{exparts*}
    \begin{answer}
      直接把它们代入，看它们是否满足全部三个方程。
      \begin{exparts}
        \partsitem 否。
        \partsitem 是。
        \partsitem 是。
      \end{exparts}
    \end{answer}
  \recommended \item
    \nearbylemma{th:GenEqPartHomo} 表明 $\vec{p}$ 可以取任意一个特解。
    若可能，求出这个方程组的通解
    \begin{equation*}
      \begin{linsys}{4}
        x  &-  &y  &   &    &+  &w  &=  &4  \\
       2x  &+  &3y &-  &z   &   &   &=  &0  \\
           &   &y  &+  &z   &+  &w  &=  &4
      \end{linsys}
    \end{equation*}
    使之以给定的向量作为其特解。
    \begin{exparts*}
      \partsitem   \( \colvec[r]{0 \\ 0 \\ 0 \\ 4} \)
      \partsitem   \( \colvec[r]{-5 \\ 1 \\ -7 \\ 10} \)
      \partsitem   \( \colvec[r]{2 \\ -1 \\ 1 \\ 1} \)
    \end{exparts*}
    \begin{answer}
      对相应的齐次方程组施行高斯消元法
      \begin{align*}
        \begin{amat}[r]{4}
           1  &-1  &0  &1  &0  \\
           2  &3   &-1 &0  &0  \\
           0  &1   &1  &1  &0
        \end{amat}
        &\grstep{-2\rho_1+\rho_2}
        \begin{amat}[r]{4}
           1  &-1  &0  &1  &0  \\
           0  &5   &-1 &-2 &0  \\
           0  &1   &1  &1  &0
        \end{amat}                                \\
        &\grstep{-(1/5)\rho_2+\rho_3}
        \begin{amat}[r]{4}
           1  &-1  &0  &1  &0  \\
           0  &5   &-1 &-2 &0  \\
           0  &0   &6/5&7/5&0
        \end{amat}
      \end{align*}
      得到这个解是齐次问题的解。
      \begin{equation*}
        \set{\colvec[r]{-5/6 \\ 1/6 \\ -7/6 \\ 1}w\suchthat w\in\Re}
      \end{equation*}
      \begin{exparts}
        \partsitem 该向量确实是一个特解，因此所求的
          通解如下。
          \begin{equation*}
            \set{\colvec[r]{0 \\ 0 \\ 0 \\ 4}+
                 \colvec[r]{-5/6 \\ 1/6 \\ -7/6 \\ 1}w\suchthat w\in\Re}
          \end{equation*}
        \partsitem 该向量是一个特解，因此所求的
          通解如下。
          \begin{equation*}
            \set{\colvec[r]{-5 \\ 1 \\ -7 \\ 10}+
                 \colvec[r]{-5/6 \\ 1/6 \\ -7/6 \\ 1}w\suchthat w\in\Re}
          \end{equation*}
        \partsitem 该向量不是该方程组的解，因为
          它不满足第三个方程。
          这样的通解不存在。
      \end{exparts}
    \end{answer}
  \item
     下面两个中，一个是非奇异的而另一个是奇异的。
     请指出哪个是哪个？
     \begin{exparts*}
       \partsitem $\begin{mat}[r]
           1  &3   \\
           4  &-12
         \end{mat}$
       \partsitem $\begin{mat}[r]
           1  &3  \\
           4  &12
         \end{mat}$
     \end{exparts*}
     \begin{answer}
       第一个非奇异而第二个奇异。
       直接用高斯消元法，看阶梯形结果的
       对角线上每个位置的数是否都非 $0$。
     \end{answer}
  \recommended \item
    奇异还是非奇异？
    \begin{exparts*}
      \partsitem \(
        \begin{mat}[r]
          1  &2  \\
          1  &3
        \end{mat}   \)
      \partsitem \(
        \begin{mat}[r]
          1  &2  \\
         -3  &-6
        \end{mat}   \)
      \partsitem \(
        \begin{mat}[r]
          1  &2  &1  \\
          1  &3  &1
        \end{mat}   \)
      \partsitem \(
        \begin{mat}[r]
          1  &2  &1  \\
          1  &1  &3  \\
          3  &4  &7
        \end{mat}   \)
      \partsitem \(
        \begin{mat}[r]
          2  &2  &1  \\
          1  &0  &5  \\
         -1  &1  &4
        \end{mat}   \)
    \end{exparts*}
    \begin{answer}
      \begin{exparts}
      \partsitem 非奇异：
        \begin{equation*}
          \grstep{-\rho_1+\rho_2}
          \begin{mat}[r]
            1  &2  \\
            0  &1
          \end{mat}
        \end{equation*}
        最终每行都含有一个首主元。
      \partsitem 奇异：
        \begin{equation*}
          \grstep{3\rho_1+\rho_2}
          \begin{mat}[r]
            1  &2  \\
            0  &0
          \end{mat}
        \end{equation*}
        最终第 \( 2 \) 行没有首主元。
      \partsitem 都不是。
        矩阵必须是方阵，这两个词才适用。
      \partsitem 奇异。
      \partsitem 非奇异。
     \end{exparts}
    \end{answer}
  \recommended \item
    给定的向量是否属于由给定的集合生成的集合？
      \begin{exparts}
        \partsitem \( \colvec[r]{2 \\ 3}, \)\
          \( \set{\colvec[r]{1 \\ 4},
                \colvec[r]{1 \\ 5}} \)
        \partsitem \( \colvec[r]{-1 \\ 0 \\ 1}, \)\
          \( \set{\colvec[r]{2 \\ 1 \\ 0},
                \colvec[r]{1 \\ 0 \\ 1}} \)
        \partsitem \( \colvec[r]{1 \\ 3 \\ 0}, \)\
          \( \set{\colvec[r]{1 \\ 0 \\ 4},
                \colvec[r]{2 \\ 1 \\ 5},
                \colvec[r]{3 \\ 3 \\ 0},
                \colvec[r]{4 \\ 2 \\ 1}} \)
        \partsitem \( \colvec[r]{1 \\ 0 \\ 1 \\ 1}, \)\
          \( \set{\colvec[r]{2 \\ 1 \\ 0 \\ 1},
                \colvec[r]{3 \\ 0 \\ 0 \\ 2}} \)
      \end{exparts}
      \begin{answer}
        在每种情形下，我们都要判断该向量是否为集合中那些向量的
        线性组合。
        \begin{exparts}
          \partsitem 是。
            求解
            \begin{equation*}
              c_1\colvec[r]{1 \\ 4}+c_2\colvec[r]{1 \\ 5}=\colvec[r]{2 \\ 3}
            \end{equation*}
            结合
            \begin{equation*}
              \begin{amat}[r]{2}
                1  &1  &2  \\
                4  &5  &3
              \end{amat}
              \grstep{-4\rho_1+\rho_2}
              \begin{amat}[r]{2}
                1  &1  &2  \\
                0  &1  &-5
              \end{amat}
            \end{equation*}
            可以得出存在给出该组合的 $c_1$ 和 $c_2$。
          \partsitem 否。
            化简
            \begin{equation*}
              \begin{amat}[r]{2}
                2  &1  &-1 \\
                1  &0  &0  \\
                0  &1  &1
              \end{amat}
              \grstep{-(1/2)\rho_1+\rho_2}
              \begin{amat}[r]{2}
                2  &1     &-1 \\
                0  &-1/2  &1/2  \\
                0  &1     &1
              \end{amat}
              \grstep{2\rho_2+\rho_3}
              \begin{amat}[r]{2}
                2  &1     &-1 \\
                0  &-1/2  &1/2  \\
                0  &0     &2
              \end{amat}
            \end{equation*}
            表明
            \begin{equation*}
              c_1\colvec[r]{2 \\ 1 \\ 0}+c_2\colvec[r]{1 \\ 0 \\ 1}
                =\colvec[r]{-1 \\ 0 \\ 1}
            \end{equation*}
            无解。
          \partsitem 是。
            化简
            \begin{align*}
              \begin{amat}[r]{4}
                1  &2  &3  &4  &1  \\
                0  &1  &3  &2  &3  \\
                4  &5  &0  &1  &0
              \end{amat}
              &\grstep{-4\rho_1+\rho_3}
              \begin{amat}[r]{4}
                1  &2  &3  &4  &1  \\
                0  &1  &3  &2  &3  \\
                0  &-3 &-12&-15&-4
              \end{amat}                   \\
              &\grstep{3\rho_2+\rho_3}
              \begin{amat}[r]{4}
                1  &2  &3  &4  &1  \\
                0  &1  &3  &2  &3  \\
                0  &0  &-3 &-9 &5
              \end{amat}
            \end{align*}
            表明有无穷多种方式
            \begin{equation*}
              \set{\colvec[r]{c_1 \\ c_2 \\ c_3 \\ c_4}=
                   \colvec[r]{-10 \\ 8 \\ -5/3 \\ 0}+
                   \colvec[r]{-9 \\ 7 \\ -3 \\ 1}c_4
                    \suchthat c_4\in\Re}
            \end{equation*}
            来写出这个组合。
            \begin{equation*}
              \colvec[r]{1 \\ 3 \\ 0}=
              c_1\colvec[r]{1 \\ 0 \\ 4}+
              c_2\colvec[r]{2 \\ 1 \\ 5}+
              c_3\colvec[r]{3 \\ 3 \\ 0}+
              c_4\colvec[r]{4 \\ 2 \\ 1}
            \end{equation*}
          \partsitem 否。
            看第三个分量。
        \end{exparts}
      \end{answer}
  \item
     证明：系数矩阵非奇异的任何线性方程组都有解，并且解是唯一的。
    \begin{answer}
        由于系数矩阵是非奇异的，高斯消元法
        结束时得到的阶梯形中每个变量都在某个方程中居于首位。
        回代给出唯一解。

      （另一种看出解唯一的方法是注意：
      当系数矩阵非奇异时，由定义，
      相应的齐次方程组有唯一解。
      由于通解是一个特解与
      每个齐次解的和，所以通解
      至多只有一个元素。）
     \end{answer}
  \item
    在引理
    \nearbylemma{le:HomoSltnSpanVecs}
    的证明中，如果没有非 \( 0=0 \) 的方程，会发生什么？
    \begin{answer}
      此时解集是整个 \( \Re^n \)，而且我们可以
      用所要求的形式表示它。
      \begin{equation*}
        \set{c_1\colvec[r]{1 \\ 0 \\ \vdotswithin{1} \\ 0}
             +c_2\colvec[r]{0 \\ 1 \\ \vdotswithin{0} \\ 0}
             +\cdots
             +c_n\colvec[r]{0 \\ 0 \\ \vdotswithin{0} \\ 1}
             \suchthat c_1,\ldots,c_n\in\Re}
      \end{equation*}
     \end{answer}
  \recommended \item
    证明：若 \( \vec{s} \) 和 \( \vec{t} \)
    满足某个齐次方程组，则下面这些向量也满足它。
    \begin{exparts*}
      \partsitem \( \vec{s}+\vec{t} \)
      \partsitem \( 3\vec{s} \)
      \partsitem \( k\vec{s}+m\vec{t} \)，其中 \( k,m\in\Re \)
    \end{exparts*}
    这个论证有什么问题：“这三个结果表明，如果齐次
    方程组有一个解，那么它就有许多个解\Dash 解的
    任意倍数是另一个解，任意一些解的和也是
    一个解\Dash 因此不存在恰好有一个解的
    齐次方程组。”？
    \begin{answer}
      设 \( \vec{s},\vec{t}\in\Re^n \)，并把它们写成如下形式。
      \begin{equation*}
        \vec{s}=\colvec{s_1 \\ \vdotswithin{s_1} \\ s_n}
          \qquad
        \vec{t}=\colvec{t_1 \\ \vdotswithin{t_1} \\ t_n}
      \end{equation*}
      再设 \( a_{i,1}x_1+\cdots+a_{i,n}x_n=0 \) 是齐次方程组中的
      第 \( i \) 个方程。
      \begin{exparts}
        \partsitem 验证很容易。
          \begin{multline*}
            a_{i,1}(s_1+t_1)+\cdots+a_{i,n}(s_n+t_n)      \\
            =
            (a_{i,1}s_1+\cdots+a_{i,n}s_n)
            +(a_{i,1}t_1+\cdots+a_{i,n}t_n)
            =
            0+0
          \end{multline*}
        \partsitem 这与前一个类似。
          \begin{equation*}
            a_{i,1}(3s_1)+\cdots+a_{i,n}(3s_n)
            =3(a_{i,1}s_1+\cdots+a_{i,n}s_n)
            =3\cdot 0=0
          \end{equation*}
        \partsitem 这个也不难多少。
          \begin{multline*}
            a_{i,1}(ks_1+mt_1)+\cdots+a_{i,n}(ks_n+mt_n)  \\
            =
            k(a_{i,1}s_1+\cdots+a_{i,n}s_n)
            +m(a_{i,1}t_1+\cdots+a_{i,n}t_n)
            =
            k\cdot 0+m\cdot 0
          \end{multline*}
      \end{exparts}
     这个论证的错误在于，只涉及
     零向量的任何线性组合都会得到零向量。
   \end{answer}
  \item
    证明：如果一个系数和常数全为有理数的方程组
    有解，那么它至少有一个全为有理数的解。
    它必定有无穷多个解吗？
    \begin{answer}
      先看证明。

      若方程组的系数全为有理数，则
      带回代的高斯消元法将只用到有理数（例如，
      \( -(m/n)\rho_i+\rho_j \)）。
      于是，我们可以只用有理数
      作为每个向量的分量来表示解集。
      因此，特解全为有理数。

      至于无穷多解的问题，
      有理向量解有无穷多个，当且仅当
      相应的齐次方程组有无穷多个
      实向量解。
      这是因为把任意参数取成有理数都会产生一个
      全为有理数的解。
   \end{answer}
\end{exercises}

\endinput




\subsection{Comparing Set Descriptions}
\emph{This subsection is optional.
Later material will not require the work here.}

A set can be described in many different ways.
Here are two different descriptions of a single set:
\begin{equation*}
  \set{\colvec[r]{1 \\ 2 \\ 3}z\suchthat z\in\Re}
  \quad\text{and}\quad
  \set{\colvec[r]{2 \\ 4 \\ 6}w\suchthat w\in\Re}.
\end{equation*}
For instance, this set contains 
\begin{equation*}
  \colvec[r]{5 \\ 10 \\ 15}
\end{equation*}
(take $z=5$ and $w=5/2$) but does not contain
\begin{equation*}
  \colvec[r]{4 \\ 8 \\ 11}
\end{equation*}
(the first component gives $z=4$ but that clashes with the third component,
similarly the first component gives $w=4/5$ but the third component 
gives something different).
Here is a third description of the same set: 
\begin{equation*}
  \set{\colvec[r]{3 \\ 6 \\ 9}+\colvec[r]{-1 \\ -2 \\ -3}y\suchthat y\in\Re}.
\end{equation*}

We need to decide when two descriptions are describing the same set.
More pragmatically stated,
how can a person tell when an answer to a homework question describes
the same set as the one described in the back of the book?

Sets are equal if and only if they have the same members.
A common way to show that two sets, $S_1$ and $S_2$, are equal is to show 
mutual inclusion:\index{sets!mutual inclusion}\index{mutual inclusion}
any member of $S_1$ is also in $S_2$, and 
any member of $S_2$ is also in $S_1$.\appendrefs{set equality}

\begin{example}
To show that 
\begin{equation*}
  S_1=
  \set{\colvec[r]{1 \\ -1 \\ 0}c+\colvec[r]{1 \\ 1 \\ 0}d\suchthat c,d\in\Re}
\end{equation*}
equals
\begin{equation*}
  S_2=
  \set{\colvec[r]{4 \\ 1 \\ 0}m+\colvec[r]{-1 \\ -3 \\ 0}n\suchthat m,n\in\Re}
\end{equation*}
we show first that $S_1\subseteq S_2$ and then that $S_2\subseteq S_1$. 

For the first half we must check that any vector
from \( S_1 \) is also in \( S_2 \).
We first consider two examples to use them as models for the general argument.
If we make up a member of $S_1$ by trying \( c=1 \) and \( d=1 \),
then to show that it is in $S_2$ we need \( m \) and $n$ such that
\begin{equation*}
  \colvec[r]{4 \\ 1 \\ 0}m
  +\colvec[r]{-1 \\ -3 \\ 0}n
  =\colvec[r]{2 \\ 0 \\ 0}
\end{equation*}
that is, this relation holds between $m$ and $n$.
\begin{equation*}
  \begin{linsys}{2}
    4m  &-  &n  &=  &2  \\
    1m  &-  &3n &=  &0  \\
        &   &0  &=  &0 
  \end{linsys}  
\end{equation*}
Similarly,
if we try \( c=2 \) and \( d=-1 \), then to show that the resulting
member of $S_1$ is in $S_2$ we need \( m \) and $n$ such that
such that 
\begin{equation*}
  \colvec[r]{4 \\ 1 \\ 0}m
  +\colvec[r]{-1 \\ -3 \\ 0}n
  =\colvec[r]{3 \\ -3 \\ 0}
\end{equation*}
that is, this holds.
\begin{equation*}
  \begin{linsys}{2}
    4m  &-  &n  &=  &3  \\
    1m  &-  &3n &=  &-3 \\
        &   &0  &=  &0 
   \end{linsys}
\end{equation*}
In the general case,
to show that any vector from \( S_1 \) is a member of \( S_2 \) we must show
that for any \( c \) and \( d \) there are appropriate \( m \) and \( n \).
We follow the pattern of the examples; fix
\begin{equation*}
  \colvec{c+d \\ -c+d \\ 0}\in S_1
\end{equation*}
and look for \( m \) and \( n \) such that
\begin{equation*}
  \colvec[r]{4 \\ 1 \\ 0}m
  +\colvec[r]{-1 \\ -3 \\ 0}n
  =\colvec{c+d \\ -c+d \\ 0}
\end{equation*}
that is, this is true.
\begin{equation*}
  \begin{linsys}{2}
    4m  &-  &n  &=  &c+d\hfill  \\
     m  &-  &3n &=  &-c+d\hfill  \\
        &   &0  &=  &0\hfill  
  \end{linsys}
\end{equation*}
Applying Gauss's Method
\begin{equation*}
  \begin{linsys}{2}
    4m  &-  &n  &=  &c+d\hfill  \\
     m  &-  &3n &=  &-c+d\hfill  
  \end{linsys}
  \grstep{-(1/4)\rho_1+\rho_2}
  \begin{linsys}{2}
    4m  &-  &n        &=  &c+d\hfill            \\
        &   &-(11/4)n &=  &-(5/4)c+(3/4)d\hfill  
   \end{linsys}
\end{equation*}
gives \( n=(5/11)c-(3/11)d \) and \( m=(4/11)c+(2/11)d \).
This shows that for any choice of $c$ and $d$ there are appropriate 
$m$ and $n$.
We conclude any member of $S_1$ is a member of $S_2$ because 
it can be rewritten in this way:
\begin{equation*}
   \colvec{c+d \\ -c+d \\ 0}
   =\colvec[r]{4 \\ 1 \\ 0}((4/11)c+(2/11)d)+
   \colvec[r]{-1 \\ -3 \\ 0}((5/11)c-(3/11)d).
\end{equation*}

For the other inclusion, \( S_2\subseteq S_1 \), we want to do the opposite.
We want to show that for any choice of $m$ and $n$ there are appropriate
$c$ and $d$.
So fix $m$ and $n$ and solve for \( c \) and \( d \):
\begin{equation*}
   \begin{linsys}{2}
     c  &+ &d  &= &4m-n\hfill \\
    -c  &+ &d  &= &m-3n\hfill 
   \end{linsys}
   \grstep{\rho_1+\rho_2}
   \begin{linsys}{2}
     c  &+ &d  &= &4m-n\hfill \\
        &  &2d &= &5m-4n\hfill 
    \end{linsys}
\end{equation*}
shows that \( d=(5/2)m-2n \) and \( c=(3/2)m+n \).
Thus any vector from \( S_2 \)
\begin{equation*}
  \colvec[r]{4 \\ 1 \\ 0}m+\colvec[r]{-1 \\ -3 \\ 0}n
\end{equation*}
is also of the right form for \( S_1 \)
\begin{equation*}
  \colvec[r]{1 \\ -1 \\ 0}((3/2)m+n)
    +\colvec[r]{1 \\ 1 \\ 0}((5/2)m-2n).
\end{equation*}
\end{example}

\begin{example}
Of course, sometimes sets are not equal.
The method of the prior example will help us see the relationship
between the two sets.
These 
\begin{equation*}
  P=
  \set{\colvec{x+y \\ 2x \\ y}\suchthat x,y\in\Re}
  \quad\text{and}\quad
  R=
  \set{\colvec{m+p \\ n \\ p}\suchthat m,n,p\in\Re}
\end{equation*}
are not equal sets.
While $P$ is a subset of $R$, it is a proper subset of $R$ because
$R$ is not a subset of $P$.

To see that, observe first that given a vector from \( P \)
we can express it in the form for \( R \)\Dash if
we fix $x$ and $y$, we can solve for appropriate $m$, $n$, and $p$:
\begin{equation*}
  \begin{linsys}{3}
     m  &   &   &+  &p  &=  &x+y\hfill  \\
        &   &n  &   &   &=  &2x\hfill   \\
        &   &   &   &p  &=  &y\hfill    
  \end{linsys}
\end{equation*}
shows that we can express any
\begin{equation*}
  \vec{v}=
  \colvec[r]{1 \\ 2 \\ 0}x+
  \colvec[r]{1 \\ 0 \\ 1}y
\end{equation*}
as a member of \( R \) with
\( m=x \), \( n=2x \), and \( p=y \):
\begin{equation*}
  \vec{v}=
  \colvec[r]{1 \\ 0 \\ 0}x+
  \colvec[r]{0 \\ 1 \\ 0}2x+
  \colvec[r]{1 \\ 0 \\ 1}y.
\end{equation*}
Thus \( P\subseteq R \).

But, for the other direction, the reduction
resulting from fixing $m$, $n$, and $p$ and looking for $x$ and $y$
\begin{align*}
  \begin{linsys}{2}
     x  &+  &y  &=  &m+p\hfill  \\
    2x  &   &   &=  &n\hfill    \\
        &   &y  &=  &p\hfill    
  \end{linsys}
  &\grstep{-2\rho_1+\rho_2}
  \begin{linsys}{2}
     x  &+  &y  &=  &m+p\hfill  \\
        &   &-2y&=  &-2m+n-2p\hfill \\
        &   &y  &=  &p\hfill    
   \end{linsys}                                  \\
  &\grstep{(1/2)\rho_2+\rho_3}
  \begin{linsys}{2}
     x  &+  &y  &=  &m+p\hfill  \\
        &   &-2y&=  &-2m+n-2p\hfill \\
        &   &0  &=  &m+(1/2)n\hfill 
    \end{linsys}
\end{align*}
shows that the only vectors
\begin{equation*}
  \colvec{m+p \\ n \\ p}\in R
\end{equation*}
representable in the form
\begin{equation*}
  \colvec{x+y \\ 2x \\ y}
\end{equation*}
are those where \( 0=m+(1/2)n \).
For instance,
\begin{equation*}
  \colvec[r]{0 \\ 1 \\ 0}
\end{equation*}
is in \( R \) but not in \( P \).
\end{example}

\begin{exercises}
  \item 
    Decide if the vector is a member of the set.
    \begin{exparts}
      \partsitem $\colvec[r]{2 \\ 3}$, 
         $\set{\colvec[r]{1 \\ 2}k\suchthat k\in\Re}$
      \partsitem $\colvec[r]{-3 \\ 3}$, 
         $\set{\colvec[r]{1 \\ -1}k\suchthat k\in\Re}$
      \partsitem $\colvec[r]{-3 \\ 3 \\ 4}$, 
             $\set{\colvec[r]{1 \\ -1 \\ 2}k\suchthat k\in\Re}$
      \partsitem $\colvec[r]{-3 \\ 3 \\ 4}$, 
             $\set{\colvec[r]{1 \\ -1 \\ 2}k+\colvec[r]{0 \\ 0 \\ 2}m
                \suchthat k,m\in\Re}$
      \partsitem $\colvec[r]{1 \\ 4 \\ 14}$, 
             $\set{\colvec[r]{2 \\ 2 \\ 5}k+\colvec[r]{-1 \\ 0 \\ 2}m
                \suchthat k,m\in\Re}$
      \partsitem $\colvec[r]{1 \\ 4 \\ 6}$, 
             $\set{\colvec[r]{2 \\ 2 \\ 5}k+\colvec[r]{-1 \\ 0 \\ 2}m
                \suchthat k,m\in\Re}$
    \end{exparts}
    \begin{answer}
      \begin{exparts}
        \partsitem No.
        \partsitem Yes.
        \partsitem No.
        \partsitem Yes.
        \partsitem Yes; use Gauss's Method to get $k=4$ and $m=-3$.
        \partsitem No; use Gauss's Method to conclude that there is no solution.
      \end{exparts}
    \end{answer}
  \item 
     Produce two descriptions of this set that are different than this one. 
     \begin{equation*}
       \set{\colvec[r]{2 \\ -5}k\suchthat k\in\Re}
     \end{equation*}
     \begin{answer}
       One easy thing to do is to double and triple the vector:
       \begin{equation*}
         \set{\colvec[r]{4 \\ -10}k\suchthat k\in\Re}
         \quad\text{and}\quad
         \set{\colvec[r]{6 \\ -55}k\suchthat k\in\Re}.
       \end{equation*}
      \end{answer}
  \recommended \item 
    Show that the three descriptions given at the start of this
    subsection all describe the same set.
    \begin{answer}
      Instead of showing all three equalities, we can show that the first
      equals the second, and that the second equals the third.
      Both equalities are easy, using the methods of this subsection.
    \end{answer}
  \recommended \item 
    Show that these sets are equal
    \begin{equation*}
      \set{\colvec[r]{1 \\ 4 \\ 1 \\ 1}
           +\colvec[r]{-1 \\ 0 \\ 1 \\ 0}z\suchthat z\in\Re  }
      \quad\text{and}\quad
      \set{\colvec[r]{0 \\ 4 \\ 2 \\ 1}
           +\colvec[r]{-1 \\ 0 \\ 1 \\ 0}k\suchthat k\in\Re  },
    \end{equation*}
    and that both describe the solution set of this system.
    \begin{equation*}  
       \begin{linsys}{4}
         x  &-  &y  &+  &z  &+  &w  &=  &-1  \\
            &   &y  &   &   &-  &w  &=  &3   \\
         x  &   &   &+  &z  &+  &2w &=  &4
       \end{linsys}
    \end{equation*}
    \begin{answer}
      That system reduces like this:
      \begin{align*}
         &\grstep{-\rho_1+\rho_2}
         \begin{linsys}{4}
           x  &-  &y  &+  &z  &+  &w  &=  &-1  \\
              &   &y  &   &   &-  &w  &=  &3   \\
              &   &y  &   &   &+  &w  &=  &5   
           \end{linsys}                              \\
         &\grstep{-\rho_2+\rho_3}
         \begin{linsys}{4}
           x  &-  &y  &+  &z  &+  &w  &=  &-1  \\
              &   &y  &   &   &-  &w  &=  &3   \\
              &   &   &   &   &   &2w &=  &2   
          \end{linsys}
      \end{align*}
      showing that \( w=1 \), \( y=4 \) and \( x=2-z \).   
    \end{answer}
  \recommended \item 
    Decide if the sets are equal.
    \begin{exparts}
      \partsitem \( \set{\colvec[r]{1 \\ 2}
                        +\colvec[r]{0 \\ 3}t
                     \suchthat t\in\Re} \)
            and
            \( \set{\colvec[r]{1 \\ 8}
                        +\colvec[r]{0 \\ -1}s
                     \suchthat s\in\Re} \)
      \partsitem \( \set{\colvec[r]{1 \\ 3 \\ 1}t
                        +\colvec[r]{2 \\ 1 \\ 5}s
                     \suchthat t,s\in\Re} \)
            and
            \( \set{\colvec[r]{4 \\ 7 \\ 7}m
                        +\colvec[r]{-4 \\ -2 \\ -10}n
                     \suchthat m,n\in\Re} \)
      \partsitem \( \set{\colvec[r]{1 \\ 2}t
                     \suchthat t\in\Re} \)
            and
            \( \set{\colvec[r]{2 \\ 4}m
                        +\colvec[r]{4 \\ 8}n
                     \suchthat m,n\in\Re} \)
      \partsitem \( \set{\colvec[r]{1 \\ 0 \\ 2}s
                        +\colvec[r]{-1 \\ 1 \\ 0}t
                     \suchthat s,t\in\Re} \)
            and
            \( \set{\colvec[r]{-1 \\ 1 \\ 1}m
                        +\colvec[r]{0 \\ 1 \\ 3}n
                     \suchthat m,n\in\Re} \)
      \partsitem \( \set{\colvec[r]{1 \\ 3 \\ 1}t
                        +\colvec[r]{2 \\ 4 \\ 6}s
                     \suchthat t,s\in\Re} \)
            and
            \( \set{\colvec[r]{3 \\ 7 \\ 7}t
                        +\colvec[r]{1 \\ 3 \\ 1}s
                     \suchthat t,s\in\Re} \)
    \end{exparts}
    \begin{answer}
      For each item, we call the first set \( S_1 \) and the
      other \( S_2 \).
     \begin{exparts}
      \partsitem They are equal.

        To see that \( S_1\subseteq S_2 \), we must show that any
        element of the
        first set is in the second, that is, for any vector of the form
        \begin{equation*}
          \vec{v}=\colvec[r]{1 \\ 2}
                  +\colvec[r]{0 \\ 3}t
        \end{equation*}
        there is an appropriate \( s \) such that
        \begin{equation*}
          \vec{v}=\colvec[r]{1 \\ 8}
                  +\colvec[r]{0 \\ -1}s.
        \end{equation*}
        Restated, given \( t \) we must find \( s \) so that this holds.
        \begin{equation*}
          \begin{linsys}{2}
            1  &+  &0s  &=  &1+0t\hfill  \\
            8  &-  &1s  &=  &2+3t\hfill  
          \end{linsys}
        \end{equation*}
        That system reduces to
        \begin{equation*}
          \begin{linsys}{1}
            1  &= &1 \hfill \\
            s  &= &6-3t
          \end{linsys}
        \end{equation*}
        That is,
        \begin{equation*}
          \colvec[r]{1 \\ 2}
          +\colvec[r]{0 \\ 3}t
          =\colvec[r]{1 \\ 8}
          +\colvec[r]{0 \\ -1}(6-3t)
        \end{equation*}
        and so we can state any vector in the form for \( S_1 \) can 
        also in the form
        needed for inclusion in \( S_2 \).

        For \( S_2\subseteq S_1 \), we look for \( t \) so that
        these equations hold.
        \begin{equation*}
          \begin{linsys}{2}
            1  &+  &0t  &=  &1+0s\hfill  \\
            2  &+  &3t  &=  &8-1s\hfill  
          \end{linsys}
        \end{equation*}
        Rewrite that as
        \begin{equation*}
          \begin{linsys}{1}
            1 &= &1\hfill   \\
            t &= &2-(1/3)s
          \end{linsys}
        \end{equation*}
        and so
        \begin{equation*}
          \colvec[r]{1 \\ 8}
          +\colvec[r]{0 \\ -1}s
          =\colvec[r]{1 \\ 2}
          +\colvec[r]{0 \\ 3}(2-(1/3)s).
        \end{equation*}
      \partsitem These two are equal.

        To show that \( S_1\subseteq S_2 \), we check that for any \( t,s \)
        we can find an appropriate \( m,n \) so that these hold.
        \begin{equation*}
          \begin{linsys}{2}
           4m  &-  &4n   &=  &1t+2s\hfill  \\
           7m  &-  &2n   &=  &3t+1s\hfill  \\
           7m  &-  &10n  &=  &1t+5s\hfill  
          \end{linsys}
        \end{equation*}
        Use Gauss's Method
        \begin{align*}
          \begin{amat}{2}
            4  &-4  &1t+2s  \\
            7  &-2  &3t+1s  \\
            7  &-10 &1t+5s
          \end{amat}
          &\grstep[(-7/4)\rho_1+\rho_3]{(-7/4)\rho_1+\rho_2}
          \begin{amat}{2}
            4  &-4  &1t+2s           \\
            0  &5   &(5/4)t-(10/4)s  \\
            0  &-3  &-(3/4)t+(6/4)s
          \end{amat}                              \\
          &\grstep{(3/5)\rho_2+\rho_3}
          \begin{amat}{2}
            4  &-4  &1t+2s           \\
            0  &5   &(5/4)t-(10/4)s  \\
            0  &0   &0
          \end{amat}
        \end{align*}
        to conclude that
        \begin{equation*}
          \colvec[r]{1 \\ 3 \\ 1}t
          +\colvec[r]{2 \\ 1 \\ 5}s
          =\colvec[r]{4 \\ 7 \\ 7}((1/2)t)
          +\colvec[r]{-4 \\ -2 \\ -10}((1/4)t-(1/2)s)
        \end{equation*}
        and so \( S_1\subseteq S_2 \).

        For \( S_2\subseteq S_1 \), solve
        \begin{equation*}
          \begin{linsys}{2}
           1t  &+  &2s   &=  &4m-4n\hfill  \\
           3t  &+  &1s   &=  &7m-2n\hfill  \\
           1t  &+  &5s   &=  &7m-10n\hfill  
          \end{linsys}
        \end{equation*}
        with Gaussian reduction
        \begin{align*}
          \begin{amat}{2}
            1  &2   &4m-4n  \\
            3  &1   &7m-2n  \\
            1  &5   &7m-10n
          \end{amat}
          &\grstep[-\rho_1+\rho_3]{-3\rho_1+\rho_2}
          \begin{amat}{2}
            1  &2   &4m-4n  \\
            0  &-5  &-5m+10n\\
            0  &3   &3m-6n
          \end{amat}                                    \\
          &\grstep{(3/5)\rho_2+\rho_3}
          \begin{amat}{2}
            1  &2   &4m-4n  \\
            0  &-5  &-5m+10n\\
            0  &0   &0
          \end{amat}
        \end{align*}
        to get
        \begin{equation*}
          \colvec[r]{4 \\ 7 \\ 7}m
          +\colvec[r]{-4 \\ -2 \\ -10}n
          =\colvec[r]{1 \\ 3 \\ 1}(2m)
          +\colvec[r]{2 \\ 1 \\ 5}(m-2n)
        \end{equation*}
        and so we can express any member of \( S_2 \) in the form needed for
        \( S_1 \).
      \partsitem These sets are equal.

        To prove that \( S_1\subseteq S_2 \), we must be able to solve
        \begin{equation*}
          \begin{linsys}{2}
           2m  &+  &4n  &=  &1t\hfill  \\
           4m  &+  &8n  &=  &2t\hfill  
          \end{linsys}
        \end{equation*}
        for \( m \) and \( n \) in terms of \( t \).
        Apply Gaussian  reduction
        \begin{equation*}
          \begin{amat}{2}
            2  &4   &1t  \\
            4  &8   &2t
          \end{amat}
          \onegrstep{-2\rho_1+\rho_2}
          \begin{amat}{2}
            2  &4   &1t  \\
            0  &0   &0
          \end{amat}
        \end{equation*}
        to conclude that
        any pair \( m,n \) where \( 2m+4n=t \) will do.
        For instance,
        \begin{equation*}
          \colvec[r]{1 \\ 2}t
          =\colvec[r]{2 \\ 4}((1/2)t)
          +\colvec[r]{4 \\ 8}(0)
        \end{equation*}
        or
        \begin{equation*}
          \colvec[r]{1 \\ 2}t
          =\colvec[r]{2 \\ 4}((-3/2)t)
          +\colvec[r]{4 \\ 8}(t).
        \end{equation*}
        Thus \( S_1\subseteq S_2 \).

        For \( S_2\subseteq S_1 \), we solve
        \begin{equation*}
          \begin{linsys}{2}
           1t  &=  &2m+4n\hfill  \\
           2t  &=  &4m+8n\hfill  
          \end{linsys}
        \end{equation*}
        with Gauss's Method
        \begin{equation*}
          \begin{amat}{1}
            1  &2m+4n  \\
            2  &4m+8n
          \end{amat}
          \grstep{-2\rho_1+\rho_2}
          \begin{amat}{1}
            1  &2m+4n  \\
            0  &0
          \end{amat}
        \end{equation*}
        to deduce that any vector in \( S_2 \) is also in \( S_1 \).
        \begin{equation*}
          \colvec[r]{2 \\ 4}m
          +\colvec[r]{4 \\ 8}n
          =\colvec[r]{1 \\ 2}(2m+4n)\in S_1.
        \end{equation*}
      \partsitem Neither set is a subset of the other.

        For \( S_1\subseteq S_2 \) to hold we must be able to solve
        \begin{equation*}
          \begin{linsys}{2}
          -1m  &+  &0n   &=  &1s-1t\hfill  \\
           1m  &+  &1n   &=  &0s+1t\hfill  \\
           1m  &+  &3n   &=  &2s+0t\hfill  
          \end{linsys}
        \end{equation*}
        for \( m \) and \( n \) in terms of \( t \) and \( s \).
        Gauss's Method
        \begin{align*}
          \begin{amat}{2}
           -1  &0   &1s-1t  \\
            1  &1   &0s+1t  \\
            1  &3   &2s+0t
          \end{amat}
          &\grstep[\rho_1+\rho_3]{\rho_1+\rho_2}
          \begin{amat}{2}
           -1  &0   &1s-1t  \\
            0  &1   &1s+0t  \\
            0  &3   &3s-1t
          \end{amat}                        \\
          &\grstep{-3\rho_2+\rho_3}
          \begin{amat}{2}
           -1  &0   &1s-1t  \\
            0  &1   &1s+0t  \\
            0  &3   &0s-1t
          \end{amat}
        \end{align*}
        shows that we can only find an appropriate pair \( m,n \) when
        \( t=0 \).
        That is,
        \begin{equation*}
          \colvec[r]{-1 \\ 1 \\ 0}
        \end{equation*}
        has no expression of the form
        \begin{equation*}
           \colvec[r]{-1 \\ 1 \\ 1}m+\colvec[r]{0 \\ 1 \\ 3}n.
        \end{equation*}

        Having shown that \( S_1 \) is not a subset of \( S_2 \), we know
        \( S_1\neq S_2 \) so, strictly speaking, we need not go further.
        But we shall also show that \( S_2 \) is not a subset of \( S_1 \).

        For \( S_2\subseteq S_1 \) to hold, we must be able to solve
        \begin{equation*}
          \begin{linsys}{2}
           1s  &-  &1t   &=  &-1m+0n\hfill  \\
           0s  &+  &1t   &=  &1m+1n\hfill  \\
           2s  &+  &0t   &=  &1m+3n\hfill  
          \end{linsys}
        \end{equation*}
        for \( s \) and \( t \).
        Apply row reduction
        \begin{align*}
          \begin{amat}{2}
            1  &-1  &-1m+0n  \\
            0  &1   &1m+1n  \\
            2  &0   &1m+3n
          \end{amat}
          &\grstep{-2\rho_1+\rho_3}
          \begin{amat}{2}
            1  &-1  &-1m+0n  \\
            0  &1   &1m+1n  \\
            0  &2   &3m+3n
          \end{amat}                                    \\
          &\grstep{-2\rho_2+\rho_3}
          \begin{amat}{2}
            1  &-1  &-1m+0n  \\
            0  &1   &1m+1n  \\
            0  &0   &1m+1n
          \end{amat}
        \end{align*}
        to deduce that the only vectors from \( S_2 \) that are also in
        \( S_1 \) are of the form
        \begin{equation*}
          \colvec[r]{-1 \\ 1 \\ 1}m
          +\colvec[r]{0 \\ 1 \\ 3}(-m).
        \end{equation*}
        For instance,
        \begin{equation*}
          \colvec{-1 \\ 1 \\ 1}
        \end{equation*}
        is in \( S_2 \) but not in \( S_1 \).
      \partsitem These sets are equal.

        First we change the parameters:
        \begin{equation*}
          S_2=\set{\colvec[r]{3 \\ 7 \\ 7}m
                   +\colvec[r]{1 \\ 3 \\ 1}n
                   \suchthat m,n\in\Re}.
        \end{equation*}

        Now, to show that \( S_1\subseteq S_2 \), we solve
        \begin{equation*}
          \begin{linsys}{2}
           3m  &+  &1n   &=  &1t+2s\hfill  \\
           7m  &+  &3n   &=  &3t+4s\hfill  \\
           7m  &+  &1n   &=  &1t+6s\hfill  
          \end{linsys}
        \end{equation*}
        with Gauss's Method
        \begin{align*}
          \begin{amat}{2}
            3  &1   &1t+2s  \\
            7  &3   &3t+4s  \\
            7  &1   &1t+6s
          \end{amat}
          &\grstep[(-7/3)\rho_1+\rho_3]{(-7/3)\rho_1+\rho_2}
          \begin{amat}{2}
            3  &1    &1t+2s  \\
            0  &2/3  &(2/3)t-(2/3)s  \\
            0  &-4/3 &(-4/3)t+(4/3)s
          \end{amat}                                    \\
          &\grstep{2\rho_2+\rho_3}
          \begin{amat}{2}
            3  &1    &1t+2s  \\
            0  &2/3  &(2/3)t-(2/3)s  \\
            0  &0    &0
          \end{amat}
        \end{align*}
        to get that
        \begin{equation*}
          \colvec[r]{1 \\ 3 \\ 1}t
          +\colvec[r]{2 \\ 4 \\ 6}s
          =\colvec[r]{3 \\ 7 \\ 7}(s)
          +\colvec[r]{1 \\ 3 \\ 1}(t-s)
        \end{equation*}
        and so \( S_1\subseteq S_2 \).

        The proof that \( S_2\subseteq S_1 \) involves solving
        \begin{equation*}
          \begin{linsys}{2}
           1t  &+  &2s   &=  &3m+1n\hfill  \\
           3t  &+  &4s   &=  &7m+3n\hfill  \\
           1t  &+  &6s   &=  &7m+1n\hfill  
          \end{linsys}
        \end{equation*}
        with Gaussian reduction
        \begin{align*}
          \begin{amat}{2}
            1  &2   &3m+1n  \\
            3  &4   &7m+3n  \\
            1  &6   &7m+1n
          \end{amat}
          &\grstep[-\rho_1+\rho_3]{-3\rho_1+\rho_2}
          \begin{amat}{2}
            1  &2   &3m+1n  \\
            0  &-2  &-2m    \\
            0  &4   &4m
          \end{amat}                                \\
          &\grstep{2\rho_2+\rho_3}
          \begin{amat}{2}
            1  &2   &3m+1n  \\
            0  &-2  &-2m    \\
            0  &0   &0
          \end{amat}
        \end{align*}
        to conclude
        \begin{equation*}
          \colvec[r]{3 \\ 7 \\ 7}m
          +\colvec[r]{1 \\ 3 \\ 1}n
          =\colvec[r]{1 \\ 3 \\ 1}(m+n)
          +\colvec[r]{2 \\ 4 \\ 6}(m)
        \end{equation*}
        and so any vector in \( S_2 \) is also in \( S_1 \).
    \end{exparts}  
   \end{answer}
\end{exercises}
